% Traduction de « C'est… que », « C'est… qui », « Ne… que » et de l'emphase ; redacción : describir a su familia (80 palabras) — Espagnol Tle A — Cours signé OnBuch+
% Programme officiel MINESEC. Compiler avec : tectonic E04-traduccion-enfasis-descripcion-familia.tex
\newcommand{\DOCMATIERE}{Espagnol}
\newcommand{\DOCNIVEAU}{Tle A}
\newcommand{\DOCMODULE}{Módulo 1 : Vie familiale et intégration sociale}
\newcommand{\DOCLECON}{E04}
\newcommand{\DOCTITRE}{Traduction de « C'est… que », « C'est… qui », « Ne… que » et de l'emphase ; redacción : describir a su familia (80 palabras)}
% OnBuch+ — préambule « cours signés » ENRICHI (compilé avec tectonic / XeLaTeX)
% Système visuel « Pop » d'OnBuch+ : fond crème (jamais blanc), bordures encre
% épaisses, ombres « dures » décalées, titres Archivo Black en capitales.
% Version MATHÉMATIQUES : graphes (pgfplots/tikz), boîtes pédagogiques
% (définition, propriété, méthode, attention, le savais-tu, exercice résolu,
% exercice, TP, bilan), page de garde et signature OnBuch+ ancrée en bas.
%
% Macros attendues dans main.tex AVANT \input{preamble} :
%   \DOCMATIERE  \DOCNIVEAU  \DOCMODULE  \DOCLECON (numéro)  \DOCTITRE
\documentclass[11pt]{article}

\usepackage{fontspec}
\usepackage[spanish,french]{babel} % français = langue principale ; l'espagnol via \esp{..} / espagnol
\usepackage[a4paper,margin=2cm,top=3.4cm,bottom=2.8cm,headheight=1.4cm,footskip=1.3cm]{geometry}
\usepackage{amsmath,amssymb,amsfonts}
\usepackage{mathtools} % \xmapsto, \coloneqq... souvent utilisés par les modèles
\usepackage{cancel} % simplifications barrées (bilans de forces)
% \abs{z} (module, valeur absolue) et \norm{u} (norme), utilisés par les modèles
\providecommand{\abs}{}\let\abs\relax\DeclarePairedDelimiter{\abs}{\lvert}{\rvert}
\providecommand{\norm}{}\let\norm\relax\DeclarePairedDelimiter{\norm}{\lVert}{\rVert}
\usepackage[table]{xcolor} % [table] : fournit aussi \rowcolors (alternance de lignes)
\usepackage{pagecolor}
\usepackage{tikz}
\usetikzlibrary{arrows.meta,positioning,calc,shapes.geometric,shapes.misc,shapes.arrows,decorations.pathmorphing,decorations.pathreplacing,decorations.markings,patterns,shadows,mindmap,trees,fit,backgrounds,matrix,intersections,angles,quotes}
\usepackage{pgfplots}
\usepackage[version=4]{mhchem} % équations nucléaires \ce{^{235}_{92}U}
\usepackage[european,siunitx]{circuitikz} % schémas électriques
\pgfplotsset{compat=1.18}
\usepgfplotslibrary{fillbetween,groupplots} % aires entre courbes (calcul intégral)
\usepackage{enumitem}
\usepackage{fancyhdr}
% [most] charge toutes les bibliothèques tcolorbox (listings, minted, raster,
% documentation...) et gonfle inutilement la mémoire TeX sur un document long
% (~300 boîtes) : on ne charge que ce qui est réellement utilisé.
\usepackage[skins,breakable]{tcolorbox}
\usepackage{graphicx}
\usepackage{colortbl}
\usepackage{booktabs}
\usepackage{tabularx}
\usepackage{diagbox} % case d'en-tête diagonale des tableaux à double entrée
% Colonnes extensibles centrée / gauche / droite (C, L, R), souvent utilisées par les modèles
\newcolumntype{C}{>{\centering\arraybackslash}X}
\newcolumntype{L}{>{\raggedright\arraybackslash}X}
\newcolumntype{R}{>{\raggedleft\arraybackslash}X}
\usepackage{array}
\usepackage{multicol}
\usepackage{siunitx}
% parse-numbers=false : accepte aussi \SI{2,5 \times 10^{-3}}{...}
\sisetup{locale=FR,output-decimal-marker={,},group-minimum-digits=4,parse-numbers=false}
\DeclareSIUnit\ohm{\ensuremath{\symup{\Omega}}} % U+2126 absent de la police
\DeclareSIUnit\division{div} % réglages d'oscilloscope (ms/div, V/div)
\DeclareSIUnit\tr{tr} % tours (vitesse de rotation, ex. \SI{1500}{\tr\per\minute})
\DeclareSIUnit\molar{M} % concentration molaire (ex. \SI{3}{\molar}, \SI{1}{\micro\molar})
\DeclareSIUnit\an{an} % année (le modèle confond parfois avec \year, un compteur TeX) — ex. \SI{5}{\kilo\meter\per\an}
\DeclareSIUnit\FCFA{FCFA} % franc CFA (ex. \SI{500}{\FCFA\per\kilo\gram})
\DeclareSIUnit\degreeBrix{\textdegree Brix} % teneur en sucre d'un sirop/jus (réfractométrie)
\DeclareSIUnit\month{mois} % le modèle utilise parfois \month, un compteur TeX, comme unité de durée
\usepackage{qrcode}
% \degree : commande fréquemment utilisée par les modèles pour un angle ou une
% température en dehors de \SI{}{} (non fournie par les paquets chargés).
\providecommand{\degree}{^{\circ}}
% \qed (fin de démonstration), utilisé par les modèles sans amsthm
\providecommand{\qed}{\hfill$\square$}
% \barre : double barre des valeurs interdites dans un tableau de variations (tkz-tab)
\providecommand{\barre}{\ensuremath{\|}}
% environnement proof (amsthm non chargé) utilisé par les modèles
\ifdefined\proof\else\newenvironment{proof}[1][Démonstration]{\par\noindent\textit{#1.}\ }{\qed\par}\fi
\usepackage{lastpage}
\usepackage{hyperref}
\hypersetup{hidelinks}

% ============================================================
% Palette « Pop » OnBuch+ — mêmes hex que la classe P (plus_ui.dart)
% ============================================================
\definecolor{poporange}{HTML}{F59321}
\definecolor{popdark}{HTML}{E07A0C}
\definecolor{popcream}{HTML}{FFF6E8}
\definecolor{popink}{HTML}{1C1714}
\definecolor{poppurple}{HTML}{7A5AE0}
\definecolor{popgreen}{HTML}{2E9E6A}
\definecolor{poppink}{HTML}{E0457B}
\definecolor{popblue}{HTML}{2F7DE1}
\definecolor{popgold}{HTML}{C98A12}
\definecolor{popmuted}{HTML}{8A7F76}
\definecolor{popline}{HTML}{EADFD2}
\definecolor{popgray}{HTML}{8A7F76}
\colorlet{popgrey}{popgray}
% Alias tolérants (le modèle nomme parfois les couleurs différemment)
\colorlet{popwhite}{white}
\colorlet{popblack}{popink}
\colorlet{popred}{poppink}
\colorlet{popyellow}{popgold}
% Teintes claires (fonds de tableaux/encadrés) — le modèle nomme parfois
% les couleurs avec un suffixe L (ex. popblueL) sans les avoir définies.
\colorlet{poporangeL}{poporange!20}
\colorlet{popdarkL}{popdark!20}
\colorlet{poppurpleL}{poppurple!20}
\colorlet{popgreenL}{popgreen!20}
\colorlet{poppinkL}{poppink!20}
\colorlet{popblueL}{popblue!20}
\colorlet{popgoldL}{popgold!20}
\colorlet{popredL}{popred!20}
\colorlet{popgrayL}{popgray!20}
% Teintes claires pour fonds de tableaux / zones
\colorlet{popblueL}{popblue!10!white}
\colorlet{poporangeL}{poporange!14!white}
\colorlet{popgreenL}{popgreen!12!white}
\colorlet{poppurpleL}{poppurple!12!white}
\colorlet{poppinkL}{poppink!12!white}
\colorlet{popgoldL}{popgold!14!white}

\pagecolor{popcream}

% ============================================================
% Typographie
% ============================================================
\defaultfontfeatures{Ligatures=TeX}
\setmainfont{PlusJakartaSans-Medium.ttf}[Path=../../../fonts/,BoldFont=PlusJakartaSans-Bold.ttf]
% Maths en sans-serif (Fira Math) avec chiffres et lettres droites de Plus
% Jakarta, pour que les formules se fondent dans le texte.
\usepackage[math-style=ISO,bold-style=ISO]{unicode-math}
\setmathfont{FiraMath-Regular.otf}[Path=../../../fonts/]
\setmathfont{PlusJakartaSans-Medium.ttf}[Path=../../../fonts/,range={up/num,up/{Latin,latin},\mathup}]
\usepackage{icomma} % virgule décimale sans espace en mode math
% Caractères mathématiques Unicode absents des polices de texte (Plus Jakarta,
% Archivo Black) : rendus par leur équivalent mathématique, en texte comme en
% formule — sinon ils s'affichent en carré vide (ex. « Arithmétique dans ℤ »).
\usepackage{newunicodechar}
\newunicodechar{ℝ}{\ensuremath{\mathbb{R}}}
\newunicodechar{ℤ}{\ensuremath{\mathbb{Z}}}
\newunicodechar{ℂ}{\ensuremath{\mathbb{C}}}
\newunicodechar{ℕ}{\ensuremath{\mathbb{N}}}
\newunicodechar{ℚ}{\ensuremath{\mathbb{Q}}}
\newunicodechar{𝔻}{\ensuremath{\mathbb{D}}}
\newunicodechar{α}{\ensuremath{\alpha}}
\newunicodechar{β}{\ensuremath{\beta}}
\newunicodechar{γ}{\ensuremath{\gamma}}
\newunicodechar{δ}{\ensuremath{\delta}}
\newunicodechar{ε}{\ensuremath{\varepsilon}}
\newunicodechar{θ}{\ensuremath{\theta}}
\newunicodechar{λ}{\ensuremath{\lambda}}
\newunicodechar{μ}{\ensuremath{\mu}}
\newunicodechar{σ}{\ensuremath{\sigma}}
\newunicodechar{τ}{\ensuremath{\tau}}
\newunicodechar{φ}{\ensuremath{\varphi}}
\newunicodechar{ψ}{\ensuremath{\psi}}
\newunicodechar{ω}{\ensuremath{\omega}}
\newunicodechar{Σ}{\ensuremath{\Sigma}}
% majuscules produites par \MakeUppercase dans les titres (α-aminés -> Α)
\newunicodechar{Α}{\ensuremath{\alpha}}
\newunicodechar{Β}{\ensuremath{\beta}}
\newunicodechar{Γ}{\ensuremath{\Gamma}}
\newunicodechar{Δ}{\ensuremath{\Delta}}
\newunicodechar{Ε}{\ensuremath{\varepsilon}}
\newunicodechar{Θ}{\ensuremath{\Theta}}
\newunicodechar{Λ}{\ensuremath{\Lambda}}
\newunicodechar{Μ}{\ensuremath{\mu}}
\newunicodechar{Τ}{\ensuremath{\tau}}
\newunicodechar{Φ}{\ensuremath{\Phi}}
\newunicodechar{Ψ}{\ensuremath{\Psi}}
\newunicodechar{Ω}{\ensuremath{\Omega}}
\newunicodechar{↦}{\ensuremath{\mapsto}}
\newunicodechar{∈}{\ensuremath{\in}}
\newunicodechar{∉}{\ensuremath{\notin}}
\newunicodechar{∧}{\ensuremath{\wedge}}
\newunicodechar{⇔}{\ensuremath{\Leftrightarrow}}
\newunicodechar{⟺}{\ensuremath{\Longleftrightarrow}}
\newunicodechar{⇒}{\ensuremath{\Rightarrow}}
\newunicodechar{⟹}{\ensuremath{\Longrightarrow}}
\newunicodechar{∘}{\ensuremath{\circ}}
\newunicodechar{∪}{\ensuremath{\cup}}
\newunicodechar{∩}{\ensuremath{\cap}}
\newunicodechar{≡}{\ensuremath{\equiv}}
\newunicodechar{⊂}{\ensuremath{\subset}}
\newunicodechar{⊥}{\ensuremath{\perp}}
\newunicodechar{∃}{\ensuremath{\exists}}
\newunicodechar{∀}{\ensuremath{\forall}}
\newunicodechar{∛}{\ensuremath{\sqrt[3]{\,}}}
\newunicodechar{∜}{\ensuremath{\sqrt[4]{\,}}}
\newunicodechar{⃗}{\ensuremath{{}^{\to}}}
\newunicodechar{ⁿ}{\ifmmode^{n}\else\textsuperscript{\ensuremath{n}}\fi}
\newunicodechar{ˣ}{\ifmmode^{x}\else\textsuperscript{\ensuremath{x}}\fi}
\newunicodechar{ᵇ}{\ifmmode^{b}\else\textsuperscript{\ensuremath{b}}\fi}
\newunicodechar{ʳ}{\ifmmode^{r}\else\textsuperscript{\ensuremath{r}}\fi}
\newunicodechar{ᵘ}{\ifmmode^{u}\else\textsuperscript{\ensuremath{u}}\fi}
\newunicodechar{ʸ}{\ifmmode^{y}\else\textsuperscript{\ensuremath{y}}\fi}
\newunicodechar{ᵃ}{\ifmmode^{a}\else\textsuperscript{\ensuremath{a}}\fi}
\newunicodechar{ᵉ}{\ifmmode^{e}\else\textsuperscript{\ensuremath{e}}\fi}
\newunicodechar{ᵏ}{\ifmmode^{k}\else\textsuperscript{\ensuremath{k}}\fi}
\newunicodechar{ᵖ}{\ifmmode^{p}\else\textsuperscript{\ensuremath{p}}\fi}
\newunicodechar{ᴺ}{\ifmmode^{N}\else\textsuperscript{\ensuremath{N}}\fi}
\newunicodechar{ᶜ}{\ifmmode^{c}\else\textsuperscript{\ensuremath{c}}\fi}
\newunicodechar{ʰ}{\ifmmode^{h}\else\textsuperscript{\ensuremath{h}}\fi}
\newunicodechar{ᵠ}{\ifmmode^{q}\else\textsuperscript{\ensuremath{q}}\fi}
\newunicodechar{ₙ}{\ifmmode_{n}\else\textsubscript{\ensuremath{n}}\fi}
\newunicodechar{ₐ}{\ifmmode_{a}\else\textsubscript{\ensuremath{a}}\fi}
\newunicodechar{ᵢ}{\ifmmode_{i}\else\textsubscript{\ensuremath{i}}\fi}
\newunicodechar{ᵣ}{\ifmmode_{r}\else\textsubscript{\ensuremath{r}}\fi}
\newunicodechar{ₖ}{\ifmmode_{k}\else\textsubscript{\ensuremath{k}}\fi}
\newunicodechar{ᵦ}{\ifmmode_{\beta}\else\textsubscript{\ensuremath{\beta}}\fi}
\newunicodechar{ₓ}{\ifmmode_{x}\else\textsubscript{\ensuremath{x}}\fi}
\newfontfamily\popdisplay{ArchivoBlack-Regular.ttf}[Path=../../../fonts/]
\newfontfamily\poplogo{SpaceGrotesk-Bold.ttf}[Path=../../../fonts/]
\newfontfamily\popsemi{PlusJakartaSans-SemiBold.ttf}[Path=../../../fonts/]
\newfontfamily\popxbold{PlusJakartaSans-ExtraBold.ttf}[Path=../../../fonts/]
\newfontfamily\popxboldls{PlusJakartaSans-ExtraBold.ttf}[Path=../../../fonts/,LetterSpace=3.0]

\setlist[enumerate]{leftmargin=1.7em,itemsep=5pt,topsep=4pt}
\setlist[itemize]{leftmargin=1.7em,itemsep=4pt,topsep=4pt,label={\color{poporange}\textbullet}}
\setlength{\parindent}{0pt}
\setlength{\parskip}{5pt}
\renewcommand{\arraystretch}{1.25}

% pgfplots : style Pop par défaut
\pgfplotsset{
  every axis/.append style={axis line style={popink,line width=1pt},tick style={popink},
    grid style={popline,line width=0.5pt},label style={font=\small},tick label style={font=\footnotesize}},
  popcurve/.style={poporange,line width=1.8pt,no markers,smooth},
}
% Même style disponible dans un \draw TikZ simple (hors pgfplots)
\tikzset{popcurve/.style={poporange,line width=1.8pt}}
\tikzset{popgrid/.style={popline,very thin}}
\colorlet{popcurve}{poporange} % le nom du style est parfois utilisé comme couleur

% ============================================================
% Pill / squiggle
% ============================================================
\newcommand{\poppill}[3]{%
  \tikz[baseline=-0.55ex]\node[draw=popink,line width=1.2pt,rounded corners=2.6pt,%
    fill=#2,inner xsep=6.5pt,inner ysep=3pt]{%
    \popxboldls\fontsize{7}{7}\selectfont\color{#3}\MakeUppercase{#1}};%
}
\newcommand{\popsquiggle}[1]{%
  \tikz[baseline=-0.4ex]{
    \draw[#1,line width=2.6pt,line cap=round]
      plot [smooth] coordinates {(0,0.09) (0.34,0.01) (0.68,0.21) (1.02,0.01) (1.36,0.10)};
  }%
}
% Mot-clé surligné (marqueur orange sous le texte)
\newcommand{\cle}[1]{{\popxbold\color{popdark}#1}}

% ============================================================
% En-tête / pied
% ============================================================
\pagestyle{fancy}
\fancyhf{}
\renewcommand{\headrulewidth}{0pt}
\renewcommand{\footrulewidth}{0pt}
\lhead{\raisebox{-0.15ex}{\poplogo\fontsize{13}{13}\selectfont\color{popink}OnBuch\color{poporange}+}}
\rhead{\poppill{\DOCMATIERE\ \textbullet\ \DOCNIVEAU\ \textbullet\ Leçon \DOCLECON}{white}{popink}}
\lfoot{\popxbold\fontsize{7.6}{7.6}\selectfont\color{popmuted}\MakeUppercase{Cours signé OnBuch+}\ \textbullet\ {\popsemi\DOCTITRE}}
\rfoot{\tikz[baseline=-0.6ex]\node[rounded corners=8.5pt,fill=popink,minimum height=17pt,minimum width=17pt,inner xsep=5pt,inner sep=0]{\popxbold\fontsize{8}{8}\selectfont\color{popcream}\thepage\,/\,\pageref*{LastPage}};}

% ============================================================
% Titres
% ============================================================
\newcommand{\chaptitre}[2]{%
  \begin{tcolorbox}[enhanced,colback=poporange,colframe=popink,boxrule=2.6pt,arc=11pt,
      drop shadow=popink,left=15pt,right=15pt,top=12pt,bottom=11pt,before skip=3pt,after skip=18pt]
    \poppill{#2}{popink}{popcream}\par\vspace{9pt}
    {\raggedright\hyphenpenalty=10000\exhyphenpenalty=10000\popdisplay\fontsize{22}{24}\selectfont\color{popcream}\MakeUppercase{#1}\par}
  \end{tcolorbox}
}
% Section numérotée : pastille orange avec le numéro + titre Archivo
\newcounter{popsec}
\newcommand{\coursec}[1]{%
  \stepcounter{popsec}\setcounter{popsub}{0}%
  \par\vspace{16pt}\noindent
  \begin{minipage}{\linewidth}
  \tikz[baseline=-0.7ex]\node[draw=popink,line width=1.6pt,fill=poporange,rounded corners=4pt,minimum size=22pt,inner sep=0]{\popdisplay\fontsize{12}{12}\selectfont\color{popcream}\thepopsec};%
  \hspace{8pt}{\popdisplay\fontsize{15}{17}\selectfont\color{popink}\MakeUppercase{#1}}\par
  \vspace{4pt}\noindent{\color{poporange}\rule{36pt}{3.6pt}}
  \end{minipage}\par\nopagebreak\vspace{8pt}
}
\newcounter{popsub}
\newcommand{\courssub}[1]{%
  \stepcounter{popsub}%
  \par\vspace{9pt}\noindent{\popxbold\fontsize{11.5}{13}\selectfont\color{popink}{\color{poporange}\thepopsec.\thepopsub}\hspace{6pt}#1}\par\nopagebreak\vspace{4pt}
}

% ============================================================
% Boîtes pédagogiques — même recette : fond blanc, bordure épaisse colorée,
% ombre dure encre, badge plein. Titre optionnel [#1].
% ============================================================
\tcbset{popbase/.style={breakable,enhanced,colback=white,boxrule=2.3pt,arc=9pt,
  shadow={3pt}{-3pt}{0pt}{popink},left=14pt,right=14pt,top=11pt,bottom=11pt,before skip=11pt,after skip=15pt,
  fontupper=\popsemi\fontsize{10.6}{14.5}\selectfont\color{popink},
  fontlower=\popsemi\fontsize{10.6}{14.5}\selectfont\color{popink}}}
% Coche dessinée (le glyphe ✓ n'existe pas dans les polices du document)
\newcommand{\popcheck}[1]{\tikz[baseline=-0.5ex]\draw[#1,line width=1.6pt,line cap=round,line join=round](-0.13,0.0)--(-0.04,-0.09)--(0.13,0.1);}
\providecommand{\coche}{\popcheck{popgreen}}
\providecommand{\resultat}[1]{\par\smallskip\noindent\fcolorbox{popgreen}{popgreenL}{\parbox{\dimexpr\linewidth-2\fboxsep-2\fboxrule\relax}{\textbf{Résultat.} #1}}\par\smallskip}
\newcommand{\popboxhead}[3]{\poppill{#1}{#2}{white}\ifx\relax#3\relax\else\hspace{6pt}{\popxbold\fontsize{10.6}{12}\selectfont\color{popink}#3}\fi\par\vspace{7pt}}

\newtcolorbox{aretenir}[1][]{popbase,colframe=popgreen,before upper={\popboxhead{À retenir}{popgreen}{#1}}}
% Texte en espagnol (document de lecture, dialogue, poème, extrait) : cadre
% d'étude. Un titre optionnel (source, auteur) s'affiche dans l'en-tête.
\newtcolorbox{textebox}[1][]{popbase,colframe=popblue,before upper={\popboxhead{Texto}{popblue}{#1}\begin{otherlanguage}{spanish}},after upper={\end{otherlanguage}}}
% Marques ✓ / ✗ (forme correcte / fautive) souvent utilisées par les modèles
\providecommand{\cross}{\textcolor{poppink}{\ensuremath{\times}}}
\providecommand{\xmark}{\cross}\providecommand{\crossmark}{\cross}\providecommand{\cmark}{\ensuremath{\checkmark}}
% Ligne de réponse / justification à compléter (inventée par les modèles)
\providecommand{\justification}[1]{\par\noindent\textit{Justification~:} \dotfill\par}
% \textipa (package tipa non chargé) : la police ne couvre pas l'API -> simple passage
\providecommand{\textipa}[1]{#1}
% Soulignement (ulem non chargé) : \uline, \ul
\providecommand{\uline}[1]{\underline{#1}}\providecommand{\ul}[1]{\underline{#1}}
% Mot ou phrase en espagnol à l'intérieur d'un paragraphe français
\newcommand{\esp}[1]{\ifmmode\text{\textit{#1}}\else\begingroup\selectlanguage{spanish}\itshape #1\endgroup\fi}
\newtcolorbox{exemplebox}[1][]{popbase,colframe=poporange,before upper={\popboxhead{Exemple}{poporange}{#1}}}
\newtcolorbox{definition}[1][]{popbase,colframe=popblue,before upper={\popboxhead{Définition}{popblue}{#1}}}
\newtcolorbox{propriete}[1][]{popbase,colframe=poppurple,before upper={\popboxhead{Propriété}{poppurple}{#1}}}
\newtcolorbox{methode}[1][]{popbase,colframe=popgold,before upper={\popboxhead{Méthode}{popgold}{#1}}}
\newtcolorbox{attention}[1][]{popbase,colframe=poppink,before upper={\popboxhead{Attention}{poppink}{#1}}}
\newtcolorbox{experience}[1][]{popbase,colframe=popblue,colback=popblueL,before upper={\popboxhead{Expérience}{popblue}{#1}}}
% « Le savais-tu ? » : carte violette pleine (contexte, culture, Cameroun)
\newtcolorbox{savaistu}[1][]{popbase,colback=poppurple,colframe=popink,
  fontupper=\popsemi\fontsize{10.4}{14.2}\selectfont\color{white},
  before upper={\poppill{Le savais-tu\,?}{white}{poppurple}\ifx\relax#1\relax\else\hspace{6pt}{\popxbold\color{white}#1}\fi\par\vspace{7pt}}}
% Exercice résolu : énoncé en haut, \tcblower puis la solution sur fond crème
\newcounter{popexo}\newcounter{popexr}
\newtcolorbox{exoresolu}[1][]{popbase,colframe=poporange,colbacklower=popcream,
  segmentation style={popink,line width=1.2pt,dashed},
  before upper={\stepcounter{popexr}\popboxhead{Exercice résolu \thepopexr}{poporange}{#1}},
  before lower={\poppill{Solution}{popink}{popcream}\par\vspace{6pt}}}
% Exercice (non résolu en place) — la correction est dans la partie Corrigés
\newtcolorbox{exercice}[1][]{popbase,colframe=popink,
  before upper={\stepcounter{popexo}\popboxhead{Exercice \thepopexo}{popink}{#1}}}
% Corrigé
\newtcolorbox{corrige}[1][]{popbase,colframe=popgreen,colback=popgreenL,
  before upper={\popboxhead{Corrigé}{popgreen}{#1}}}
% Objectifs / prérequis / bilan
\newtcolorbox{objectifs}{popbase,colframe=popink,colback=poporangeL,
  before upper={\popboxhead{Objectifs de la leçon}{popink}{}}}
\newtcolorbox{prerequis}{popbase,colframe=popink,colback=popblueL,
  before upper={\popboxhead{Prérequis}{popblue}{}}}
\newtcolorbox{bilan}{popbase,colframe=popink,colback=white,boxrule=2.8pt,
  before upper={\popboxhead{Fiche bilan}{poporange}{L'essentiel à connaître pour l'examen}}}
% Figure encadrée avec légende
\newtcolorbox{popfigure}[1][]{enhanced,breakable=false,colback=white,colframe=popline,boxrule=1.4pt,arc=8pt,
  left=10pt,right=10pt,top=10pt,bottom=8pt,before skip=10pt,after skip=12pt,halign=center,
  lower separated=false,halign lower=center,
  fontlower=\popsemi\fontsize{9}{11.5}\selectfont\color{popmuted},
  #1}
\newcommand{\legende}[1]{\par\vspace{6pt}{\popsemi\fontsize{9}{11.5}\selectfont\color{popmuted}#1}}

% ============================================================
% Signature OnBuch+
% ============================================================
\newcommand{\poplogoinline}[1]{{\poplogo\fontsize{#1}{#1}\selectfont\color{popink}OnBuch\color{poporange}+}}
\newcommand{\signatureonbuch}{%
  \begin{tcolorbox}[enhanced,colback=popink,colframe=popink,boxrule=2.2pt,arc=10pt,
      drop shadow=poporange,left=14pt,right=14pt,top=10pt,bottom=10pt,before skip=0pt,after skip=0pt]
    \tikz[baseline=-0.7ex]\node[circle,fill=popgreen,draw=popcream,line width=1.3pt,minimum size=22pt,inner sep=0]{\popcheck{white}};%
    \hspace{9pt}\begin{minipage}[c]{0.62\linewidth}
      {\popxbold\fontsize{12}{13}\selectfont\color{popcream}Signé\ \textbullet\ {\poplogo OnBuch\color{poporange}+}}\par\vspace{2pt}
      {\raggedright\popsemi\fontsize{8}{10.5}\selectfont\color{popline}\MakeUppercase{Équipe pédagogique OnBuch+}\\\MakeUppercase{Conforme au programme officiel MINESEC}\par}
    \end{minipage}\hfill
    {\poplogo\fontsize{22}{22}\selectfont\color{popcream}OnBuch\color{poporange}+}
  \end{tcolorbox}%
}

% ============================================================
% Page de garde
% #1 titre · #2 matière · #3 niveau · #4 module · #5 n° leçon · #6 durée ·
% #7 description / objectifs (texte ou itemize)
% ============================================================
\newcommand{\pagegarde}[7]{%
  \thispagestyle{empty}
  \newgeometry{a4paper,margin=1.7cm,top=1.6cm,bottom=1.4cm}%
  \begin{tikzpicture}[remember picture,overlay]
    \fill[poporange!18] ([xshift=-3.2cm,yshift=-3.6cm]current page.north east) circle (4.6cm);
    \fill[poppurple!14] ([xshift=2.4cm,yshift=7.6cm]current page.south west) circle (3.8cm);
  \end{tikzpicture}%
  {\poplogo\fontsize{30}{30}\selectfont\color{popink}OnBuch\color{poporange}+}\hfill
  \raisebox{0.6ex}{\poppill{Cours signé \textbullet\ Premium}{popink}{popcream}}\par
  \vspace{1pt}\popsquiggle{poppurple}\par
  \vspace{8pt}
  \begin{tcolorbox}[enhanced,colback=poporange,colframe=popink,boxrule=2.8pt,arc=13pt,
      drop shadow=popink,left=18pt,right=18pt,top=12pt,bottom=13pt,after skip=10pt]
    \poppill{#2 \textbullet\ #3}{popink}{popcream}\hspace{5pt}\poppill{#4}{white}{popink}\par\vspace{8pt}
    {\popdisplay\fontsize{14}{14}\selectfont\color{popink}LEÇON #5}\par\vspace{4pt}
    {\raggedright\hyphenpenalty=10000\exhyphenpenalty=10000\popdisplay\fontsize{25}{27}\selectfont\color{popcream}\MakeUppercase{#1}\par}
    \vspace{8pt}
    {\popxbold\fontsize{9.5}{9.5}\selectfont\color{popink}\MakeUppercase{Durée conseillée : #6}}
  \end{tcolorbox}
  \begin{tcolorbox}[enhanced,colback=white,colframe=popink,boxrule=2.2pt,arc=10pt,
      drop shadow=popink,left=16pt,right=16pt,top=10pt,bottom=10pt,after skip=8pt]
    \poppill{Dans cette leçon}{poppurple}{white}\par\vspace{5pt}
    \popsemi\fontsize{9.3}{12.6}\selectfont\color{popink}#7
  \end{tcolorbox}
  \noindent\begin{minipage}[t]{0.487\linewidth}
    \begin{tcolorbox}[enhanced,colback=popgreen,colframe=popink,boxrule=2.2pt,arc=10pt,
        drop shadow=popink,left=11pt,right=11pt,top=10pt,bottom=10pt,height=3.1cm,valign=center]
      \tcbox[colback=white,colframe=popink,boxrule=1.3pt,arc=5pt,boxsep=0pt,nobeforeafter,
        left=3pt,right=3pt,top=3pt,bottom=3pt,valign=center]{\qrcode[height=1.9cm]{https://play.google.com/store/apps/details?id=com.luvvix.onbuch&hl=fr}}%
      \hspace{8pt}\begin{minipage}[c]{0.5\linewidth}
        {\popdisplay\fontsize{11}{12.5}\selectfont\color{white}\MakeUppercase{Télécharge OnBuch}}\par\vspace{4pt}
        {\raggedright\popsemi\fontsize{8.4}{11}\selectfont\color{white}Gratuit \textbullet\ Google Play\\Tuteur IA, annales, concours, résultats\par}
      \end{minipage}
    \end{tcolorbox}
  \end{minipage}\hfill
  \begin{minipage}[t]{0.487\linewidth}
    \begin{tcolorbox}[enhanced,colback=poppurple,colframe=popink,boxrule=2.2pt,arc=10pt,
        drop shadow=popink,left=11pt,right=11pt,top=10pt,bottom=10pt,height=3.1cm,valign=center]
      \tikz[baseline=-0.7ex]\node[circle,fill=white,draw=popink,line width=1.3pt,minimum size=1.6cm,inner sep=0]{\popdisplay\fontsize{24}{24}\selectfont\color{poppurple}+};%
      \hspace{8pt}\begin{minipage}[c]{0.6\linewidth}
        {\popdisplay\fontsize{11}{12.5}\selectfont\color{white}\MakeUppercase{Passe à OnBuch+}}\par\vspace{4pt}
        {\raggedright\popsemi\fontsize{8.4}{11}\selectfont\color{white}Tous les cours signés, Tuteur IA illimité, fiches PDF — onglet OnBuch+ de l'app\par}
      \end{minipage}
    \end{tcolorbox}
  \end{minipage}
  \vspace{8pt}
  \popcontenu
  \vfill
  \signatureonbuch
  \restoregeometry
  \newpage
}

% Rangée « Ce cours contient » (page de garde)
\newcommand{\poptile}[3]{% #1 couleur #2 gros texte #3 légende
  \begin{tcolorbox}[enhanced,colback=#1,colframe=popink,boxrule=2pt,arc=9pt,shadow={2.5pt}{-2.5pt}{0pt}{popink},
      left=6pt,right=6pt,top=9pt,bottom=9pt,halign=center,valign=center,height=1.75cm,width=\linewidth,nobeforeafter]
    {\popdisplay\fontsize{12.5}{13}\selectfont\color{white}\MakeUppercase{#2}}\par\vspace{3pt}
    {\centering\popsemi\fontsize{7.8}{9.5}\selectfont\color{white}#3\par}
  \end{tcolorbox}}
\newcommand{\popcontenu}{%
  \par\noindent{\popxboldls\fontsize{7.5}{7.5}\selectfont\color{popmuted}CE COURS SIGNÉ CONTIENT}\par\vspace{7pt}
  \noindent\begin{minipage}[t]{0.235\linewidth}\poptile{popblue}{Cours}{complet, illustré, pas à pas}\end{minipage}\hfill
  \begin{minipage}[t]{0.235\linewidth}\poptile{poporange}{Méthodes}{et exercices résolus}\end{minipage}\hfill
  \begin{minipage}[t]{0.235\linewidth}\poptile{poppink}{Exercices}{type examen + corrigés}\end{minipage}\hfill
  \begin{minipage}[t]{0.235\linewidth}\poptile{popgreen}{Fiche bilan}{l'essentiel à retenir}\end{minipage}\par}

% Sommaire
\newtcolorbox{sommaire}{enhanced,colback=white,colframe=popink,boxrule=2.2pt,arc=10pt,
  drop shadow=popink,left=18pt,right=18pt,top=13pt,bottom=13pt,before skip=6pt,after skip=20pt,
  before upper={\poppill{Sommaire}{poppurple}{white}\par\vspace{9pt}\popsemi\fontsize{10.6}{14.5}\selectfont\color{popink}}}

% Signature de fin de cours — poussée tout en bas de la dernière page
\newcommand{\findecours}{%
  \par\vfill\nopagebreak
  \begin{center}\popsquiggle{poporange}\end{center}\vspace{4pt}
  \signatureonbuch
}
\AtBeginDocument{\def\llbracket{\mathopen{[\mkern-3.5mu[}}\def\rrbracket{\mathclose{]\mkern-3.5mu]}}} % intervalles d'entiers (glyphe ⟦ absent de la police)
% Glyphes absents de Fira Math : redéfinis à partir de glyphes disponibles
\AtBeginDocument{%
  \def\setminus{\mathbin{\backslash}}\let\smallsetminus\setminus
  \def\vdots{\mathord{\vbox{\baselineskip3.2pt\lineskiplimit0pt\kern4pt\hbox{.}\hbox{.}\hbox{.}}}}%
  \def\ddots{\mathinner{\mkern1mu\raise6.5pt\hbox{.}\mkern2mu\raise3.6pt\hbox{.}\mkern2mu\raise0.7pt\hbox{.}\mkern1mu}}%
  \def\triangle{\mathord{\scalebox{0.9}{$\symup{\Delta}$}}}%
  \def\nabla{\mathord{\raisebox{\depth}{\scalebox{1}[-1]{$\symup{\Delta}$}}}}%
  \def\checkmark{\popcheck{popdark}}%
  \def\star{\mathbin{\ast}}\def\bigstar{\mathord{\ast}}%
  \def\diamond{\mathbin{\scalebox{0.6}{$\Diamond$}}}%
  \def\bigcup{\mathop{\mathchoice{\raisebox{-0.3em}{\scalebox{1.7}{$\cup$}}}{\scalebox{1.25}{$\cup$}}{\cup}{\cup}}\displaylimits}%
  \def\bigcap{\mathop{\mathchoice{\raisebox{-0.3em}{\scalebox{1.7}{$\cap$}}}{\scalebox{1.25}{$\cap$}}{\cap}{\cap}}\displaylimits}%
  \def\lesseqqgtr{\mathrel{\lessgtr}}\def\bigoplus{\mathop{\oplus}\displaylimits}%
}
\providecommand{\ssi}{si et seulement si} % abréviation fréquente
\AtBeginDocument{\providecommand{\wideparen}{\overparen}} % arc de cercle

%% FIGFIX-BEGIN v5 — correctifs globaux des figures (docs/onbuch-generation/tools/figfix)
\usepackage{adjustbox}\usepackage{etoolbox}\tcbuselibrary{raster}
% toute figure TikZ/pgfplots plus large que la ligne est réduite à la largeur disponible
\BeforeBeginEnvironment{tikzpicture}{\begin{adjustbox}{max width=\linewidth}}
\AfterEndEnvironment{tikzpicture}{\end{adjustbox}}
% légendes pgfplots lisibles par-dessus les courbes
\pgfplotsset{every axis/.append style={legend style={fill=white,fill opacity=0.92,text opacity=1,draw=black!15,inner sep=3pt}}}
% étiquettes d'arêtes (syntaxe edge["..."]) avec fond blanc
\tikzset{every edge quotes/.append style={fill=white,inner sep=1.2pt}}
%% FIGFIX-END
\begin{document}

\pagegarde{Traduction de « C'est… que », « C'est… qui », « Ne… que » et de l'emphase ; redacción : describir a su familia (80 palabras)}{Espagnol}{Tle A}{Módulo 1 : Vie familiale et intégration sociale}{E04}{2 h}{Cette leçon de traduction étudie les équivalents espagnols des tournures d'emphase françaises (« c'est… qui/que », « ne… que ») et leurs pièges. Elle aboutit à une rédaction guidée : la description de sa famille en 80 mots.
\vspace{4pt}
\begin{multicols}{2}\small
\begin{itemize}
\item À la fin de cette leçon, je sais traduire « c'est… qui » par « es… quien / el que »
\item je sais traduire « c'est… que » par « es… que / lo que » selon la fonction du mot mis en relief
\item je sais traduire « c'est… que » avec un complément circonstanciel par « fue ayer cuando… »
\item je sais traduire « ne… que » par « solo », « únicamente », « no… más que » et « no… sino »
\item je sais distinguer « no… sino » (opposition) de « no… más que » (restriction)
\item je sais reconnaître et traduire les tournures clivées espagnoles vers le français
\end{itemize}\end{multicols}}

\begin{sommaire}
\begin{multicols}{2}
\begin{enumerate}
\item Observation : découvrir l'emphase
\item Traduire « c'est… qui »
\item Traduire « c'est… que »
\item Traduire « ne… que »
\item Lexique : la famille et le portrait
\item Rédaction guidée : describir a su familia (80 palabras)
\item Activité d'intégration
\item Méthodes et exercices résolus
\item Exercices
\item Corrigés des exercices
\item Fiche bilan
\end{enumerate}
\end{multicols}
\end{sommaire}

\begin{objectifs}
\begin{itemize}
\item À la fin de cette leçon, je sais traduire « c'est… qui » par « es… quien / el que »
\item je sais traduire « c'est… que » par « es… que / lo que » selon la fonction du mot mis en relief
\item je sais traduire « c'est… que » avec un complément circonstanciel par « fue ayer cuando… »
\item je sais traduire « ne… que » par « solo », « únicamente », « no… más que » et « no… sino »
\item je sais distinguer « no… sino » (opposition) de « no… más que » (restriction)
\item je sais reconnaître et traduire les tournures clivées espagnoles vers le français
\item je sais employer le lexique de la parenté et du portrait physique et moral
\item je sais rédiger en espagnol une description de ma famille en 80 mots
\end{itemize}
\end{objectifs}

\begin{prerequis}
\begin{itemize}
\item Le verbe ser : conjugaison au présent et au passé simple (fue)
\item Les pronoms relatifs simples : que, quien, donde
\item La négation simple : no + verbe
\item Le lexique de base de la famille (padre, madre, hermano…)
\item Les adjectifs qualificatifs et leur accord (Première)
\end{itemize}
\end{prerequis}

\newpage

\coursec{Observation : découvrir l'emphase}

\courssub{Situation d'accroche : la rentrée au lycée de Yaoundé}

C'est la rentrée au lycée de Yaoundé. Dans la classe de Terminale A, le professeur d'espagnol propose une activité simple : chacun doit présenter sa famille en quelques phrases. Aïcha lève la main. Elle veut dire : « C'est ma grand-mère qui m'a élevée » et « Je n'ai que deux frères ». Elle connaît le vocabulaire : \esp{abuela}, \esp{hermanos}, \esp{tener}... Mais au moment de parler, elle hésite. Comment rendre en espagnol ces tournures si françaises, ce « c'est... qui » et ce « ne... que » ?

Bonne nouvelle : l'espagnol possède ses propres mécanismes pour mettre un élément en relief. Aïcha dira : \esp{Es mi abuela quien me crió} et \esp{Solo tengo dos hermanos}. Cette leçon vous donne les clés pour traduire ces structures d'emphase et, au bout du chemin, rédiger le portrait de votre famille en quatre-vingts mots.

\textbf{Problématique :} comment le français et l'espagnol mettent-ils en valeur un élément de la phrase, et quelles correspondances faut-il connaître pour traduire sans faute ?

\courssub{Première observation : un texte en situation}

Lisons d'abord le petit texte qu'Aïcha finira par produire après cette leçon. Observons les phrases en gras : qu'y a-t-il de particulier dans leur construction ?

\begin{textebox}[La presentación de Aïcha]
\textbf{Es mi abuela quien me crió.} Ella vive con nosotros en Yaoundé y \textbf{es ella quien cuida de la casa} cuando mis padres trabajan. \textbf{Fue en Douala donde pasé mi infancia}, antes de venir a la capital. \textbf{Solo tengo dos hermanos}, pero \textbf{no tengo más que cariño por ellos.} Mi hermano mayor estudia medicina y \textbf{es a él a quien admiro más} de la familia. Los domingos, \textbf{es en familia como celebramos} las grandes ocasiones: comemos juntos, cantamos y recordamos a los abuelos que ya no están.
\end{textebox}

\begin{definition}[Glossaire du texte]
\begin{itemize}
\item \esp{criar} : élever (un enfant) ; \esp{me crió} : m'a élevée.
\item \esp{cuidar de la casa} : s'occuper de la maison.
\item \esp{pasar la infancia} : passer son enfance.
\item \esp{no tener más que cariño por alguien} : n'avoir que de l'affection pour quelqu'un.
\item \esp{admirar} : admirer ; \esp{a quien admiro} : que j'admire.
\item \esp{ya no están} : ne sont plus (de ce monde).
\end{itemize}
\end{definition}

\textbf{Questions d'observation :}
\begin{enumerate}
\item Dans la phrase \esp{Es mi abuela quien me crió}, quel élément est mis en valeur : la grand-mère ou l'action d'élever ?
\item Comment traduiriez-vous littéralement \esp{Fue en Douala donde pasé mi infancia} ?
\item Dans \esp{Solo tengo dos hermanos}, quel mot espagnol correspond au « ne... que » français ?
\end{enumerate}

Réponses attendues : c'est la \cle{grand-mère} qui est mise en relief ; la phrase se traduit par « C'est à Douala que j'ai passé mon enfance » ; et \esp{solo} correspond à « ne... que ». Nous venons de toucher du doigt le phénomène central de cette leçon : l'\cle{emphase}.

\courssub{Qu'est-ce que l'emphase ?}

\begin{definition}[L'emphase]
L'\cle{emphase} (de l'espagnol \esp{énfasis}) est un procédé grammatical et stylistique qui consiste à \textbf{mettre en relief} un élément de la phrase — le sujet, le complément d'objet, un complément circonstanciel — au moyen d'une tournure spéciale. Au lieu de dire simplement \esp{Mi padre trabaja aquí} (Mon père travaille ici), on dit \esp{Es mi padre quien trabaja aquí} (C'est mon père qui travaille ici) pour insister sur le fait que c'est \emph{lui}, et pas un autre.
\end{definition}

Pourquoi mettre un élément en relief ? Le plus souvent pour :
\begin{itemize}
\item \textbf{corriger} une idée fausse : \esp{No es mi tío quien vive aquí, es mi padre.} (Ce n'est pas mon oncle qui vit ici, c'est mon père.)
\item \textbf{insister} sur une personne ou un fait : \esp{Fue mi madre quien me enseñó a leer.} (C'est ma mère qui m'a appris à lire.)
\item \textbf{opposer} deux éléments : \esp{Es a ti a quien espero, no a tu hermano.} (C'est toi que j'attends, pas ton frère.)
\item \textbf{limiter} : \esp{Solo tengo dos hermanos.} (Je n'ai que deux frères.)
\end{itemize}

\courssub{Le constat : deux langues, deux mécanismes}

Observons maintenant des paires de phrases bilingues. Dans chaque cas, demandez-vous : \emph{qu'est-ce qui est mis en valeur, et comment chaque langue s'y prend-elle ?}

\begin{exemplebox}[Mise en relief du sujet]
\esp{Es mi padre quien trabaja aquí.} « C'est mon père qui travaille ici. »\par
\esp{Es mi abuela quien me crió.} « C'est ma grand-mère qui m'a élevée. »\par
\esp{Son mis hermanos quienes cocinan los domingos.} « Ce sont mes frères qui cuisinent le dimanche. »
\end{exemplebox}

\begin{exemplebox}[Mise en relief du complément d'objet]
\esp{Es a ti a quien espero.} « C'est toi que j'attends. »\par
\esp{Es a mi hermano mayor a quien admiro.} « C'est mon frère aîné que j'admire. »\par
\esp{Es el cacao lo que exporta la región.} « C'est le cacao qu'exporte la région. »
\end{exemplebox}

\begin{exemplebox}[Mise en relief d'un complément circonstanciel]
\esp{Fue en Douala donde pasé mi infancia.} « C'est à Douala que j'ai passé mon enfance. »\par
\esp{Fue ayer cuando llegó mi tío de España.} « C'est hier que mon oncle est arrivé d'Espagne. »\par
\esp{Es en familia como celebramos las fiestas.} « C'est en famille que nous fêtons les grandes occasions. »
\end{exemplebox}

Le constat est clair : le français \cle{clive} la phrase avec \textbf{c'est... qui} (devant un sujet) et \textbf{c'est... que} (devant un complément) ; l'espagnol clive avec \esp{es... quien} (devant une personne sujet), \esp{es... a quien} (devant une personne complément), \esp{es... lo que} (devant une chose complément) et \esp{fue... cuando/donde} (devant une circonstance de temps ou de lieu). Une phrase clivée est une phrase « coupée en deux » : le premier morceau porte l'élément mis en relief, le second porte le reste du message.

\begin{popfigure}
\begin{tikzpicture}[node distance=0.4cm]
\node[draw=popblue, fill=popblueL, rounded corners, align=center, text width=5.6cm] (fr) {\textbf{Français}\\ \textcolor{popink}{\textbf{C'est mon père}} \textbf{qui} travaille ici.};
\node[draw=popgreen, fill=popgreenL, rounded corners, align=center, text width=5.6cm, right=2.4cm of fr] (es) {\textbf{Espagnol}\\ \textcolor{popink}{\textbf{Es mi padre}} \textbf{quien} trabaja aquí.};
\draw[-{Stealth[length=3mm]}, very thick, poporange] (fr) -- (es) node[midway, above, align=center, text=popdark]{\footnotesize traduction};
\node[below=0.3cm of fr, align=center, text width=5.6cm] {\footnotesize élément en relief : \textcolor{popink}{\textbf{sujet}}};
\node[below=0.3cm of es, align=center, text width=5.6cm] {\footnotesize élément en relief : \textcolor{popink}{\textbf{sujet}}};
\end{tikzpicture}
\legende{De la phrase clivée française à la phrase clivée espagnole : l'élément mis en relief (en couleur) est placé juste après \esp{es}, comme après « c'est ».}
\end{popfigure}

\begin{aretenir}[Le principe de la correspondance]
La structure française \textbf{c'est + X + qui/que} se traduit presque toujours par \esp{es + X + quien/que} en espagnol. L'ordre est le même ; seuls les mots-outils changent. La difficulté consiste à choisir le bon mot-outil espagnol : \esp{quien}, \esp{a quien}, \esp{que}, \esp{lo que}, \esp{donde}, \esp{cuando}. Ce sera l'objet des sections suivantes.
\end{aretenir}

\courssub{Premier tableau comparatif des correspondances}

\begin{center}
\begin{tabularx}{\linewidth}{|X|X|X|}
\hline
\rowcolor{popblueL} \textbf{Français} & \textbf{Español} & \textbf{Remarque} \\
\hline
C'est mon père \textbf{qui} travaille. & \esp{Es mi padre \textbf{quien} trabaja.} & sujet personne : \esp{quien} \\
\hline
Ce sont mes sœurs \textbf{qui} chantent. & \esp{Son mis hermanas \textbf{quienes} cantan.} & pluriel : \esp{son... quienes} \\
\hline
C'est toi \textbf{que} j'attends. & \esp{Es a ti \textbf{a quien} espero.} & COD personne : \esp{a quien} \\
\hline
C'est le cacao \textbf{que} j'achète. & \esp{Es el cacao \textbf{lo que} compro.} & COD chose : \esp{lo que} \\
\hline
C'est à Douala \textbf{que} je suis né. & \esp{Fue en Douala \textbf{donde} nací.} & lieu : \esp{donde} \\
\hline
C'est hier \textbf{qu'}il est arrivé. & \esp{Fue ayer \textbf{cuando} llegó.} & temps : \esp{cuando} \\
\hline
Je \textbf{ne} \textbf{que} deux frères. (n'ai que) & \esp{\textbf{Solo} tengo dos hermanos.} & restriction : \esp{solo} \\
\hline
Il \textbf{ne} fait \textbf{que} parler. & \esp{\textbf{No} hace \textbf{más que} hablar.} & restriction : \esp{no... más que} \\
\hline
\end{tabularx}
\end{center}

\begin{attention}[Piège du francophone : le calque mot à mot]
Ne traduisez jamais « c'est... que » par \esp{es... que} devant une personne sujet : \esp{Es mi padre que trabaja aquí} est fautif ; il faut \esp{Es mi padre \textbf{quien} trabaja aquí}. De même, « ne... que » ne se traduit jamais par \esp{no... que} seul : « Je n'ai que deux frères » n'est pas \esp{No tengo que dos hermanos} mais \esp{\textbf{Solo} tengo dos hermanos} ou \esp{\textbf{No} tengo \textbf{más que} dos hermanos}. Attention aussi au temps du verbe \esp{ser} : « C'est hier que... » exige le passé : \esp{\textbf{Fue} ayer cuando...}.
\end{attention}

\courssub{Carte mentale : les quatre visages de l'emphase}

\begin{popfigure}
\begin{tikzpicture}[mindmap, grow cyclic, every node/.style={concept}, concept color=popblue!40,
root concept/.append style={concept color=poporange!60, font=\bfseries},
level 1/.append style={level distance=3.4cm, sibling angle=90},
level 2/.append style={level distance=2.6cm, sibling angle=50}]
\node[root concept, align=center]{Emphase\\ \esp{Énfasis}}
  child[concept color=popgreen!50]{ node[align=center]{Sujet\\ \esp{Es mi padre quien...}} }
  child[concept color=poppurple!40]{ node[align=center]{COD\\ \esp{Es a ti a quien...}\\ \esp{Es el libro lo que...}} }
  child[concept color=popgold!50]{ node[align=center]{Circonstance\\ \esp{Fue ayer cuando...}\\ \esp{Fue en Douala donde...}} }
  child[concept color=poppink!50]{ node[align=center]{Restriction\\ \esp{Solo tengo dos...}\\ \esp{No... más que}} };
\end{tikzpicture}
\legende{Carte mentale de l'emphase : quatre types d'éléments peuvent être mis en relief — le sujet, le COD, la circonstance, et la restriction (« ne... que »).}
\end{popfigure}

\courssub{Exercice résolu d'observation}

\begin{exoresolu}[Repérer l'élément mis en relief]
Énoncé : pour chaque phrase espagnole, soulignez l'élément mis en relief, précisez sa nature (sujet, COD, circonstance, restriction) et traduisez en français avec une tournure clivée.\par
1. \esp{Es mi tía quien vende tela en el mercado de Yaoundé.}\par
2. \esp{Es a mis primos a quienes visito los domingos.}\par
3. \esp{Fue en 1960 cuando Camerún se independizó.}\par
4. \esp{Solo hablo dos lenguas: francés y español.}
\tcblower
Solution détaillée :\par
1. Élément en relief : \esp{mi tía} — \textbf{sujet} (c'est elle qui fait l'action). Traduction : « C'est ma tante qui vend du tissu au marché de Yaoundé. » Le mot-outil \esp{quien} confirme qu'il s'agit d'une personne sujet.\par
2. Élément en relief : \esp{a mis primos} — \textbf{COD} (ce sont eux que je visite). Traduction : « C'est mes cousins que je rends visite le dimanche » ou mieux : « C'est chez mes cousins que je vais le dimanche ». Le \esp{a quienes} (pluriel) signale un complément de personne.\par
3. Élément en relief : \esp{en 1960} — \textbf{circonstance de temps}. Traduction : « C'est en 1960 que le Cameroun est devenu indépendant. » Le mot-outil \esp{cuando} et le prétérit \esp{fue} confirment la mise en relief d'un moment passé.\par
4. Élément en relief : \esp{dos lenguas} — \textbf{restriction}. Traduction : « Je ne parle que deux langues : le français et l'espagnol. » \esp{Solo} placé devant le verbe équivaut exactement à « ne... que ».
\end{exoresolu}

\begin{savaistu}[L'emphase par simple déplacement du mot]
L'espagnol possède une arme d'emphase que le français n'a pas : l'\cle{ordre des mots} libre. En déplaçant simplement un complément en tête de phrase, on le met en relief sans aucune tournure spéciale : \esp{A Juan lo vi ayer.} « Jean, je l'ai vu hier. » (littéralement : « À Jean, je l'ai vu hier »). Le complément anticipé est alors repris par un pronom (\esp{lo}, \esp{la}, \esp{le}...). Autre exemple : \esp{El fútbol lo juegan mis hermanos todos los sábados.} « Le football, mes frères y jouent tous les samedis. » En français, nous devons utiliser une virgule et une reprise ; l'espagnol fait tout cela dans la phrase, naturellement. C'est pourquoi l'espagnol est souvent plus souple que le français pour l'emphase.
\end{savaistu}

\begin{exercice}[À vous d'observer]
Dans le texte d'Aïcha (début de la section), retrouvez trois tournures d'emphase que nous n'avons pas encore commentées. Pour chacune : 1) recopiez-la ; 2) nommez l'élément mis en relief ; 3) proposez une traduction française avec « c'est... qui/que » ou « ne... que ».
\end{exercice}

\begin{corrige}[Exercice : À vous d'observer]
Réponses possibles :\par
1. \esp{Es ella quien cuida de la casa} : mise en relief du sujet \esp{ella} ; « C'est elle qui s'occupe de la maison. »\par
2. \esp{Es a él a quien admiro más} : mise en relief du COD de personne \esp{a él} ; « C'est lui que j'admire le plus. »\par
3. \esp{No tengo más que cariño por ellos} : restriction ; « Je n'ai que de l'affection pour eux. » On note ici la variante \esp{no... más que}, équivalente de \esp{solo}.
\end{corrige}

\begin{aretenir}[Bilan de l'observation]
\begin{itemize}
\item L'\cle{emphase} met en relief un élément de la phrase (sujet, COD, circonstance, restriction).
\item Le français clive avec \textbf{c'est... qui} (sujet) et \textbf{c'est... que} (complément) ; l'espagnol avec \esp{es... quien}, \esp{es... a quien}, \esp{es... lo que}, \esp{fue... cuando/donde}.
\item « Ne... que » se traduit par \esp{solo}, \esp{únicamente}, \esp{no... más que} ou \esp{no... sino}.
\item L'espagnol peut aussi créer l'emphase par le seul ordre des mots : \esp{A Juan lo vi ayer}.
\item Règle d'or : repérez d'abord \emph{ce qui est mis en relief} en français, puis choisissez le mot-outil espagnol correspondant.
\end{itemize}
\end{aretenir}

\coursec{Traduire « c'est… qui »}

Dans la section précédente, nous avons \emph{observé} comment l'espagnol, comme le français, aime mettre un élément de la phrase \emph{en relief} : au lieu de dire simplement \esp{María canta}, on peut dire \esp{Es María quien canta}, exactement comme le français dit « C'est Marie qui chante ». Nous allons maintenant étudier de manière systématique la première de ces tournures : la phrase clivée avec \esp{qui}. C'est la structure la plus fréquente à l'examen, et aussi celle où les francophones commettent le plus d'erreurs, car le français et l'espagnol se ressemblent… sans être identiques.

\courssub{Observation : un dialogue en famille}

Lisons d'abord ce petit dialogue original. Repérez les phrases en gras : quel élément est mis en relief ? Comment l'espagnol traduit-il le « qui » français ?

\begin{textebox}[Un domingo en familia]
— ¿Quién ha roto el jarrón del salón?\\
— \textbf{Es Pedro quien lo ha roto}, mamá. Yo estaba en el patio.\\
— ¡Mientes! \textbf{Fue tu hermana quien lo rompió} ayer por la tarde.\\
— No, abuela, \textbf{son los niños de los vecinos quienes juegan} siempre cerca de la ventana.\\
— Bueno, basta de discusiones. \textbf{Es el abuelo quien decide} en esta casa, y él dice que nadie sale al cine esta noche.\\
— ¡Qué injusto! \textbf{Soy yo quien siempre ayuda} en la cocina, y nunca me recompensan.\\
— Tranquilo, hijo. \textbf{Es tu madre quien tiene la última palabra}, no yo.
\end{textebox}

\textbf{Traduction :} « — Qui a cassé le vase du salon ? — C'est Pedro qui l'a cassé, maman. Moi, j'étais dans la cour. — Tu mens ! C'est ta sœur qui l'a cassé hier après-midi. — Non, grand-mère, ce sont les enfants des voisins qui jouent toujours près de la fenêtre. — Bon, assez de disputes. C'est le grand-père qui décide dans cette maison, et il dit que personne ne va au cinéma ce soir. — C'est injuste ! C'est moi qui aide toujours à la cuisine, et on ne me récompense jamais. — Calme-toi, mon fils. C'est ta mère qui a le dernier mot, pas moi. »

\textbf{Glossaire :} \esp{el jarrón} : le vase ; \esp{el patio} : la cour ; \esp{mentir} (mientes) : mentir ; \esp{romper} : casser ; \esp{los vecinos} : les voisins ; \esp{basta de} : assez de ; \esp{recompensar} : récompenser ; \esp{la última palabra} : le dernier mot.

\textbf{Questions d'observation :}
\begin{enumerate}
\item Soulignez l'élément mis en relief dans chaque phrase en gras : s'agit-il d'une personne ou d'une chose ?
\item Quel mot espagnol correspond au « qui » français ? Change-t-il selon le nombre (singulier/pluriel) ?
\item Observez le verbe \esp{ser} : quelles formes apparaissent (\esp{es}, \esp{son}, \esp{fue}, \esp{soy}) ? De quoi dépendent-elles ?
\end{enumerate}

\textbf{Réponses attendues :} 1) Tous les éléments mis en relief sont des personnes (Pedro, tu hermana, los niños, el abuelo, yo, tu madre). 2) Le « qui » français se traduit par \esp{quien} au singulier et \esp{quienes} au pluriel. 3) \esp{Ser} apparaît sous les formes \esp{es} (singulier), \esp{son} (pluriel), \esp{fue} (passé, singulier) et \esp{soy} (avec « yo ») : il s'accorde avec l'élément mis en relief.

\courssub{Cas 1 : l'élément mis en relief est une personne (sujet animé)}

Lorsque l'élément mis en relief est une \cle{personne} (ou un être animé), l'espagnol n'utilise jamais \esp{que} seul : il exige le pronom relatif \esp{quien} (singulier) ou \esp{quienes} (pluriel), ou bien la tournure avec l'article + \esp{que}.

\begin{propriete}[« C'est + personne + qui »]
\begin{itemize}
\item \textbf{Singulier :} \esp{Es + personne + quien + verbe} \\
\esp{Es María quien canta.} « C'est Marie qui chante. »
\item \textbf{Pluriel :} \esp{Son + personnes + quienes + verbe} \\
\esp{Son mis hermanos quienes llegan.} « Ce sont mes frères qui arrivent. »
\item \textbf{Variante équivalente :} \esp{es + el que / la que} (singulier), \esp{son + los que / las que} (pluriel) \\
\esp{Es María la que canta.} / \esp{Son mis hermanos los que llegan.}
\end{itemize}
Le verbe \esp{ser} s'accorde avec l'élément mis en relief : \esp{es} au singulier, \esp{son} au pluriel.
\end{propriete}

\begin{attention}[Quien ou quién : l'accent fait la différence]
Le pronom relatif \esp{quien} (celui qui) s'écrit \textbf{sans accent} : \esp{Es María quien canta.} L'accentué \esp{quién} n'existe que dans les interrogations directes et indirectes : \esp{¿Quién canta?} « Qui chante ? » ; \esp{No sé quién canta.} « Je ne sais pas qui chante. » Ne confondez pas les deux à l'écrit.
\end{attention}

\begin{exemplebox}[Exemples traduits]
\begin{itemize}
\item \esp{Es mi padre quien trabaja en el mercado de Yaoundé.} « C'est mon père qui travaille au marché de Yaoundé. »
\item \esp{Es la profesora quien corrige los exámenes.} « C'est la professeure qui corrige les examens. »
\item \esp{Son los alumnos de Terminale quienes organizan la fiesta.} « Ce sont les élèves de Terminale qui organisent la fête. »
\item \esp{Son mis primas las que viven en Douala.} « Ce sont mes cousines qui habitent à Douala. »
\item \esp{Fue el médico quien me curó.} « C'est le médecin qui m'a soigné. » (passé)
\end{itemize}
\end{exemplebox}

\begin{attention}[Le piège n° 1 du francophone]
En français, « qui » ne change jamais. En espagnol, après une \textbf{personne} mise en relief, on ne peut \textbf{jamais} dire \esp{es mi padre que trabaja}. C'est une faute grave. Il faut obligatoirement : \esp{Es mi padre \textbf{quien} trabaja} ou \esp{Es mi padre \textbf{el que} trabaja}. Retenez : \textbf{personne → quien / el que ; jamais « que » seul}.
\end{attention}

\courssub{Cas 2 : l'élément mis en relief est une chose (sujet inanimé)}

Lorsque l'élément mis en relief est une \cle{chose}, une idée ou un animal, \esp{quien} est impossible. On utilise alors \esp{el que} (avec l'article qui s'accorde en genre et en nombre avec la chose) ou, pour une idée globale, \esp{lo que}.

\begin{propriete}[« C'est + chose + qui »]
\begin{itemize}
\item \textbf{Chose déterminée :} \esp{Es + chose + el que / la que / los que / las que + verbe} \\
\esp{Es el sol el que me despierta.} « C'est le soleil qui me réveille. »
\item \textbf{Idée, proposition :} \esp{Es + lo que + verbe} \\
\esp{Es lo que me preocupa.} « C'est ce qui m'inquiète. »
\item La tournure \esp{es… que} est parfois entendue dans la langue courante pour les choses, surtout en Amérique latine, mais à l'examen préférez toujours \esp{el que}, qui est la norme soignée.
\end{itemize}
\end{propriete}

\begin{exemplebox}[Exemples traduits]
\begin{itemize}
\item \esp{Es el ruido el que no me deja dormir.} « C'est le bruit qui m'empêche de dormir. »
\item \esp{Es la lluvia la que destruye las cosechas de cacao.} « C'est la pluie qui détruit les récoltes de cacao. »
\item \esp{Son los bendskins los que causan muchos accidentes.} « Ce sont les bendskins qui causent beaucoup d'accidents. »
\item \esp{Fue tu carta la que me alegró el día.} « C'est ta lettre qui m'a réjoui la journée. »
\end{itemize}
\end{exemplebox}

\courssub{Accord du verbe ser et concordance des temps}

Deux questions se posent à chaque traduction : \textbf{singulier ou pluriel ?} et \textbf{présent ou passé ?}

\begin{propriete}[Accord et concordance]
\begin{itemize}
\item Le verbe \esp{ser} s'accorde en \textbf{nombre} avec l'élément mis en relief : \esp{Es Juan…} / \esp{Son Juan y Pedro…}
\item Il s'accorde en \textbf{temps} avec le moment de l'énonciation : présent \esp{es/son} ; passé \esp{fue/fueron} (prétérit, action ponctuelle) ou \esp{era/eran} (imparfait, habitude ou description).
\item Le second verbe garde le temps logique de la phrase : \esp{Fue ayer cuando llegó.} « C'est hier qu'il est arrivé. »
\end{itemize}
\end{propriete}

\begin{center}
\begin{tabularx}{\linewidth}{|X|X|X|}
\hline
\rowcolor{popblueL} Français & Español (quien) & Español (el que) \\
\hline
C'est Marie qui chante. & Es María quien canta. & Es María la que canta. \\
\hline
Ce sont mes frères qui arrivent. & Son mis hermanos quienes llegan. & Son mis hermanos los que llegan. \\
\hline
C'est le soleil qui me réveille. & --- (inanimé) & Es el sol el que me despierta. \\
\hline
C'était mon oncle qui téléphonait. & Era mi tío quien llamaba. & Era mi tío el que llamaba. \\
\hline
Ce furent ses paroles qui me touchèrent. & --- (inanimé) & Fueron sus palabras las que me conmovieron. \\
\hline
\end{tabularx}
\end{center}

\begin{popfigure}
\begin{tikzpicture}[node distance=0.5cm and 1.6cm,
  box/.style={draw=popblue, fill=popblueL, rounded corners, align=center, minimum width=3.1cm, minimum height=1cm, font=\small},
  dec/.style={draw=poporange, fill=poporangeL, rounded corners, align=center, minimum width=3.4cm, minimum height=1cm, font=\small},
  arr/.style={-{Stealth[length=3mm]}, thick, popdark}]
\node[box, align=center] (start){Élément mis\\en relief};
\node[dec, below left=0.6cm and 0.2cm of start, align=center] (anim){Animé\\(personne)};
\node[dec, below right=0.6cm and 0.2cm of start, align=center] (inanim){Inanimé\\(chose, idée)};
\node[box, below=0.6cm of anim, fill=popgreenL, draw=popgreen, align=center] (r1){es + quien\\son + quienes\\(ou el que / los que)};
\node[box, below=0.6cm of inanim, fill=popgreenL, draw=popgreen, align=center] (r2){es + el que / la que\\es + lo que\\(jamais quien)};
\draw[arr] (start) -- (anim);
\draw[arr] (start) -- (inanim);
\draw[arr] (anim) -- (r1);
\draw[arr] (inanim) -- (r2);
\end{tikzpicture}
\legende{Schéma de décision : comment traduire « c'est… qui » selon la nature de l'élément mis en relief.}
\end{popfigure}

\courssub{Cas particulier : « c'est moi qui… », « c'est nous qui… »}

Quand l'élément mis en relief est un \cle{pronom personnel} (moi, toi, lui, nous, vous, eux), le verbe \esp{ser} se conjugue \textbf{selon la personne} du pronom, et non plus à la troisième personne. C'est une différence majeure avec le français, où « c'est » reste invariable.

\begin{propriete}[« C'est moi qui… » → « Soy yo quien… »]
\begin{center}
\begin{tabular}{|l|l|}
\hline
\rowcolor{popblueL} Français & Español \\
\hline
C'est moi qui… & Soy yo quien… \\
\hline
C'est toi qui… & Eres tú quien… \\
\hline
C'est lui/elle qui… & Es él/ella quien… \\
\hline
C'est nous qui… & Somos nosotros quienes… \\
\hline
C'est vous qui… & Sois vosotros / Son ustedes quienes… \\
\hline
C'est eux qui… & Son ellos quienes… \\
\hline
\end{tabular}
\end{center}
Le second verbe se met à la 3\up{e} personne après \esp{quien}, comme en français : \esp{Soy yo quien \textbf{ha ganado}} « c'est moi qui \emph{a} gagné ».
\end{propriete}

\begin{exemplebox}[Exemples traduits]
\begin{itemize}
\item \esp{Soy yo quien ha ganado el partido.} « C'est moi qui ai gagné le match. »
\item \esp{Somos nosotros quienes decidimos.} « C'est nous qui décidons. »
\item \esp{Eres tú quien debe pedir perdón.} « C'est toi qui dois demander pardon. »
\item \esp{Son ellos quienes pagan la fiesta.} « Ce sont eux qui paient la fête. »
\item \esp{Fuiste tú quien me avisó.} « C'est toi qui m'as prévenu. » (passé)
\end{itemize}
\end{exemplebox}

\begin{attention}[Deux pièges avec les pronoms]
\begin{enumerate}
\item Ne dites jamais \esp{es yo quien…} : le verbe \esp{ser} doit se conjuguer → \esp{\textbf{Soy} yo quien…}. Le français « c'est » invariable est un faux ami de structure.
\item Après \esp{quien}, le verbe se met à la 3\up{e} personne même avec « moi » : \esp{Soy yo quien \textbf{ha} ganado} (et non \esp{he ganado}), exactement comme en français « c'est moi qui \emph{a} gagné ».
\end{enumerate}
\end{attention}

\begin{savaistu}[Pourquoi « quien » et pas « que » ?]
En espagnol, le pronom relatif \esp{que} renvoie de préférence aux choses ; pour les personnes, la langue soignée préfère \esp{quien}, hérité du latin \emph{quem}. C'est pourquoi la phrase clivée, qui met l'accent sur la personne, exige \esp{quien}. En Guinée équatoriale, seul pays hispanophone d'Afrique, on entend exactement les mêmes tournures : \esp{Es mi hermano quien habla francés} « C'est mon frère qui parle français ».
\end{savaistu}

\courssub{Méthode et entraînement}

\begin{methode}[Traduire « c'est… qui » en quatre étapes]
\begin{enumerate}
\item \textbf{Identifier} l'élément mis en relief (ce qui se trouve entre « c'est » et « qui »).
\item \textbf{Animé ou inanimé ?} Personne → \esp{quien/quienes} ; chose → \esp{el que/la que} ; idée → \esp{lo que}.
\item \textbf{Singulier ou pluriel ?} → \esp{es} ou \esp{son} ; \esp{quien} ou \esp{quienes} ; \esp{el que} ou \esp{los que}.
\item \textbf{Temps de l'énonciation ?} Présent → \esp{es/son} ; passé → \esp{fue/fueron} ; pronom personnel → conjuguer \esp{ser} selon la personne (\esp{soy yo}, \esp{somos nosotros}…).
\end{enumerate}
\end{methode}

\begin{exoresolu}[Thème guidé]
Traduisez en espagnol : « Ce sont mes grands-parents qui m'ont élevé. »
\tcblower
\textbf{Étape 1 :} élément mis en relief = « mes grands-parents ». \textbf{Étape 2 :} personnes → \esp{quienes} (pluriel). \textbf{Étape 3 :} pluriel → \esp{son}. \textbf{Étape 4 :} l'énonciation est au présent, le second verbe au passé composé → \esp{han criado}.\\
\textbf{Solution :} \esp{Son mis abuelos quienes me han criado.} (variante correcte : \esp{Son mis abuelos los que me han criado.})
\end{exoresolu}

\begin{exoresolu}[Thème : phrases à traduire]
Traduisez : 1) « C'est Aïcha qui prépare le repas. » 2) « C'est le vent qui fait tomber les mangues. » 3) « C'est nous qui avons gagné le match de football. » 4) « C'était les voisins qui chantaient toute la nuit. »
\tcblower
1) Personne, singulier, présent : \esp{Es Aïcha quien prepara la comida.}\\
2) Chose, singulier : \esp{Es el viento el que hace caer los mangos.} (\esp{quien} impossible : inanimé.)\\
3) Pronom « nous » → \esp{ser} conjugué : \esp{Somos nosotros quienes hemos ganado el partido de fútbol.}\\
4) Personnes, pluriel, imparfait : \esp{Eran los vecinos quienes cantaban toda la noche.}
\end{exoresolu}

\begin{exercice}[À vous de jouer]
Traduisez en espagnol :
\begin{enumerate}
\item C'est mon frère aîné qui conduit le taxi-brousse.
\item Ce sont les femmes du marché qui vendent les meilleurs légumes.
\item C'est la télévision qui passionne mes petits frères.
\item C'est moi qui ai réparé la radio.
\item C'était son père qui travaillait à la plantation de cacao.
\end{enumerate}
\end{exercice}

\begin{corrige}[Exercice « À vous de jouer »]
\begin{enumerate}
\item \esp{Es mi hermano mayor quien conduce el taxi colectivo.} (personne, singulier)
\item \esp{Son las mujeres del mercado quienes venden las mejores verduras.} (personnes, pluriel)
\item \esp{Es la televisión la que apasiona a mis hermanitos.} (chose → \esp{la que} ; attention à la \emph{a} personnelle devant \esp{mis hermanitos})
\item \esp{Soy yo quien ha arreglado la radio.} (pronom « moi » → \esp{soy})
\item \esp{Era su padre quien trabajaba en la plantación de cacao.} (imparfait : habitude passée)
\end{enumerate}
\end{corrige}

\begin{aretenir}[L'essentiel de la section]
\begin{itemize}
\item Personne mise en relief : \esp{es/son + quien/quienes} (ou \esp{el que/los que}) — jamais \esp{que} seul.
\item Chose mise en relief : \esp{es + el que / la que} ; idée : \esp{lo que}.
\item \esp{Ser} s'accorde avec l'élément mis en relief en nombre (\esp{es/son}) et en temps (\esp{es/fue/era}).
\item Avec un pronom : \esp{soy yo quien…}, \esp{somos nosotros quienes…}
\item \esp{Quien} relatif s'écrit sans accent ; \esp{quién} avec accent seulement dans l'interrogation.
\end{itemize}
\end{aretenir}

\coursec{Traduire « c'est… que »}

Dans la section précédente, nous avons appris à traduire « c'est… qui » lorsque l'élément mis en relief est le \emph{sujet} de la phrase : \esp{Es mi hermana quien cocina} (« C'est ma sœur qui cuisine »). Mais que se passe-t-il lorsque l'élément mis en relief n'est \emph{pas} le sujet ? Le français dit alors « c'est… \emph{que} » : « C'est hier qu'il est arrivé », « C'est à Douala que j'habite », « C'est ce livre que je cherche ». L'espagnol, lui, choisit un relatif différent selon la \emph{nature} de l'élément mis en relief. C'est toute la logique de cette section : observer, classer, puis appliquer.

\courssub{Observation : quatre phrases, quatre relatifs}

\begin{textebox}[Un mensaje de Amina a su prima]
Hola, prima:\\
¿Sabes qué? Fue ayer cuando llegó mi hermano mayor de España. Es en el aeropuerto de Yaundé donde lo esperábamos toda la familia. Es a él a quien mi abuela quiere más, ¡qué suerte tiene! Y es el regalo que trajo lo que más me gustó: un libro de fotos de Madrid. Lo que me sorprende es que habla español casi sin acento. Te cuento más mañana.\\
Un abrazo, Amina.
\end{textebox}

\textbf{Glossaire :} \esp{el aeropuerto} : l'aéroport ; \esp{lo esperábamos} : nous l'attendions ; \esp{el regalo} : le cadeau ; \esp{me sorprende} : me surprend ; \esp{el acento} : l'accent ; \esp{un abrazo} : une embrassade.

Observons les quatre tournures d'emphase du message :

\begin{center}
\begin{tabularx}{\linewidth}{|X|X|X|}\hline
\rowcolor{popblueL} Élément mis en relief & Español & Français \\ \hline
temps (ayer) & \esp{Fue ayer cuando llegó.} & C'est hier qu'il est arrivé. \\ \hline
lieu (en el aeropuerto) & \esp{Es en el aeropuerto donde lo esperábamos.} & C'est à l'aéroport que nous l'attendions. \\ \hline
personne, COD (a él) & \esp{Es a él a quien mi abuela quiere más.} & C'est lui que ma grand-mère préfère. \\ \hline
chose, COD (el regalo) & \esp{Es el regalo lo que más me gustó.} & C'est le cadeau qui m'a le plus plu. \\ \hline
\end{tabularx}
\end{center}

La leçon est claire : le français utilise partout le même mot « que », mais l'espagnol \emph{adapte le relatif} : \cle{cuando} pour le temps, \cle{donde} pour le lieu, \cle{a quien} pour la personne, \cle{lo que} ou \cle{que} pour la chose.

\courssub{Cas 1 : le COD mis en relief}

\textbf{a) Le COD est une personne.}

\begin{propriete}[« C'est + personne + que » (COD)]
Quand l'élément mis en relief est un \emph{complément d'objet direct} désignant une \emph{personne}, on traduit par : \esp{Es a + personne + a quien (ou que) + verbe}. La préposition \esp{a} (l'« a » personnel) est obligatoire devant la personne. Au pluriel, on emploie \esp{a quienes} : \esp{Es a mis hermanos a quienes espero.} (« C'est mes frères que j'attends. »)
\end{propriete}

\begin{exemplebox}[Exemples traduits]
\esp{Es a ti a quien vi en el mercado.} « C'est toi que j'ai vu(e) au marché. »

\esp{Es a mi padre a quien debo llamar.} « C'est mon père que je dois appeler. »

\esp{Es a la profesora a quien todos respetan.} « C'est la professeure que tout le monde respecte. »
\end{exemplebox}

\textbf{b) Le COD est une chose.}

\begin{propriete}[« C'est + chose + que » (COD)]
Quand l'élément mis en relief est une \emph{chose}, on traduit par : \esp{Es + chose + lo que + verbe} (tournure la plus élégante) ou \esp{Es + chose + que + verbe} (plus simple, admise).
\end{propriete}

\begin{exemplebox}[Exemples traduits]
\esp{Es este libro lo que busco.} « C'est ce livre que je cherche. »

\esp{Es el cacao lo que exporta la región.} « C'est le cacao que la région exporte. »

\esp{Es tu ayuda lo que necesito.} « C'est ton aide qu'il me faut. »
\end{exemplebox}

\courssub{Cas 2 : le complément de lieu mis en relief}

\begin{propriete}[« C'est + lieu + que »]
Quand l'élément mis en relief est un \emph{complément de lieu}, on traduit par : \esp{Es (en) + lieu + donde + verbe}. Le relatif est \cle{donde}, jamais \esp{que}. On garde la préposition du lieu (\esp{en}, \esp{a}) devant le nom de lieu, et non devant \esp{donde}.
\end{propriete}

\begin{exemplebox}[Exemples traduits]
\esp{Es en Yaundé donde vivo.} « C'est à Yaoundé que j'habite. »

\esp{Es en Duala donde está el puerto más grande del país.} « C'est à Douala que se trouve le plus grand port du pays. »

\esp{Fue en Malabo donde se conocieron mis padres.} « C'est à Malabo que mes parents se sont rencontrés. »
\end{exemplebox}

\courssub{Cas 3 : le complément de temps mis en relief}

\begin{propriete}[« C'est + temps + que »]
Quand l'élément mis en relief est un \emph{complément de temps}, on traduit par : \esp{Fue + temps + cuando + verbe} (si la subordonnée est au passé) ou \esp{Es + temps + cuando + verbe} (présent ou futur). Le relatif est \cle{cuando}, jamais \esp{que}.
\end{propriete}

\begin{exemplebox}[Exemples traduits]
\esp{Fue ayer cuando llegó.} « C'est hier qu'il est arrivé. »

\esp{Fue en 1960 cuando el Camerún se independizó.} « C'est en 1960 que le Cameroun est devenu indépendant. »

\esp{Es mañana cuando empiezan los exámenes.} « C'est demain que commencent les examens. »
\end{exemplebox}

\courssub{Cas 4 : « Ce… que » en tête de phrase}

Quand la phrase française commence par « Ce que…, c'est… », l'espagnol utilise \cle{lo que} en tête, suivi de \esp{ser} :

\begin{exemplebox}[Exemples traduits]
\esp{Lo que quiero es aprobar el examen.} « Ce que je veux, c'est réussir l'examen. »

\esp{Lo que me gusta es el fútbol.} « Ce que j'aime, c'est le football. »

\esp{Lo que dice mi abuela es siempre verdad.} « Ce que dit ma grand-mère est toujours vrai. »
\end{exemplebox}

\begin{attention}[Piège du francophone]
Ne traduis jamais mécaniquement le « que » français par \esp{que} ! Devant un \emph{temps}, il faut \esp{cuando} ; devant un \emph{lieu}, \esp{donde} ; devant une \emph{personne-COD}, \esp{a quien}. Dire \esp{*Es ayer que llegó} est une erreur typique : la bonne forme est \esp{Fue ayer cuando llegó}. Autre piège : ne pas oublier l'« a » personnel devant la personne mise en relief : \esp{Es a ti a quien vi}, et non \esp{*Es ti a quien vi}.
\end{attention}

\courssub{Es ou Fue ? L'accord des temps}

\begin{propriete}[Choix entre es et fue]
Le verbe \esp{ser} de la tournure clivée s'accorde en temps avec le verbe de la subordonnée :
\begin{itemize}
\item subordonnée au \emph{présent} ou au \emph{futur} → \esp{Es… donde/cuando/lo que…}
\item subordonnée au \emph{passé} → \esp{Fue… donde/cuando/lo que…}
\end{itemize}
\end{propriete}

\esp{Es en Yaundé donde vivo.} (présent) — \esp{Fue en Yaundé donde nací.} (passé : « C'est à Yaoundé que je suis né. »)

\courssub{Organigramme de décision}

\begin{popfigure}
\begin{tikzpicture}[
  node distance=0.55cm and 0.7cm,
  startstop/.style={rectangle, rounded corners, draw=popdark, fill=popblueL, align=center, font=\small, inner sep=5pt},
  decision/.style={rectangle, draw=poporange, fill=poporangeL, align=center, font=\small, inner sep=5pt},
  result/.style={rectangle, rounded corners, draw=popgreen, fill=popgreenL, align=center, font=\small, inner sep=5pt},
  arr/.style={-{Stealth[length=2.5mm]}, thick, popdark}
]
\node[startstop] (q) {Quel élément est mis en relief ?};
\node[decision, below left=0.8cm and 2.2cm of q, align=center] (per){une personne\\(COD)};
\node[decision, below left=0.8cm and -0.4cm of q, align=center] (cho){une chose\\(COD)};
\node[decision, below right=0.8cm and -0.4cm of q] (lieu) {un lieu};
\node[decision, below right=0.8cm and 2.2cm of q] (temps) {un temps};
\node[result, below=0.6cm of per, align=center] (r1){\esp{Es a + pers.}\\\esp{a quien / que}};
\node[result, below=0.6cm of cho, align=center] (r2){\esp{Es + chose}\\\esp{lo que / que}};
\node[result, below=0.6cm of lieu, align=center] (r3){\esp{Es/Fue en + lieu}\\\esp{donde}};
\node[result, below=0.6cm of temps, align=center] (r4){\esp{Fue + temps}\\\esp{cuando}};
\draw[arr] (q) -| (per);
\draw[arr] (q) -- (cho);
\draw[arr] (q) -- (lieu);
\draw[arr] (q) -| (temps);
\draw[arr] (per) -- (r1);
\draw[arr] (cho) -- (r2);
\draw[arr] (lieu) -- (r3);
\draw[arr] (temps) -- (r4);
\end{tikzpicture}
\legende{Choisir le bon relatif emphatique selon l'élément mis en relief.}
\end{popfigure}

\courssub{Tableau récapitulatif des relatifs emphatiques}

\begin{center}
\begin{tabularx}{\linewidth}{|l|X|X|X|}\hline
\rowcolor{popblueL} Élément & Français & Español & Exemple \\ \hline
personne (COD) & c'est… que & \esp{a quien / que} & \esp{Es a ti a quien vi.} \\ \hline
chose (COD) & c'est… que & \esp{lo que / que} & \esp{Es este libro lo que busco.} \\ \hline
lieu & c'est… que & \esp{donde} & \esp{Es en Yaundé donde vivo.} \\ \hline
temps & c'est… que & \esp{cuando} & \esp{Fue ayer cuando llegó.} \\ \hline
« ce que…, c'est » & ce que… c'est & \esp{lo que… es} & \esp{Lo que quiero es aprobar.} \\ \hline
\end{tabularx}
\end{center}

\begin{aretenir}[L'essentiel]
« C'est… que » ne se traduit jamais mot à mot. On identifie l'élément mis en relief (personne, chose, lieu, temps), on choisit le relatif correspondant (\esp{a quien}, \esp{lo que}, \esp{donde}, \esp{cuando}), et l'on accorde \esp{ser} au temps de la subordonnée (\esp{es} au présent, \esp{fue} au passé).
\end{aretenir}

\begin{exoresolu}[Traduis en espagnol]
Énoncé : 1. C'est hier que le match a commencé. 2. C'est à Douala que ma tante travaille. 3. C'est toi que je cherchais. 4. Ce que j'aime, c'est la musique camerounaise.
\tcblower
Solution détaillée :

1. Élément mis en relief : « hier » = temps ; subordonnée au passé → \esp{Fue ayer cuando empezó el partido.}

2. Élément : « à Douala » = lieu ; subordonnée au présent → \esp{Es en Duala donde trabaja mi tía.}

3. Élément : « toi » = personne-COD → \esp{Es a ti a quien buscaba.} (on garde l'imparfait \esp{buscaba}).

4. « Ce que…, c'est… » en tête → \esp{Lo que me gusta es la música camerunesa.}
\end{exoresolu}

\begin{exercice}[À toi de traduire]
Traduis en espagnol : 1. C'est en 1972 que le pays a changé de nom. 2. C'est ce marché que je préfère. 3. C'est à l'école que nous avons appris l'espagnol. 4. C'est demain que part le bus. 5. Ce que dit le professeur est important.
\end{exercice}

\begin{corrige}[Exercice : traduction de « c'est… que »]
1. \esp{Fue en 1972 cuando el país cambió de nombre.} (temps + passé → \esp{fue… cuando}).

2. \esp{Es este mercado lo que prefiero.} (chose-COD → \esp{lo que}).

3. \esp{Es en la escuela donde aprendimos español.} (lieu → \esp{donde}).

4. \esp{Es mañana cuando sale el autobús.} (temps + présent → \esp{es… cuando}).

5. \esp{Lo que dice el profesor es importante.} (« ce que… » en tête → \esp{lo que}).
\end{corrige}

\begin{savaistu}[Une tournure qui rapporte des points]
Au baccalauréat, les correcteurs des épreuves de thème valorisent particulièrement les tournures clivées bien construites : \esp{Fue ayer cuando…}, \esp{Es en Duala donde…}, \esp{Lo que quiero es…}. Une seule de ces structures, correctement employée, montre que tu maîtrises l'emphase espagnole — un marqueur de très bon niveau. Entraîne-toi à en placer une dans chaque thème !
\end{savaistu}

\coursec{Traduire « ne… que »}

Après avoir appris à mettre en relief un élément de la phrase avec les tournures clivées (\esp{Es Ana quien canta}, \esp{Fue ayer cuando llegó}), nous abordons une autre structure très fréquente du français : la \cle{restriction} exprimée par « ne… que ». Cette tournure, si naturelle en français, est un piège classique pour les hispanisants, car l'espagnol ne la traduit \emph{jamais} mot à mot. Observons d'abord quelques phrases avant de dégager les règles.

\courssub{Observation : découvrir les solutions espagnoles}

\begin{textebox}[Un diálogo en el mercado de Yaoundé]
— ¿Cuántos hijos tienes, señora Mballa?\\
— Solo tengo dos: un niño y una niña. No tengo más que dos, ¡y ya es bastante trabajo!\\
— ¿Y tu hija mayor? ¿Qué hace?\\
— No hace más que estudiar. No quiere ser comerciante, sino médica.\\
— ¡Qué bien! ¿Y el pequeño?\\
— Él no piensa más que en el fútbol. No solo juega en el barrio, sino también en el equipo del liceo.\\
— ¿Vienes al mercado todos los días?\\
— No, únicamente los sábados. Los otros días no vendo más que cacao en mi aldea.
\end{textebox}

\textbf{Glossaire :} \esp{el barrio} = le quartier ; \esp{el equipo} = l'équipe ; \esp{vender} = vendre ; \esp{la aldea} = le village ; \esp{el trabajo} = le travail.

\textbf{Questions d'observation :}
\begin{enumerate}
\item Relevez toutes les phrases du dialogue qui traduisent un « ne… que » français.
\item Combien de structures espagnoles différentes observez-vous ? Lesquelles ?
\item Dans « \esp{No quiere ser comerciante, sino médica} », s'agit-il d'une simple restriction ? Justifiez.
\end{enumerate}

On constate que l'espagnol dispose de \textbf{trois solutions principales} pour rendre « ne… que », plus une construction spéciale avec un verbe. Étudions-les une à une.

\courssub{Solution 1 : la restriction simple avec \esp{solo}, \esp{solamente}, \esp{únicamente}}

\begin{propriete}[Restriction simple : \esp{solo} avant l'élément restreint]
Pour traduire un « ne… que » de \textbf{simple restriction} (le cas le plus fréquent), l'espagnol supprime la négation et place \cle{solo} (ou ses synonymes \esp{solamente}, \esp{únicamente}) \textbf{juste avant l'élément restreint}, exactement comme « que » en français.
\begin{center}
\begin{tabularx}{\linewidth}{|X|X|X|}
\hline
\rowcolor{popblueL} Français & Español & Remarque \\
\hline
ne + verbe + \textbf{que} + élément & verbe affirmatif + \textbf{solo} + élément & Pas de négation \esp{no} \\
\hline
\end{tabularx}
\end{center}
\esp{Solo} est invariable. \esp{Solamente} et \esp{únicamente} sont des variantes plus soutenues ou plus insistantes, parfaitement interchangeables avec \esp{solo} dans cet emploi.
\end{propriete}

\begin{exemplebox}[Exemples traduits]
\esp{Solo tengo diez años.} « Je n'ai que dix ans. »\\
\esp{Únicamente quedan tres plazas.} « Il ne reste que trois places. »\\
\esp{Solamente hablo español en clase.} « Je ne parle qu'espagnol en classe. »\\
\esp{El autobús solo pasa los lunes.} « Le car ne passe que le lundi. »\\
\esp{Solo conozco a su hermano.} « Je ne connais que son frère. »
\end{exemplebox}

Remarquez la correspondance des positions : dans « Je ne parle \textbf{qu}'espagnol », le « que » précède immédiatement « espagnol » ; dans \esp{solo hablo español}, \esp{solo} occupe la même place relative. C'est le meilleur réflexe à acquérir : \textbf{repérer ce que « que » restreint, puis placer \esp{solo} juste devant}.

\begin{attention}[Le « que » de « ne… que » ne se traduit jamais seul]
Ne dites jamais \esp{no tengo que dos hermanos} : c'est un calque du français, incompréhensible en espagnol. La négation \esp{no} et la restriction ne cohabitent pas dans cette solution. On choisit : soit \esp{Solo tengo dos hermanos}, soit \esp{No tengo más que dos hermanos} (voir ci-dessous), mais jamais un mélange des deux.
\end{attention}

\courssub{Solution 2 : la restriction insistante avec \esp{no… más que}}

\begin{propriete}[Restriction forte : \esp{no… más que}]
Pour insister sur la restriction, surtout au \textbf{registre soutenu} ou pour exprimer une nuance de regret, de plainte ou d'évidence, l'espagnol emploie la tournure négative \cle{no… más que} (littéralement « ne… pas plus que »). La négation \esp{no} se place devant le verbe, et \esp{más que} devant l'élément restreint.
\end{propriete}

\begin{exemplebox}[Exemples traduits]
\esp{No tengo más que diez años.} « Je n'ai que dix ans. » (insistance : « à peine dix ans »)\\
\esp{No hay más que Dios.} « Il n'y a que Dieu. »\\
\esp{No me quedan más que cinco mil francos.} « Il ne me reste que cinq mille francs. »\\
\esp{No dijo más que la verdad.} « Il n'a dit que la vérité. »\\
\esp{No quiero más que tu felicidad.} « Je ne veux que ton bonheur. »

La différence entre \esp{Solo tengo diez años} et \esp{No tengo más que diez años} est une différence de \textbf{ton}, non de sens : la seconde phrase souligne davantage la petitesse du nombre. Au baccalauréat, les deux sont acceptées ; la seconde montre une maîtrise plus fine de la langue.

\end{exemplebox}

\courssub{Solution 3 : l'opposition avec \esp{no… sino}}

\begin{propriete}[Opposition et correction : \esp{no… sino}]
Quand le français « ne… que » cache en réalité une \textbf{opposition} ou une \textbf{correction} (« pas X, mais Y »), l'espagnol emploie \cle{no… sino}. La première proposition est niée, la seconde la corrige ou s'y substitue. Si le second élément contient un verbe conjugué, on utilise \esp{sino que}.
\end{propriete}

\begin{exemplebox}[Exemples traduits]
\esp{No bebe vino, sino agua.} « Il ne boit pas de vin, mais de l'eau. »\\
\esp{No son tres, sino cuatro.} « Ce ne sont pas trois, mais quatre. »\\
\esp{No vive en Douala, sino en Bafoussam.} « Il n'habite pas à Douala, mais à Bafoussam. »\\
\esp{No vino a trabajar, sino que vino a descansar.} « Il n'est pas venu travailler, mais se reposer. »
\end{exemplebox}

\begin{attention}[\esp{no… sino} n'est jamais une simple restriction]
\esp{No… sino} \textbf{corrige} une première affirmation : il suppose deux éléments opposés (X nié, Y affirmé). On ne peut pas traduire « Je n'ai que deux frères » par \esp{no tengo dos hermanos sino…} : il n'y a rien à corriger. Test : si l'on peut dire « pas X mais Y », alors \esp{sino} convient ; sinon, utilisez \esp{solo} ou \esp{no… más que}.
\end{attention}

\courssub{Cas particulier : « ne… que » + verbe → \esp{no hacer más que} + infinitif}

Quand « ne… que » porte sur une \textbf{action répétée ou permanente} (« il ne fait que… »), l'espagnol possède une tournure idiomatique élégante : \cle{no hacer más que + infinitif}.

\begin{exemplebox}[Exemples traduits]
\esp{No hace más que hablar.} « Il ne fait que parler. »\\
\esp{No hace más que estudiar.} « Elle ne fait qu'étudier. »\\
\esp{No hacéis más que quejaros.} « Vous ne faites que vous plaindre. »\\
\esp{El niño no hace más que jugar con el móvil.} « L'enfant ne fait que jouer avec le portable. »
\end{exemplebox}

On peut aussi dire \esp{solo habla}, \esp{únicamente estudia}, mais la tournure avec \esp{hacer} rend fidèlement la nuance d'agacement ou d'habitude du français.



\courssub{Complément utile pour le Bac : « non seulement… mais aussi »}

\begin{propriete}[\esp{no solo… sino también}]
La tournure française « non seulement… mais aussi » se traduit par \cle{no solo… sino también} (ou \esp{no solamente… sino también}). Attention : on retrouve ici \esp{sino}, car la seconde partie \emph{ajoute} un élément qui dépasse la première affirmation.
\end{propriete}

\begin{exemplebox}[Exemples traduits]
\esp{No solo juega al fútbol, sino también al baloncesto.} « Il ne joue pas seulement au football, mais aussi au basket-ball. »\\
\esp{El español no solo se habla en España, sino también en Guinea Ecuatorial.} « L'espagnol ne se parle pas seulement en Espagne, mais aussi en Guinée équatoriale. »
\end{exemplebox}

\courssub{Tableaux récapitulatifs}

\begin{center}
\begin{tabularx}{\linewidth}{|X|X|X|}
\hline
\rowcolor{popblueL} Français & Español & Contexte d'emploi \\
\hline
Je n'ai que dix ans. & Solo tengo diez años. & Restriction simple, registre courant \\
\hline
Je n'ai que dix ans (à peine). & No tengo más que diez años. & Restriction insistante, soutenu \\
\hline
Il ne boit pas de vin, mais de l'eau. & No bebe vino, sino agua. & Opposition, correction \\
\hline
Elle ne fait qu'étudier. & No hace más que estudiar. & Action répétée, habitude \\
\hline
Non seulement il chante, mais il danse. & No solo canta, sino también baila. & Addition graduelle \\
\hline
\end{tabularx}
\end{center}

\begin{popfigure}
\begin{tikzpicture}[node distance=0.4cm, every node/.style={align=center}]
\node[draw=popblue, fill=popblueL, rounded corners, font=\bfseries] (fr) {Français : « ne… que »};
\node[draw=popgreen, fill=popgreenL, rounded corners, text width=4cm, below left=1.2cm and 1.5cm of fr] (solo) {\textbf{solo / únicamente}\\ restriction simple\\ \emph{Solo tengo dos hijos.}};
\node[draw=poporange, fill=poporangeL, rounded corners, text width=4cm, below=1.5cm of fr] (masque) {\textbf{no… más que}\\ restriction insistante\\ \emph{No hay más que Dios.}};
\node[draw=poppurple, fill=poppurpleL, rounded corners, text width=4cm, below right=1.2cm and 1.5cm of fr] (sino) {\textbf{no… sino}\\ opposition, correction\\ \emph{No es médico, sino enfermero.}};
\draw[-{Stealth}, thick, popgreen] (fr) -- (solo);
\draw[-{Stealth}, thick, poporange] (fr) -- (masque);
\draw[-{Stealth}, thick, poppurple] (fr) -- (sino);
\node[draw=popgold, fill=popgoldL, rounded corners, text width=4.5cm, below=0.8cm of masque] (hacer) {\textbf{no hacer más que + inf.}\\ action répétée\\ \emph{No hace más que hablar.}};
\draw[-{Stealth}, thick, popgold] (masque) -- (hacer);
\end{tikzpicture}
\legende{Carte de décision : les quatre traductions de « ne… que » selon le contexte.}
\end{popfigure}

\begin{exoresolu}[Traduire en choisissant la bonne structure]
Traduisez en espagnol : 1) « Nous n'avons qu'un professeur d'espagnol. » 2) « Ce n'est pas un examen, mais un contrôle. » 3) « Mon frère ne fait que regarder la télévision. » 4) « Il ne reste que deux jours avant les vacances. »
\tcblower
\textbf{Solution détaillée :}
\begin{enumerate}
\item Restriction simple : \esp{Solo tenemos un profesor de español.} (ou, avec insistance : \esp{No tenemos más que un profesor de español.})
\item Opposition « pas X mais Y » : \esp{No es un examen, sino un control.}
\item « Ne… que » + verbe d'action : \esp{Mi hermano no hace más que ver la televisión.}
\item Restriction simple : \esp{Solo quedan dos días antes de las vacaciones.} (ou \esp{No quedan más que dos días…})
\end{enumerate}
\textbf{Méthode :} on identifie d'abord le type de « ne… que » (restriction ou opposition), puis on place l'élément espagnol exactement devant l'élément restreint, sans jamais traduire « que » isolément.
\end{exoresolu}

\begin{exercice}[À vous de traduire]
Traduisez en espagnol : 1) « Je n'ai que trois sœurs. » 2) « Il n'y a que toi qui puisses m'aider. » (utilisez \esp{no… más que}) 3) « Elle n'écrit pas des lettres, mais des poèmes. » 4) « Les élèves ne font que rire. » 5) « Non seulement il parle français, mais aussi espagnol. »
\end{exercice}


\begin{corrige}[Exercice : traduction de « ne… que »]
1) \esp{Solo tengo tres hermanas.} ou \esp{No tengo más que tres hermanas.} 2) \esp{No hay más que tú que puedas ayudarme.} 3) \esp{No escribe cartas, sino poemas.} 4) \esp{Los alumnos no hacen más que reír.} 5) \esp{No solo habla francés, sino también español.}
\end{corrige}

\begin{savaistu}[\esp{solo} s'écrit sans accent depuis 2010]
Avant 2010, la Real Academia Española recommandait d'écrire \esp{sólo} avec un accent quand il signifiait « únicamente », pour le distinguer de l'adjectif \esp{solo} (« seul »). Depuis la réforme de 2010, l'accent a disparu dans tous les cas : on écrit toujours \esp{solo}, car le contexte suffit à lever toute ambiguïté. Ainsi, \esp{Solo tengo diez años} (« je n'ai que dix ans ») et \esp{Estoy solo} (« je suis seul ») s'écrivent désormais de la même façon. Au baccalauréat, n'ajoutez donc jamais d'accent sur ce mot.
\end{savaistu}

\begin{aretenir}[L'essentiel de la section]
\begin{itemize}
\item Restriction simple : \esp{solo} / \esp{solamente} / \esp{únicamente} placé juste avant l'élément restreint, sans négation.
\item Restriction insistante : \esp{no… más que}.
\item Opposition ou correction : \esp{no… sino} (\esp{sino que} devant un verbe conjugué).
\item « Ne faire que » + infinitif : \esp{no hacer más que} + infinitif.
\item « Non seulement… mais aussi » : \esp{no solo… sino también}.
\item Interdit absolu : traduire « que » tout seul (\esp{no tengo que dos hermanos} est un calque fautif).
\end{itemize}
\end{aretenir}

\courssub{Lexique : la famille et le portrait}

Pour décrire sa famille, il faut maîtriser le vocabulaire de la parenté, les adjectifs de description physique et morale, ainsi que les verbes de relation (\esp{ser}, \esp{tener}, \esp{llevarse}).

\begin{definition}[Vocabulaire de la parenté \textit{(léxico de la parentela)}]
\begin{center}
\begin{tabularx}{\linewidth}{|X|X|X|}
\hline
\rowcolor{popblueL} Français & Español & Note \\
\hline
Le père / La mère & El padre / La madre & \esp{Los padres} = les parents \\
\hline
Le grand-père / La grand-mère & El abuelo / La abuela & \esp{Los abuelos} = les grands-parents \\
\hline
Le frère / La sœur & El hermano / La hermana & \esp{Los hermanos} = les frères et sœurs \\
\hline
L'oncle / La tante & El tío / La tía & \esp{Los tíos} = les oncles et tantes \\
\hline
Le cousin / La cousine & El primo / La prima & \esp{Los primos} = les cousins \\
\hline
Le fils / La fille & El hijo / La hija & \esp{Los hijos} = les enfants \\
\hline
Le neveu / La nièce & El sobrino / La sobrina & \\
\hline
Le beau-père / La belle-mère & El suegro / La suegra & Parents du conjoint \\
\hline
Le gendre / La bru & El yerno / La nuera & Conjoints des enfants \\
\hline
Le beau-frère / La belle-sœur & El cuñado / La cuñada & Frères/sœurs du conjoint \\
\hline
Le demi-frère / La demi-sœur & El hermanastro / La hermanastra & Un parent commun \\
\hline
Le parrain / La marraine & El padrino / La madrina & Important en pays hispanophones \\
\hline
\end{tabularx}
\end{center}
\end{definition}

\begin{definition}[Portrait physique \textit{(retrato físico)}]
\begin{center}
\begin{tabularx}{\linewidth}{|X|X|}
\hline
\rowcolor{popgreenL} Español & Français \\
\hline
\esp{alto / baja}, \esp{grande / pequeño} & grand / petit \\
\esp{delgado / delgada}, \esp{flaco / flaca} & mince, maigre \\
\esp{gordo / gorda}, \esp{relleno / rellena} & gros, rond \\
\esp{joven}, \esp{mayor}, \esp{anciano / anciana} & jeune, âgé, vieux \\
\esp{guapo / guapa}, \esp{bonito / bonita}, \esp{feo / fea} & beau, joli, laid \\
\esp{ojos azules / marrones / verdes} & yeux bleus / marron / verts \\
\esp{pelo rubio / castaño / negro / pelirrojo} & cheveux blonds / châtains / noirs / roux \\
\esp{liso / rizado / ondulado} & lisse / bouclé / ondulé \\
\esp{llevar barba / bigote / gafas} & porter barbe / moustache / lunettes \\
\esp{parecerse a} & ressembler à \\
\hline
\end{tabularx}
\end{center}
\end{definition}

\begin{definition}[Portrait moral \textit{(retrato moral / carácter)}]
\begin{center}
\begin{tabularx}{\linewidth}{|X|X|X|}
\hline
\rowcolor{poppurpleL} Positivo & Negativo & Neutre / Nuancé \\
\hline
\esp{simpático / a} (aimable) & \esp{antipático / a} (désagréable) & \esp{tímido / a} (timide) \\
\esp{amable} (gentil) & \esp{maleducado / a} (mal élevé) & \esp{callado / a} (silencieux) \\
\esp{trabajador / a} (travailleur) & \esp{perezoso / a} (paresseux) & \esp{serio / a} (sérieux) \\
\esp{generoso / a} (généreux) & \esp{egoísta} (égoïste) & \esp{callado / a} (discret) \\
\esp{sincero / a} (sincère) & \esp{mentiroso / a} (menteur) & \esp{testarudo / a} (têtu) \\
\esp{alegre / divertido / a} (joyeux) & \esp{triste / aburrido / a} (triste/ennuyeux) & \esp{celoso / a} (jaloux) \\
\esp{responsable} (responsable) & \esp{irresponsable} & \esp{independiente} (indépendant) \\
\esp{inteligente / listo / a} (intelligent) & \esp{tonto / a} (bête) & \esp{soñador / a} (rêveur) \\
\hline
\end{tabularx}
\end{center}
\end{definition}

\begin{propriete}[Verbes clés et structures pour la description]
\begin{itemize}
\item \textbf{\esp{Ser}} : traits permanents, identité, physique, caractère (\esp{Es alto, es médico, es simpático}).
\item \textbf{\esp{Tener}} : âge, parties du corps, possession (\esp{Tiene 15 años, tiene ojos verdes, tiene coche}).
\item \textbf{\esp{Llevar}} : accessoires, coiffure, barbe (\esp{Lleva gafas, lleva el pelo corto, lleva barba}).
\item \textbf{\esp{Llevarse bien / mal (con)}} : s'entendre bien / mal (avec). \emph{Ex:} \esp{Me llevo muy bien con mi hermana.}
\item \textbf{\esp{Vivir con / en}} : habiter avec / à. \emph{Ex:} \esp{Vivo con mis padres en Yaoundé.}
\item \textbf{\esp{Dedicarse a / Trabajar de / Estudiar}} : profession / études. \emph{Ex:} \esp{Mi padre se dedica al comercio.}
\end{itemize}
\end{propriete}

\courssub{Rédaction guidée : \esp{describir a su familia} (80 palabras)}

\begin{methode}[Étapes pour réussir la rédaction (80 mots ± 10\%)]
\begin{enumerate}
\item \textbf{Analyser le sujet :} « Décrivez votre famille » $\rightarrow$ structure (combien, qui), portraits (physique + moral), relations / activités communes.
\item \textbf{Organiser en 3 paragraphes logiques :}
      \begin{itemize}
      \item \textit{Intro (15-20 mots) :} Composition du foyer, lieu de vie.
      \item \textit{Développement (40-50 mots) :} Portrait de 2-3 membres clés (père, mère, frère/sœur). Utiliser \esp{ser}, \esp{tener}, \esp{llevar}. Varier le vocabulaire moral/physique.
      \item \textit{Conclusion (15-20 mots) :} Relations familiales (\esp{llevarse bien}), activités partagées, sentiment global.
      \end{itemize}
\item \textbf{Rédiger en espagnol directement} (pas de traduction mot à mot). Surveiller les accords (genre/nombre) et \esp{ser}/\esp{estar}.
\item \textbf{Compter les mots} (les contractions \esp{del}, \esp{al} comptent pour 1 mot ; les noms propres pour 1).
\item \textbf{Relire :} Orthographe (accents, \esp{ñ}, \esp{ü}), ponctuation (\esp{¿ ? ¡ !}), cohérence des temps (présent indicatif généralement).
\end{enumerate}
\end{methode}

\begin{exemplebox}[Modèle de rédaction (environ 80 mots)]
\esp{Vivo en Douala con mis padres y mis dos hermanos menores. Mi padre, un hombre alto y fuerte de cuarenta y cinco años, es comerciante y muy trabajador; tiene el pelo negro y lleva bigote. Mi madre, baja y delgada, es ama de casa, alegre y generosa; siempre lleva el pelo recogido. Mi hermano Juan, de doce años, es travieso y deportista: no hace más que jugar al fútbol. Mi hermana Ana, de diez, es tímida y estudiosa, lee mucho. Nos llevamos todos muy bien y los domingos comemos juntos en la aldea de mis abuelos. Es una familia unida y feliz.}
\end{exemplebox}

\textbf{Analyse du modèle :}
\begin{itemize}
\item \textit{Structure :} Intro (famille, lieu) $\rightarrow$ Père (physique, métier, moral) $\rightarrow$ Mère (physique, rôle, moral) $\rightarrow$ Frère (âge, moral, action \esp{no hace más que}) $\rightarrow$ Sœur (âge, moral) $\rightarrow$ Conclusion (relations, tradition, sentiment).
\item \textit{Grammaire :} \esp{ser} (identité, traits), \esp{tener} (âge, cheveux), \esp{llevar} (bigote, pelo), \esp{no hace más que} (restriction action), \esp{llevarse bien} (relation).
\item \textit{Compte :} ~82 mots. Parfait.
\end{itemize}



\begin{exercice}[À vous d'écrire]
Rédigez en espagnol une description de votre famille (ou d'une famille imaginaire) en \textbf{80 mots environ}. Respectez la structure en 3 parties et utilisez au moins : \esp{ser}, \esp{tener}, \esp{llevar}, \esp{llevarse bien}, un adjectif physique, un adjectif moral, et une tournure de restriction (\esp{solo} ou \esp{no… más que}).
\end{exercice}


\begin{corrige}[Exercice : Redacción - Describir a su familia]
\textbf{Proposition de correction (exemple élève) :}
\esp{Mi familia es grande y unida. Vivimos en Yaoundé: mis padres, mis tres hermanos y yo. Mi padre es ingeniero, alto y canoso, muy serio pero justo. Mi madre es profesora, baja y morena, siempre alegre y paciente; lleva gafas. Mi hermano mayor, Carlos, tiene veinte años, es guapo y deportista, solo piensa en el baloncesto. Mis hermanas gemelas, Ana y Bea, son quinceañeras, simpáticas y estudiosas; no hacen más que estudiar. Yo soy el menor, responsable y callado. Nos queremos mucho y los fines de semana visitamos a los abuelos en el pueblo.}

\textbf{Points de vigilance pour l'auto-évaluation :}
\begin{itemize}
\item \checkmark 80 mots $\pm$ 10.
\item \checkmark Vocabulaire famille, physique, moral présent.
\item \checkmark \esp{Ser} / \esp{Tener} / \esp{Llevar} bien employés.
\item \checkmark Tournure de restriction utilisée (\esp{solo piensa}, \esp{no hacen más que}).
\item \checkmark Accords (genre/nombre) : \esp{ingeniero}, \esp{canoso}, \esp{serio}, \esp{justo} (masc. sg.) ; \esp{profesora}, \esp{baja}, \esp{morena}, \esp{alegre}, \esp{paciente} (fém. sg.) ; \esp{gemelas}, \esp{quinceañeras}, \esp{simpáticas}, \esp{estudiosas} (fém. pl.).
\item \checkmark Ponctuation espagnole correcte.
\end{itemize}
\end{corrige}


\coursec{Lexique : la famille et le portrait}

Dans la section précédente, nous avons appris à traduire la restriction \esp{ne… que} par \esp{solo}, \esp{no… más que} ou \esp{no… sino}. Nous savons donc maintenant dire : \esp{No tengo más que dos hermanos.} « Je n'ai que deux frères et sœurs. » Mais pour rédiger notre description de famille en 80 mots — l'objectif final de cette leçon — il ne suffit pas de maîtriser les structures grammaticales : il faut posséder un vocabulaire riche et précis. C'est l'objet de cette section : le lexique de la parenté, de la situation familiale, du portrait physique et du portrait moral, ainsi que les verbes et connecteurs indispensables à la rédaction.

\courssub{Observer d'abord : un texte de présentation familiale}

Lisons d'abord ce court texte de présentation, rédigé par une élève de Terminale de Yaoundé. Soulignez mentalement tous les mots qui désignent des membres de la famille.

\begin{textebox}[Mi familia, por Aminatou]
Me llamo Aminatou y vivo en Yaoundé con mi familia. Somos una familia numerosa y muy unida. Mis padres se llaman Ibrahim y Aïcha. Mi padre es alto y delgado, de ojos negros; es estricto pero muy cariñoso con nosotros. Mi madre es baja y alegre; es trabajadora y paciente, y siempre nos apoya en los estudios.

Tengo dos hermanos y una hermana. Mi hermano mayor se parece mucho a mi padre: es moreno y de pelo rizado. Mi hermana pequeña es muy graciosa y se lleva bien con todo el mundo. Los fines de semana, visitamos a mis abuelos, que viven en un pueblo cerca de Douala. Mi abuelo es viudo: mi abuela murió hace tres años. Allí también vemos a mis tíos, a mis primos y a mis sobrinos. Mi tío Ahmed está casado con mi tía Mariam, y tienen tres hijos. Me llevo muy bien con mis primos: somos como hermanos.

Mi familia no es solo grande: es también generosa. Cuando alguien tiene un problema, todos le ayudan. Por eso me siento orgullosa de mi familia.
\end{textebox}

\textbf{Glossaire.} \esp{numerosa} : nombreuse ; \esp{unida} : unie, soudée ; \esp{estricto} : strict ; \esp{cariñoso} : affectueux ; \esp{apoyar} : soutenir ; \esp{se parece a} : ressemble à ; \esp{moreno} : brun, au teint mat ; \esp{de pelo rizado} : aux cheveux bouclés ; \esp{graciosa} : drôle, charmante ; \esp{se lleva bien con} : s'entend bien avec ; \esp{viudo} : veuf ; \esp{murió} : est morte ; \esp{orgullosa} : fière.

\textbf{Questions de compréhension.}
\begin{enumerate}
\item ¿Cuántos hermanos tiene Aminatou ?
\item ¿Cómo es físicamente su padre ? ¿Y moralmente ?
\item ¿Dónde viven los abuelos de Aminatou ?
\item ¿Con quién se lleva muy bien Aminatou ?
\end{enumerate}

\courssub{La parenté : les mots de la famille}

\begin{definition}[La parenté]
La \cle{parenté} (\esp{el parentesco}) désigne les liens familiaux entre les personnes. En espagnol, le vocabulaire de la famille est très régulier : presque tous les noms masculins en \esp{-o} ont un féminin en \esp{-a} : \esp{el tío / la tía}, \esp{el primo / la prima}, \esp{el sobrino / la sobrina}.
\end{definition}

\begin{textebox}[Vocabulaire de la parenté]
\begin{center}
\begin{tabularx}{\linewidth}{|X|X|}
\hline
\rowcolor{popblueL} \textbf{Español} & \textbf{Français} \\
\hline
el abuelo / la abuela & le grand-père / la grand-mère \\
\hline
el abuelo paterno / materno & le grand-père paternel / maternel \\
\hline
el padre / la madre & le père / la mère \\
\hline
los padres & les parents (père et mère) \\
\hline
el esposo / la esposa & le mari / la femme (épouse) \\
\hline
el hermano / la hermana & le frère / la sœur \\
\hline
el hermano mayor / menor & le frère aîné / cadet \\
\hline
el tío / la tía & l'oncle / la tante \\
\hline
el primo / la prima & le cousin / la cousine \\
\hline
el sobrino / la sobrina & le neveu / la nièce \\
\hline
el cuñado / la cuñada & le beau-frère / la belle-sœur \\
\hline
el suegro / la suegra & le beau-père / la belle-mère (parents du conjoint) \\
\hline
el padrastro / la madrastra & le beau-père / la belle-mère (remariage) \\
\hline
el hijo / la hija & le fils / la fille \\
\hline
hijo único / hija única & fils unique / fille unique \\
\hline
el nieto / la nieta & le petit-fils / la petite-fille \\
\hline
los parientes & les parents (famille élargie) \\
\hline
\end{tabularx}
\end{center}
\end{textebox}

\begin{attention}[Los padres, los parientes : attention aux faux amis !]
Le mot \esp{los padres} ne signifie pas « les parents » au sens large : il désigne uniquement \cle{le père et la mère}. Pour parler de la famille élargie (oncles, tantes, cousins…), on dit \esp{los parientes} ou \esp{los familiares}. De même, \esp{el padrastro} et \esp{el suegro} se traduisent tous les deux par « beau-père », mais \esp{el padrastro} est le nouveau mari de la mère, tandis que \esp{el suegro} est le père du conjoint. Enfin, ne confondez pas \esp{el sobrino} (le neveu, fils de votre frère ou sœur) et \esp{el nieto} (le petit-fils, fils de votre enfant) : ces deux mots sont souvent intervertis par les francophones.
\end{attention}

\begin{propriete}[L'accord des adjectifs et le masculin pluriel]
En espagnol comme en français, un groupe mixte se met au masculin pluriel : \esp{mis hermanos} peut signifier « mes frères » ou « mes frères et sœurs » ; \esp{mis tíos} = « mes oncles et tantes ». L'adjectif s'accorde en genre et en nombre avec le nom : \esp{mi hermana mayor} « ma sœur aînée », \esp{mis primos pequeños} « mes petits cousins ».
\end{propriete}

\begin{popfigure}
\begin{center}
\begin{tikzpicture}[
  every node/.style={draw, rounded corners=2pt, align=center, font=\small, inner sep=4pt},
  gen1/.style={fill=poppurpleL},
  gen2/.style={fill=popblueL},
  gen3/.style={fill=popgreenL},
  link/.style={draw=popmuted, thick}
]
\node[gen1, align=center] (abu) at (0,3.4){los abuelos\\(grands-parents)};
\node[gen2, align=center] (padres) at (-2.6,1.7){los padres\\(parents)};
\node[gen2, align=center] (tios) at (2.6,1.7){los tíos\\(oncles et tantes)};
\node[gen3, align=center] (yo) at (-3.6,0){\textbf{yo}\\(moi)};
\node[gen3, align=center] (herm) at (-1.4,0){mis hermanos\\(frères et sœurs)};
\node[gen3, align=center] (primos) at (2.6,0){mis primos\\(cousins)};
\draw[link] (abu) -- (padres);
\draw[link] (abu) -- (tios);
\draw[link] (padres) -- (yo);
\draw[link] (padres) -- (herm);
\draw[link] (tios) -- (primos);
\end{tikzpicture}
\end{center}
\legende{Arbre généalogique simple : trois générations de la famille.}
\end{popfigure}

\courssub{La situation familiale}

\begin{definition}[El estado civil]
La \cle{situation familiale} (\esp{el estado civil}) indique la situation matrimoniale d'une personne. Voici les termes essentiels :
\end{definition}

\begin{center}
\begin{tabularx}{\linewidth}{|X|X|X|}
\hline
\rowcolor{popblueL} \textbf{Español} & \textbf{Français} & \textbf{Exemple} \\
\hline
soltero / soltera & célibataire & \esp{Mi tío es soltero.} \\
\hline
casado / casada & marié(e) & \esp{Está casada con un médico.} \\
\hline
divorciado / divorciada & divorcé(e) & \esp{Mis tíos están divorciados.} \\
\hline
viudo / viuda & veuf / veuve & \esp{Mi abuelo es viudo.} \\
\hline
familia numerosa & famille nombreuse & \esp{Somos una familia numerosa.} \\
\hline
familia unida & famille unie & \esp{Es una familia muy unida.} \\
\hline
\end{tabularx}
\end{center}

\begin{attention}[Ser ou estar avec la situation familiale ?]
On dit \esp{ser soltero} (état durable) mais \esp{estar casado}, \esp{estar divorciado}, \esp{estar separado} (résultat d'un changement d'état). Comparez : \esp{Mi hermana es soltera.} « Ma sœur est célibataire. » / \esp{Mi hermana está casada desde 2020.} « Ma sœur est mariée depuis 2020. »
\end{attention}

\courssub{Le portrait physique}

Pour décrire l'apparence d'une personne, l'espagnol utilise le verbe \esp{ser} (caractéristiques permanentes) et la préposition \esp{de} pour les traits précis.

\begin{propriete}[Décrire le physique]
\begin{itemize}
\item \textbf{ser + adjectif} : \esp{Mi padre es alto y delgado.} « Mon père est grand et mince. »
\item \textbf{de + nom} pour les yeux et les cheveux : \esp{Es de ojos negros y de pelo rizado.} « Il a les yeux noirs et les cheveux bouclés. »
\item \textbf{tener + article + nom} est aussi possible : \esp{Tiene los ojos negros.} « Elle a les yeux noirs. »
\end{itemize}
\end{propriete}

\begin{center}
\begin{tabularx}{\linewidth}{|X|X|}
\hline
\rowcolor{popblueL} \textbf{Español} & \textbf{Français} \\
\hline
alto / alta & grand(e) \\
\hline
bajo / baja & petit(e) \\
\hline
delgado / delgada & mince \\
\hline
gordo / gorda & gros(se) \\
\hline
moreno / morena & brun(e), au teint mat \\
\hline
de ojos negros / azules & aux yeux noirs / bleus \\
\hline
de pelo rizado / liso / corto & aux cheveux bouclés / lisses / courts \\
\hline
joven / mayor & jeune / âgé(e), aîné(e) \\
\hline
guapo / guapa & beau / belle \\
\hline
\end{tabularx}
\end{center}

\begin{attention}[Faux amis et calques]
\begin{itemize}
\item \esp{grande} signifie « grand » en taille ou en importance, mais pour la taille d'une personne on préfère \esp{alto} : \esp{Mi hermano es alto.} et non \esp{es grande}.
\item \esp{bajo} = « petit » (de taille), jamais « bas » au sens figuré.
\item Ne dites pas \esp{tiene ojos negros} sans article : la forme correcte est \esp{tiene los ojos negros} ou \esp{es de ojos negros}.
\end{itemize}
\end{attention}

\courssub{Le portrait moral et les verbes des relations}

\begin{textebox}[Le caractère et les relations familiales]
\begin{center}
\begin{tabularx}{\linewidth}{|X|X|}
\hline
\rowcolor{popblueL} \textbf{Español} & \textbf{Français} \\
\hline
cariñoso / cariñosa & affectueux / affectueuse \\
\hline
estricto / estricta & strict(e) \\
\hline
generoso / generosa & généreux / généreuse \\
\hline
trabajador / trabajadora & travailleur / travailleuse \\
\hline
paciente / impaciente & patient(e) / impatient(e) \\
\hline
alegre & joyeux / joyeuse \\
\hline
llevarse bien con & bien s'entendre avec \\
\hline
llevarse mal con & mal s'entendre avec \\
\hline
parecerse a & ressembler à \\
\hline
querer a & aimer (une personne) \\
\hline
apoyar & soutenir \\
\hline
criar & élever (des enfants) \\
\hline
cuidar (de) & s'occuper (de), prendre soin (de) \\
\hline
\end{tabularx}
\end{center}
\end{textebox}

\begin{propriete}[Construction des verbes de relation]
\begin{itemize}
\item \esp{parecerse a + personne} : verbe pronominal, toujours avec la préposition \esp{a} : \esp{Me parezco a mi padre.} « Je ressemble à mon père. »
\item \esp{llevarse bien/mal con + personne} : \esp{Nos llevamos muy bien.} « Nous nous entendons très bien. »
\item \esp{querer a + personne} : le \esp{a} personnel est obligatoire devant la personne aimée : \esp{Quiero mucho a mi abuela.} « J'aime beaucoup ma grand-mère. »
\end{itemize}
\end{propriete}

\begin{exemplebox}[Exemples traduits]
\esp{Mi hermana mayor se parece a mi madre.} « Ma sœur aînée ressemble à ma mère. »

\esp{Me llevo muy bien con mis primos.} « Je m'entends très bien avec mes cousins. »

\esp{Mis padres nos criaron con mucho cariño.} « Mes parents nous ont élevés avec beaucoup d'affection. »

\esp{Mi abuela cuida de nosotros cuando nuestros padres trabajan.} « Ma grand-mère s'occupe de nous quand nos parents travaillent. »
\end{exemplebox}

\courssub{Les connecteurs pour la rédaction}

Une description de 80 mots doit être fluide : les connecteurs relient les phrases entre elles et évitent la juxtaposition de phrases courtes.

\begin{center}
\begin{tabularx}{\linewidth}{|X|X|X|}
\hline
\rowcolor{popblueL} \textbf{Conector} & \textbf{Français} & \textbf{Emploi} \\
\hline
además & de plus & ajouter une idée \\
\hline
también & aussi & ajouter une idée \\
\hline
sin embargo & cependant & opposition \\
\hline
en cambio & en revanche & contraste entre deux personnes \\
\hline
por eso & c'est pourquoi & conséquence \\
\hline
pero & mais & opposition simple \\
\hline
\end{tabularx}
\end{center}

\begin{exemplebox}[Connecteurs en contexte]
\esp{Mi padre es estricto; sin embargo, es muy cariñoso.} « Mon père est strict ; cependant, il est très affectueux. »

\esp{Mi hermano es tranquilo; en cambio, mi hermana es muy alegre.} « Mon frère est calme ; en revanche, ma sœur est très joyeuse. »

\esp{Mi familia siempre me apoya; por eso me siento orgulloso de ella.} « Ma famille me soutient toujours ; c'est pourquoi j'en suis fier. »
\end{exemplebox}

\begin{exoresolu}[Compléter un portrait familial]
Énoncé. Complétez ce texte avec les mots suivants : \esp{parece}, \esp{llevo}, \esp{unida}, \esp{abuelos}, \esp{además}.

« Somos una familia muy \ldots\ldots\ldots. Mi hermana se \ldots\ldots\ldots\ a mi madre: es alta y de ojos negros. Me \ldots\ldots\ldots\ muy bien con mis primos. \ldots\ldots\ldots, los domingos visitamos a mis \ldots\ldots\ldots\ en su pueblo. »

\tcblower
Solution détaillée.
\begin{enumerate}
\item \esp{unida} : on décrit le caractère de la famille (« une famille très unie ») ; accord féminin singulier avec \esp{familia}.
\item \esp{parece} : structure \esp{parecerse a} à la 3\textsuperscript{e} personne du singulier, avec le pronom réfléchi \esp{se} déjà placé devant : « ma sœur ressemble à ma mère ».
\item \esp{llevo} : \esp{llevarse bien con} à la 1\textsuperscript{re} personne du singulier : « je m'entends très bien avec mes cousins ».
\item \esp{Además} : connecteur d'addition, en tête de phrase, avec majuscule et accent sur le \esp{e} final.
\item \esp{abuelos} : complément du verbe \esp{visitamos} avec le \esp{a} personnel déjà présent dans la phrase : « le dimanche, nous rendons visite à nos grands-parents dans leur village ».
\end{enumerate}
Texte complet : « Somos una familia muy \textbf{unida}. Mi hermana se \textbf{parece} a mi madre: es alta y de ojos negros. Me \textbf{llevo} muy bien con mis primos. \textbf{Además}, los domingos visitamos a mis \textbf{abuelos} en su pueblo. »
\end{exoresolu}

\begin{exercice}[À vous de jouer]
\begin{enumerate}
\item Traduisez : « Mon oncle est divorcé ; cependant, il s'entend bien avec ma tante. »
\item Traduisez : « Mes grands-parents nous ont élevés avec beaucoup d'affection. »
\item Complétez : \esp{Mi primo es bajo; \ldots\ldots\ldots, mi prima es muy alta.}
\item Trouvez l'erreur : \esp{Mis padres y mis parientes son mi padre y mi madre.}
\end{enumerate}
\end{exercice}

\begin{corrige}[Exercice « À vous de jouer »]
\begin{enumerate}
\item \esp{Mi tío está divorciado; sin embargo, se lleva bien con mi tía.} (Attention : \esp{estar divorciado}, et \esp{llevarse} est pronominal : le pronom \esp{se} est obligatoire.)
\item \esp{Mis abuelos nos criaron con mucho cariño.} (\esp{criar} au prétérit, 3\textsuperscript{e} personne du pluriel : \esp{criaron}.)
\item \esp{en cambio} (contraste entre deux personnes différentes ; \esp{sin embargo} conviendrait pour une opposition dans le portrait d'une même personne).
\item Erreur de faux ami : \esp{los padres} désigne le père et la mère ; \esp{los parientes} désigne la famille élargie. Phrase correcte : \esp{Mis padres son mi padre y mi madre; mis parientes son toda mi familia.}
\end{enumerate}
\end{corrige}

\begin{aretenir}[L'essentiel du lexique familial]
\begin{itemize}
\item \esp{los padres} = le père et la mère ; \esp{los parientes} = la famille élargie.
\item Physique : \esp{ser alto/bajo/delgado}, \esp{de ojos negros}, \esp{de pelo rizado}, \esp{tener los ojos negros}.
\item Moral : \esp{cariñoso, estricto, generoso, trabajador, paciente, alegre}.
\item Relations : \esp{parecerse a}, \esp{llevarse bien/mal con}, \esp{querer a}, \esp{apoyar}, \esp{criar}, \esp{cuidar de}.
\item Situation : \esp{ser soltero} mais \esp{estar casado/divorciado} ; \esp{ser viudo}.
\item Connecteurs : \esp{además} (addition), \esp{sin embargo} et \esp{en cambio} (opposition), \esp{por eso} (conséquence).
\end{itemize}
\end{aretenir}

Vous disposez désormais de tout le vocabulaire nécessaire : dans la prochaine section, nous mettrons ces mots en pratique pour rédiger, en 80 mots, la description de votre propre famille.

\coursec{Rédaction guidée : describir a su familia (80 palabras)}

Après avoir maîtrisé le lexique de la parenté (\emph{padre, madre, hermanos, abuelos, tíos, primos}) et les adjectifs de description physique (\emph{alto, bajo, delgado, moreno, rubio, ojos verdes}) et moral (\emph{simpático, trabajador, estricto, cariñoso, generoso}), ainsi que les tournures d'emphase (\esp{Es... quien/el que...}, \esp{Fue... cuando...}) et de restriction (\esp{solo}, \esp{no... más que}, \esp{no... sino}), nous abordons maintenant l'épreuve de rédaction courte type Baccalauréat. La consigne est classique : \esp{« Describa a su familia (80 palabras) »}. Il ne s'agit pas de lister des noms, mais de produire un texte cohérent, structuré, grammaticalement correct et respectant la contrainte de longueur. Cette section vous guide pas à pas, de la planification à la relecture finale, en passant par l'analyse d'un modèle commenté.

\courssub{Analyse de la consigne et stratégie globale}

\begin{definition}[Consigne type Bac : « Describa a su familia (80 palabras) »]
La consigne impose trois contraintes majeures :
\begin{enumerate}
    \item \textbf{Le sujet :} La famille (la vôtre ou une famille fictive crédible). Il faut en donner la composition, décrire les membres principaux et évoquer les relations.
    \item \textbf{Le format :} Une rédaction courte, paragraphe unique ou deux courts paragraphes, sans titres ni puces.
    \item \textbf{La longueur :} \cle{80 palabras} (mots). Au Cameroun, comme dans la plupart des académies francophones, la tolérance est généralement de \textbf{± 10\%}, soit \textbf{72 à 88 mots}. En dessous de 60 mots ou au-delà de 100 mots, la note de « correction de la langue » et de « respect de la consigne » chute drastiquement.
\end{enumerate}
\end{definition}

\begin{propriete}[Stratégie du « Plan en 3 temps » pour 80 mots]
Pour tenir la longueur sans être trop bref ni hors sujet, structurez votre texte en trois blocs logiques d'environ 25--30 mots chacun :
\begin{center}
\begin{tabularx}{\linewidth}{|X|X|X|}
\hline
\rowcolor{popblueL}
\textbf{Bloc 1 : Présentation générale} & \textbf{Bloc 2 : Portraits (Physique \& Moral)} & \textbf{Bloc 3 : Relations \& Conclusion} \\
\hline
Composition du foyer (\emph{Somos cinco...}).\newline
Emphase d'introduction (\emph{Es mi madre quien...}). & Description de 2--3 membres clés.\newline
\cle{Ser} (identité, caractère) vs \cle{Estar} (état, aspect momentané).\newline
Restriction (\emph{Solo tengo un hermano...}). & Relations (\emph{Me llevo bien con...}).\newline
Sentiment final avec emphase (\emph{Es lo que más quiero}). \\
\hline
\end{tabularx}
\end{center}
\end{propriete}

\courssub{Étape 1 : Présentation générale — Composición del núcleo familiar}

Commencez par situer le foyer. Évitez la liste sèche (\emph{Mi padre se llama... Mi madre se llama...}). Préférez une phrase synthétique intégrant le nombre total et la structure.

\begin{exemplebox}[Ouvertures types]
\begin{itemize}
    \item \esp{Mi familia es numerosa y muy unida. Somos seis: mis padres, mis tres hermanos y yo.} (16 mots)
    \item \esp{Vivo en una familia monoparental con mi madre y mi hermana menor.} (13 mots)
    \item \esp{Mi familia no es muy grande, pero es el pilar de mi vida. Somos cuatro.} (15 mots)
\end{itemize}
\end{exemplebox}

\begin{aretenir}[Vocabulaire clé de la composition]
\begin{center}
\begin{tabular}{|l|l|l|}
\hline
\rowcolor{poporangeL}
\textbf{Français} & \textbf{Español} & \textbf{Note} \\
\hline
Famille nombreuse & \esp{familia numerosa} & \emph{Numerosa} s'accorde avec \emph{familia} (fém.) \\
Famille unie & \esp{familia unida} & Participe passé comme adjectif \\
Famille monoparentale & \esp{familia monoparental} & Terme technique, utile au Bac \\
Parents & \esp{padres} & Masculin pluriel = père + mère \\
Frères et sœurs & \esp{hermanos} & Masculin pluriel = frères + sœurs \\
Enfant unique & \esp{hijo único / hija única} & Accord au genre \\
Aîné / Cadet & \esp{hermano mayor / hermano menor} & \emph{Mayor/Menor} invariables en genre (sg.) \\
\hline
\end{tabular}
\end{center}
\end{aretenir}

\courssub{Étape 2 : Portraits physique et moral — Retratos físico y moral}

C'est le cœur de la description. Choisissez \textbf{deux ou trois membres} maximum (père, mère, un frère/sœur) pour ne pas dépasser 80 mots. Pour chaque membre, donnez \textbf{un trait physique} et \textbf{un trait moral}. C'est ici qu'il faut mobiliser \cle{ser} vs \cle{estar} et l'\cle{emphase}.

\begin{propriete}[Ser vs Estar dans le portrait : règle d'or]
\begin{center}
\begin{tabularx}{\linewidth}{|X|X|}
\hline
\rowcolor{popgreenL}
\textbf{\cle{Ser} + Adjectif (Identité, nature, caractère permanent)} & \textbf{\cle{Estar} + Adjectif (État, aspect momentané, résultat)} \\
\hline
\esp{Mi padre \textbf{es} alto y moreno.} (Traits innés) & \esp{Mi padre \textbf{está} cansado hoy.} (État passager) \\
\esp{Mi madre \textbf{es} cariñosa y trabajadora.} (Caractère) & \esp{Mi madre \textbf{está} enferma.} (Maladie temporaire) \\
\esp{Mi hermano \textbf{es} simpático.} (Personnalité) & \esp{Mi hermano \textbf{está} contento.} (Humeur du moment) \\
\esp{Ella \textbf{es} guapa.} (Beauté comme trait stable) & \esp{Ella \textbf{está} guapa con ese vestido.} (Elle a l'air belle \emph{avec cette robe}) \\
\hline
\end{tabularx}
\end{center}
\textbf{Attention aux adjectifs changeants de sens :} \esp{ser bueno} (gentil) vs \esp{estar bueno} (appétissant / sexy) ; \esp{ser malo} (méchant) vs \esp{estar malo} (malade / mauvais au goût) ; \esp{ser rico} (riche) vs \esp{estar rico} (délicieux).
\end{propriete}

\begin{methode}[Insérer l'emphase et la restriction (Obligatoire au Bac)]
Pour valider les acquis du module, votre rédaction \textbf{doit} contenir au moins :
\begin{enumerate}
    \item \textbf{Une tournure d'emphase} (valeur 1 à 2 points au barème « structures complexes ») :
    \begin{itemize}
        \item \esp{Es mi madre \textbf{quien} cuida de todos.} (C'est ma mère \textbf{qui} s'occupe de tous)
        \item \esp{Fue ayer \textbf{cuando} nos reunimos todos.} (C'est hier \textbf{que} nous nous sommes réunis)
        \item \esp{Lo que más quiero \textbf{es} mi familia.} (Ce que j'aime le plus \textbf{c'est} ma famille)
    \end{itemize}
    \item \textbf{Une tournure de restriction} (\emph{ne... que}) :
    \begin{itemize}
        \item \esp{\textbf{Solo} tengo una hermana.} (Je n'ai \textbf{qu'une} sœur)
        \item \esp{No tengo \textbf{más que} gratitud.} (Je n'ai \textbf{que} de la gratitude)
        \item \esp{No es \textbf{sino} mi mejor amigo.} (Il n'est \textbf{que} mon meilleur ami / c'est rien d'autre que)
    \end{itemize}
\end{enumerate}
\end{methode}

\courssub{Étape 3 : Relations et conclusion — Relaciones y conclusión}

Terminez par le lien affectif et une phrase de clôture forte, si possible emphatique.

\begin{exemplebox}[Connecteurs et formules de conclusion]
\begin{itemize}
    \item \esp{Me llevo \textbf{bien / mal} con mis hermanos.} (Je m'entends bien/mal avec...)
    \item \esp{Aunque a veces discutimos, nos queremos mucho.} (Bien qu'on se dispute parfois...)
    \item \esp{Mi familia \textbf{es lo más importante} para mí.} (Ma famille est ce qui compte le plus)
    \item \esp{\textbf{Es} mi familia \textbf{quien} me da fuerzas.} (C'est ma famille qui me donne des forces - Emphase finale)
    \item \esp{No cambiaría mi familia \textbf{por nada del mundo}.} (Je n'échangerais ma famille pour rien au monde)
\end{itemize}
\end{exemplebox}

\begin{popfigure}
\begin{tikzpicture}[
    node distance=1.2cm and 2cm,
    every node/.style={font=\small},
    pasoS/.style={rectangle, rounded corners, draw=popblue, thick, fill=popblueL, minimum width=3.5cm, minimum height=1.8cm, align=center},
    arrow/.style={-Stealth, thick, popblue}
]
\node[pasoS, align=center] (e1){\textbf{PASO 1}\\\textsc{Presentación}\\\esp{Somos 5. Es mi madre quien...}};
\node[pasoS, right=of e1, align=center] (e2){\textbf{PASO 2}\\\textsc{Retratos}\\\esp{Padre: estricto, trabajador.}\\\esp{Hermana: simpática, alta. Solo tengo...}};
\node[pasoS, right=of e2, align=center] (e3){\textbf{PASO 3}\\\textsc{Relaciones}\\\esp{Nos llevamos bien. Es lo que más quiero.}};
\draw[arrow] (e1) -- (e2);
\draw[arrow] (e2) -- (e3);
\node[below=0.3cm of e1, text=popdark, font=\tiny] {~25 palabras};
\node[below=0.3cm of e2, text=popdark, font=\tiny] {~30 palabras};
\node[below=0.3cm of e3, text=popdark, font=\tiny] {~25 palabras};
\end{tikzpicture}
\legende{Schéma de la structure en 3 étapes pour une rédaction de 80 mots équilibrés.}
\end{popfigure}

\courssub{Comptage des mots : stratégies pour viser 75--85 mots}

\begin{definition}[Qu'est-ce qu'un « palabra » en espagnol ?]
Au Bac, on compte comme un mot tout élément séparé par un espace.
\begin{itemize}
    \item Les contractions \textbf{comptent pour 1 mot} : \esp{del} (de + el), \esp{al} (a + el).
    \item Les clitiques pronominaux \textbf{comptent pour 1 mot} s'ils sont accolés (infinitif, gérondif, affirmatif) : \esp{cuida\textbf{rme}} (1), \esp{d\textbf{ámelo}} (1). Mais \textbf{2 mots} s'ils sont proclitiques (conjugué) : \esp{me \textbf{cuida}} (2), \esp{te \textbf{quiero}} (2).
    \item La ponctuation ne compte pas.
    \item Les chiffres écrits en lettres comptent (\esp{cinco} = 1) ; en chiffres (\esp{5}) comptent souvent pour 1 mais déconseillés.
\end{itemize}
\end{definition}

\begin{methode}[Technique du « Brouillon compté »]
\begin{enumerate}
    \item Rédigez votre phrase sur brouillon \textbf{en espagnol directement} (pas de traduction mot à mot depuis le français).
    \item Comptez les mots \textbf{à la fin de chaque phrase} et notez le cumul en marge.
    \item Si vous êtes à 65 mots après le bloc 2, vous savez qu'il vous reste exactement 15 mots pour le bloc 3. Ajustez le vocabulaire (ajoutez un adjectif, une subordonnée, ou coupez un adjectif).
    \item \textbf{Astuce d'ajustement :} Les adverbes en \emph{-mente} (\esp{rápidamente}, \esp{cariñosamente}) comptent pour 1 mot et ajoutent de la nuance. Les locutions (\esp{por eso}, \esp{por lo tanto}) comptent pour 2--3 mots.
\end{enumerate}
\end{methode}

\courssub{Modèle de rédaction commenté (80 palabras exactos)}

Voici un modèle type « copie idéale » respectant toutes les contraintes : longueur, structures ciblées (emphase, restriction), lexique varié, grammaire parfaite (\emph{ser/estar}, accords).

\begin{textebox}[Modèle : Descripción de mi familia (80 palabras)]
Mi familia es numerosa y muy unida. Somos seis: mis padres, tres hermanos y yo. \textbf{Es mi madre quien} cuida de todos nosotros; es cariñosa y trabajadora. Mi padre, en cambio, es estricto pero generoso. \textbf{Solo tengo una hermana}: se llama Amina y se parece mucho a mi abuela. Me llevo muy bien con mis hermanos, aunque a veces discutimos. No tengo \textbf{más que} gratitud por mi familia querida, porque \textbf{realmente es lo que más quiero} en la vida entera.
\end{textebox}

\begin{center}
\textbf{Traduction et Glossaire}
\end{center}
\begin{itemize}
    \item \esp{Mi familia es numerosa y muy unida.} = Ma famille est nombreuse et très unie.
    \item \esp{Somos seis: mis padres, tres hermanos y yo.} = Nous sommes six : mes parents, trois frères/sœurs et moi. (\emph{hermanos} = frères et sœurs).
    \item \esp{\textbf{Es mi madre quien} cuida de todos nosotros;} = \textbf{C'est ma mère qui} s'occupe de nous tous ; (\textbf{Emphase : C'est... qui} $\rightarrow$ \esp{Es... quien}).
    \item \esp{es cariñosa y trabajadora.} = elle est affectueuse et travailleuse. (\textbf{Accords fém. sing.} : \emph{cariños\textbf{a}}, \emph{trabajador\textbf{a}}).
    \item \esp{Mi padre, en cambio, es estricto pero generoso.} = Mon père, en revanche, est strict mais généreux. (\textbf{Accords masc. sing.} ; \emph{en cambio} = au contraire / en revanche).
    \item \esp{\textbf{Solo tengo una hermana}:} = \textbf{Je n'ai qu'une sœur} : (\textbf{Restriction : Ne... que} $\rightarrow$ \esp{Solo} + indicatif).
    \item \esp{se llama Amina y se parece mucho a mi abuela.} = elle s'appelle Amina et elle ressemble beaucoup à ma grand-mère. (\emph{se parece a} = pronominal, \emph{mucho} = beaucoup).
    \item \esp{Me llevo muy bien con mis hermanos,} = Je m'entends très bien avec mes frères/sœurs. (\emph{llevarse bien con} = s'entendre bien avec).
    \item \esp{aunque a veces discutimos.} = bien qu'on se dispute parfois. (\emph{although} + indicatif = fait réel ; \emph{discutir} = se disputer/bavarder).
    \item \esp{No tengo \textbf{más que} gratitud por mi familia querida,} = Je n'ai \textbf{que} de la gratitude pour ma famille aimée. (\textbf{Restriction : No... más que} = Ne... que).
    \item \esp{porque \textbf{realmente es lo que más quiero} en la vida entera.} = parce que \textbf{c'est vraiment ce que j'aime le plus} au monde entier. (\textbf{Emphase : C'est... que} $\rightarrow$ \esp{Es lo que...} ; \emph{realmente} = vraiment ; \emph{entera} = entière).
\end{itemize}

\begin{exemplebox}[Analyse grammaticale du modèle (Points de vigilance)]
\begin{tabularx}{\linewidth}{|l|X|}
\hline
\rowcolor{poppurpleL}
\textbf{Point de grammaire} & \textbf{Application dans le texte} \\
\hline
\textbf{Emphase \esp{Es... quien}} & \esp{Es mi madre quien cuida...} (Sujet personne) \\
\textbf{Emphase \esp{Es lo que}} & \esp{Es lo que más quiero} (Sujet chose/idée = \emph{lo que}) \\
\textbf{Restriction \esp{Solo} + V} & \esp{Solo tengo una hermana} (Position préverbale = « Je n'ai \emph{qu'une} ») \\
\textbf{Restriction \esp{No... más que}} & \esp{No tengo más que gratitud} (Négation + \emph{más que} = « seulement ») \\
\textbf{Ser + Adj (Caractère)} & \esp{es cariñosa, es estricto, es generoso} \\
\textbf{Accord Adj (Genre/Nombre)} & \esp{cariños\textbf{a}} (fém.), \emph{estric\textbf{o}} (masc.), \emph{generos\textbf{o}} (masc.) \\
\textbf{Verbes pronominaux} & \esp{se llama, se parece, me llevo} \\
\textbf{Connecteurs logiques} & \esp{en cambio, aunque, porque, por} \\
\hline
\end{tabularx}
\end{exemplebox}

\courssub{Exercice résolu : De l'intention au texte final}

\begin{exoresolu}[Sujet : « Describa a su familia (80 palabras) »]
\textbf{Énoncé :} Vous êtes \emph{Paul}, 17 ans, lycéen à Yaoundé. Vous vivez avec votre père (comptable, calme, grand), votre mère (enseignante, dynamique, coquette), votre grand frère (étudiant, drôle, sportif) et votre petite sœur (élève, timide, mignonne). Vous vous entendez bien malgré les disputes pour la télé. Rédigez la description.

\textbf{Étape 1 : Inventaire \& Sélection (Brouillon)}
\begin{itemize}
    \item Membres : Padre, madre, hermano mayor, hermana menor. (Total 5).
    \item Adjectifs choisis : Padre (tranquilo, alto) ; Madre (dinámica, guapa) ; Hermano (divertido, deportista) ; Hermana (tímida, mona).
    \item Structures obligatoires : Emphase sur la mère (\esp{Es mi madre quien...}). Restriction sur la sœur (\esp{Solo tengo una hermana}).
\end{itemize}

\textbf{Étape 2 : Rédaction phrase par phrase avec comptage}
\begin{enumerate}
    \item \esp{Mi familia es unida y muy divertida. Somos cinco: mis padres, mi hermano mayor, mi hermana menor y yo.} (19 mots)
    \item \esp{\textbf{Es mi madre quien} anima la casa; es dinámica y guapa.} (11 mots $\rightarrow$ Total 30)
    \item \esp{Mi padre es tranquilo y alto. Mi hermano es deportista y muy divertido.} (13 mots $\rightarrow$ Total 43)
    \item \esp{\textbf{Solo tengo una hermana}, que es tímida pero muy mona.} (10 mots $\rightarrow$ Total 53)
    \item \esp{Nos llevamos bien, aunque a veces peleamos por el mando de la tele.} (13 mots $\rightarrow$ Total 66)
    \item \esp{\textbf{Es mi familia lo que} más valoro en el mundo.} (10 mots $\rightarrow$ \textbf{Total 76 mots}). \textbf{VALIDÉ (76 $\in$ [72, 88])}.
\end{enumerate}

\textbf{Texte final (76 palabras) :}
\esp{Mi familia es unida y muy divertida. Somos cinco: mis padres, mi hermano mayor, mi hermana menor y yo. Es mi madre quien anima la casa; es dinámica y guapa. Mi padre es tranquilo y alto. Mi hermano es deportista y muy divertido. Solo tengo una hermana, que es tímida pero muy mona. Nos llevamos bien, aunque a veces peleamos por el mando de la tele. Es mi familia lo que más valoro en el mundo.}
\tcblower
\textbf{Commentaire du correcteur :}
\begin{itemize}
    \item \textbf{Longueur :} 76 mots (parfait, bien dans la fourchette).
    \item \textbf{Structures :} \esp{Es mi madre quien} (Emphase sujet personne) + \esp{Solo tengo} (Restriction) + \esp{Es mi familia lo que} (Emphase sujet chose). \textbf{3 structures complexes} = Note maximale sur ce critère.
    \item \textbf{Grammaire :} Accords parfaits (\emph{dinámic\textbf{a}}, \emph{tranquil\textbf{o}}, \emph{tímid\textbf{a}}, \emph{mon\textbf{a}}). \emph{Ser} pour traits permanents. \emph{Pelear por} (se disputer pour). \emph{Valoro} (présent indicatif, 1ère pers. sg).
    \item \textbf{Vocabulaire :} Précis (\emph{mando de la tele} = télécommande, \emph{anima la casa} = anime la maison).
    \item \textbf{Cohérence :} Connecteurs (\emph{aunque, pero, y}). Conclusion emphatique forte.
\end{itemize}
\end{exoresolu}

\courssub{Relecture finale : La check-list « Zéro Faute »}

Avant de rendre votre copie, appliquez systématiquement cette grille de relecture (5 minutes).

\begin{attention}[Pièges fréquents des francophones — Check-list de relecture]
\begin{enumerate}
    \item \textbf{Accords des adjectifs :} \esp{Mi madre es alt\textbf{a}} (pas \emph{alto}) ; \esp{Mis hermanos son alt\textbf{os}} (pluriel). \textbf{Vérifiez chaque adjectif : Genre (o/a) \& Nombre (s).}
    \item \textbf{Ser vs Estar :} \esp{Mi padre es cansado} $\times$ (Il n'est pas « fatigué de nature »). $\rightarrow$ \esp{Mi padre \textbf{está} cansado} (il est fatigué maintenant) OU \esp{Mi padre \textbf{es} trabajador} (il est travailleur de nature).
    \item \textbf{Emphase : \esp{Quien} vs \esp{Que/Lo que}.}
    \begin{itemize}
        \item Antécédent \textbf{Personne} (explicite ou implicite) $\rightarrow$ \esp{Es \textbf{quien}} / \esp{Es \textbf{el que}}. Ex: \esp{Es mi padre \textbf{quien} conduce.}
        \item Antécédent \textbf{Chose / Idée / Proposition} $\rightarrow$ \esp{Es \textbf{lo que}} / \esp{Es \textbf{lo cual}}. Ex: \esp{Es \textbf{lo que} quiero.}
        \item Circonstance de temps $\rightarrow$ \esp{Fue ayer \textbf{cuando}}...
    \end{itemize}
    \item \textbf{Restriction : \esp{Solo} vs \esp{Solamente} vs \esp{No... más que}.}
    \begin{itemize}
        \item \esp{Solo tengo un hermano} = Je n'ai qu'un frère (Position avant verbe = restriction sur l'objet).
        \item \esp{Tengo solo un hermano} = Je n'ai \textbf{qu'un} frère (insistance sur le chiffre 1).
        \item \esp{No tengo más que un hermano} = Registre plus formel, correct.
        \item \textbf{Erreur fatale :} \esp{No tengo solo un hermano} = « Je n'ai \textbf{pas seulement} un frère » (j'en ai d'autres) $\neq$ « Je n'ai \textbf{qu'un} frère ».
    \end{itemize}
    \item \textbf{Faux amis « Famille » :}
    \begin{itemize}
        \item \esp{Parientes} = Parents (au sens large : cousins, oncles...), \textbf{PAS} « parents » (père/mère). Parents = \esp{Padres}.
        \item \esp{Padres} = Père + Mère. Ne pas écrire \esp{mis padres y mi madre} (redondant).
        \item \esp{Hermanos} = Frères + Sœurs. Ne pas écrire \esp{mis hermanos y mis hermanas}.
    \end{itemize}
    \item \textbf{Comptage final :} Recomptez \textbf{à voix basse} en pointant chaque mot au stylo. Visez 75--85.
\end{enumerate}
\end{attention}

\courssub{Entraînement : À vous de jouer !}

\begin{exercice}[Sujet d'entraînement Bac]
\textbf{Consigne :} \esp{« Describa a su familia (80 palabras). »}
\textbf{Profil suggéré :} Vous vivez à Douala. Famille de 4 personnes (père chauffeur, mère commerçante au marché, vous, petite sœur). Père : costaud, drôle. Mère : courageuse, belle. Sœur : étudieuse, ressemble à la mère. Vous : sportif. Ambiance : bruyante mais aimante. Contrainte : Utiliser \esp{Es mi padre quien...} et \esp{No... más que...}.

\textbf{Consignes de travail :}
\begin{enumerate}
    \item Faites votre liste de vocabulaire (métiers, adjectifs).
    \item Appliquez la méthode en 3 étapes (Présentation $\rightarrow$ Portraits $\rightarrow$ Relations).
    \item Rédigez sur brouillon, comptez, ajustez.
    \item Recopiez au net en vérifiant la check-list « Zéro Faute ».
\end{enumerate}
\end{exercice}

\begin{corrige}[Exercice : Describa a su familia]
\textbf{Proposition de rédaction (76 palabras) :}
\esp{Mi familia es pequeña pero muy ruidosa y alegre. Somos cuatro: mis padres, mi hermanita y yo. \textbf{Es mi padre quien} nos hace reír cada día; es fuerte y muy divertido. Mi madre es valiente y guapa; vende comida en el mercado. \textbf{No tengo más que una hermana}, se llama Grace, es estudiosa y se parece a mi madre. Yo soy deportista. Aunque hay mucho ruido, nos queremos mucho. \textbf{Es el amor lo que} nos une.}

\textbf{Analyse :}
\begin{itemize}
    \item \textbf{Mots :} 76 (Conforme).
    \item \textbf{Emphase 1 :} \esp{Es mi padre quien nos hace reír} (Sujet personne).
    \item \textbf{Restriction :} \esp{No tengo más que una hermana} (Forme formelle correcte).
    \item \textbf{Emphase 2 :} \esp{Es el amor lo que nos une} (Sujet chose \emph{lo que}).
    \item \textbf{Métiers :} \esp{chofer} (chauffeur) $\rightarrow$ \esp{vende comida en el mercado} (commerçante marché).
    \item \textbf{Adjectifs :} \emph{fuerte, divertido, valiente, guapa, estudiosa, deportista} (Accords vérifiés).
    \item \textbf{Connecteurs :} \emph{pero, aunque, y}.
\end{itemize}
\textbf{Variante possible pour « père chauffeur » :} \esp{Mi padre trabaja como taxista / conductor.} (Évitez \esp{chofer} seul qui peut prêter à confusion selon les pays, \esp{conductor} ou \esp{taxista} plus clair).
\end{corrige}

\begin{savaistu}[Le « Peso » des mots au Bac camerounais]
Au Cameroun, l'épreuve d'expression écrite (Redacción) est notée sur 10 à 15 points selon les académies. La grille de correction réserve généralement :
\begin{itemize}
    \item 4--5 pts : Respect de la consigne (sujet, nombre de mots \textbf{72--88}).
    \item 4--5 pts : Correction grammaticale (accords, temps, \emph{ser/estar}, structures complexes).
    \item 2--3 pts : Richesse lexicale et connecteurs.
    \item 1--2 pts : Cohérence et présentation.
\end{itemize}
\textbf{Conseil de pro :} Un texte de 79 mots parfait grammaticalement avec 2 emphases et 1 restriction aura \textbf{toujours} une meilleure note qu'un texte de 100 mots bourré de fautes d'accords ou hors sujet. \textbf{La concision est une vertu rhétorique.}
\end{savaistu}

\newpage

\coursec{Activité d'intégration}

\coursec{Leçon E04 : Traduction de « C'est… que », « C'est… qui », « Ne… que » et de l'emphase ; redacción : describir a su familia (80 palabras)}

\courssub{Observation : Découvrir l'emphase en contexte authentique}
\begin{textebox}[Extrait du blog « Mi raíz, mi ala » — María Nsue, étudiante à Malabo (Guinée équatoriale)]
\esp{Mi abuela Pilar es el alma de la familia. \textbf{Es ella quien} guarda las recetas de la abuela bisabuela. \textbf{Fue en su cocina donde} yo aprendí a hacer el \textit{pepesup} con dieciséis años. Mi padre trabaja en la exportación de cacao; \textbf{solo} él viaja a España cada trimestre. \textbf{No tenemos más que} un coche viejo, pero \textbf{no hay más que} alegría cuando nos reunimos. \textbf{Es el domingo cuando} la casa se llena de primos, tíos y risas.}
\end{textebox}

\begin{center}
\textbf{Glossaire / Traduction}
\end{center}
\begin{itemize}
    \item \esp{el alma de la familia} : l'âme de la famille (pilier central).
    \item \esp{la bisabuela} : l'arrière-grand-mère.
    \item \esp{exportación de cacao} : exportation de cacao (secteur clé au Cameroun et en Guinée équatoriale).
    \item \esp{cada trimestre} : chaque trimestre.
    \item \esp{pepesup} : plat traditionnel (soupe de poisson/viande, épices, légumes-feuilles) partagé Cameroun / Guinée équatoriale.
\end{itemize}

\begin{experience}[Analyse guidée : Repérer les structures d'insistance]
Lisez le texte et répondez en français :
\begin{enumerate}
    \item Quels sont les trois membres de la famille mis en avant ? Par quelle tournure espagnole ? (\textit{Indice : cherchez \esp{Es... quien}, \esp{Fue... donde}, \esp{solo}, \esp{no... más que}}).
    \item Traduisez littéralement : \esp{Es ella quien guarda...} ; \esp{Fue en su cocina donde...} ; \esp{solo él viaja...} ; \esp{No tenemos más que...}.
    \item Quelle différence entendez-vous entre \esp{Es ella quien...} (présent) et \esp{Fue en su cocina donde...} (passé) ?
\end{enumerate}
\end{experience}

\courssub{Traduire « C'est… qui » : l'emphase sur le sujet (\esp{Ser + quien / el que / lo que})}
\begin{propriete}[Règle fondamentale : La phrase clivée sujet]
Pour traduire \textbf{« C'est [X] qui + verbe »} (mise en relief du sujet réel), l'espagnol utilise une phrase clivée : \\
\esp{Ser (conjugué) + Élément mis en relief + QUIEN / EL QUE / LA QUE / LOS QUE / LAS QUE / LO QUE + Verbe conjugué}.
\begin{itemize}
    \item \textbf{\esp{Quien}} (invariable, personne, antécédent souvent unique, indéfini ou nom propre) : \esp{\textbf{Es quien} me ayuda.} (C'est lui/elle qui m'aide). \esp{\textbf{Es mi hermano quien} llama.} (C'est mon frère qui appelle).
    \item \textbf{\esp{El que / La que / Los que / Las que}} (accord en genre/nombre avec l'antécédent explicite) : \esp{Es mi madre \textbf{la que} cocina.} (C'est ma mère qui cuisine). \esp{Son mis padres \textbf{los que} deciden.} (Ce sont mes parents qui décident).
    \item \textbf{\esp{Lo que}} (chose, idée, action, infinitif) : \esp{Es \textbf{lo que} quiero.} (C'est ce que je veux). \esp{Es \textbf{lo que} me duele.} (C'est ce qui me fait mal).
\end{itemize}
\cle{Règle d'accord} : Le verbe final s'accorde avec l'antécédent réel (le sujet logique). \esp{Es mi padre el que \textbf{viene}} (sg) ; \esp{Son mis hermanos los que \textbf{vienen}} (pl).
\end{propriete}

\begin{exemplebox}[Exemples traduits — Thème (FR $\rightarrow$ ES)]
\begin{enumerate}
    \item C'est mon oncle qui paie. $\rightarrow$ \esp{Es mi tío \textbf{el que} paga.}
    \item Ce sont mes cousins qui viennent. $\rightarrow$ \esp{Son mis primos \textbf{los que} vienen.}
    \item C'est moi qui ai tort. $\rightarrow$ \esp{\textbf{Soy yo quien} me equivoco.} (\textit{Note : \esp{quien} pour 1ère/2ème pers., verbe à la 1ère pers.})
    \item C'est ce livre que je veux. $\rightarrow$ \esp{Es \textbf{lo que} quiero.} (\textit{Chose $\rightarrow$ \esp{lo que}})
\end{enumerate}
\end{exemplebox}

\begin{attention}[Pièges francophones — Emphase sujet]
\begin{itemize}
    \item \textbf{Confusion \esp{quien} / \esp{que}} : \esp{*Es mi padre que viene} $\rightarrow$ \textbf{INCORRECT}. Après nom de personne explicite + emphase sujet : \esp{el que / la que} (ou \esp{quien} si antécédent unique/indéfini, mais \esp{el que} est plus sûr avec antécédent explicite). \esp{Es mi padre \textbf{el que} viene.}
    \item \textbf{Oubli de l'accord} : \esp{*Son mis padres el que vienen} $\rightarrow$ \esp{Son mis padres \textbf{los que} \textbf{vienen}}.
    \item \textbf{Calque « C'est... que » pour le sujet} : \esp{*Es mi padre que viene} (anglicisme/calque). En espagnol, \esp{que} seul ne fait pas l'emphase sujet ; il faut \esp{quien / el que}.
    \item \textbf{Position de l'accent} : \esp{Es mi aBUEla quien...} (accent d'intensité sur l'élément mis en relief).
\end{itemize}
\end{attention}

\begin{exercice}[Entraînement 1 : Thème — Emphase sujet]
Traduisez en espagnol :
\begin{enumerate}
    \item C'est ma grand-mère maternelle qui raconte les histoires.
    \item Ce sont les aînés qui décident.
    \item C'est toi qui as raison.
    \item C'est ce plat (el plato) que je préfère.
\end{enumerate}
\end{exercice}
\begin{corrige}[Exercice 1]
\begin{enumerate}
    \item \esp{Es mi abuela materna \textbf{la que} cuenta las historias.} (ou \esp{quien cuenta})
    \item \esp{Son los mayores \textbf{los que} deciden.}
    \item \esp{\textbf{Eres tú quien} tiene razón.} (\esp{quien} + verbe 3e pers. sg car \esp{quien} = « la personne qui » ; ou \esp{Eres tú el que tienes razón}).
    \item \esp{Es \textbf{lo que} prefiero.}
\end{enumerate}
\end{corrige}

\courssub{Traduire « C'est… que » : l'emphase sur les circonstances (\esp{Fue/Es + cuando/donde/como})}
\begin{propriete}[Règle fondamentale : La phrase clivée circonstancielle]
Pour traduire \textbf{« C'est [circonstance] que + indicatif »} (temps, lieu, manière), l'espagnol utilise : \\
\esp{Ser (conjugué) + Circonstance mise en relief + CUANDO / DONDE / COMO + Indicatif}.
\begin{itemize}
    \item \textbf{Temps (\esp{Cuando})} : \esp{\textbf{Fue ayer} \textbf{cuando} \textbf{llegó}.} (C'est hier qu'il est arrivé). \esp{\textbf{Es ahora} \textbf{cuando} \textbf{estudiamos}.} (C'est maintenant qu'on étudie).
    \item \textbf{Lieu (\esp{Donde})} : \esp{\textbf{Fue en Yaoundé} \textbf{donde} \textbf{nos conocimos}.} (C'est à Yaoundé que nous nous sommes rencontrés). \esp{\textbf{Es aquí} \textbf{donde} \textbf{vivo}.} (C'est ici que j'habite).
    \item \textbf{Manière (\esp{Como})} : \esp{\textbf{Fue así} \textbf{como} \textbf{lo hizo}.} (C'est ainsi qu'il l'a fait). \esp{\textbf{Es como} \textbf{como} \textbf{se hace}.} (C'est comme ça qu'on fait).
\end{itemize}
\cle{Choix du temps de \esp{Ser}} : \esp{Fue} (prétérit indéfini) pour un événement ponctuel, achevé, daté dans le passé. \esp{Era} (imparfait) pour une description de fond, une habitude passée (\esp{Era en verano cuando íbamos a la playa}). \esp{Es} pour le présent/général.
\end{propriete}

\begin{exemplebox}[Exemples traduits — Thème \& Version]
\begin{tabularx}{\linewidth}{|X|X|}
\hline
\textbf{Français} & \textbf{Espagnol} \\
\hline
C'est en 2020 que nous sommes arrivés. & \esp{\textbf{Fue en 2020} \textbf{cuando} \textbf{llegamos}.} \\
\hline
C'est à Douala que j'ai grandi. & \esp{\textbf{Fue en Douala} \textbf{donde} \textbf{crecí}.} \\
\hline
C'est ainsi qu'elle chante. & \esp{\textbf{Es así} \textbf{como} \textbf{canta}.} \\
\hline
C'était l'été que nous allions à la plage. & \esp{\textbf{Era en verano} \textbf{cuando} \textbf{íbamos} a la playa.} \\
\hline
\end{tabularx}
\end{exemplebox}

\begin{attention}[Pièges francophones — Emphase circonstancielle]
\begin{itemize}
    \item \textbf{\esp{Fue} vs \esp{Era}} : \esp{*Era ayer cuando llegó} $\rightarrow$ \textbf{FAUX}. Action ponctuelle datée $\rightarrow$ \esp{Fue}. \esp{Era} = description d'époque / habitude.
    \item \textbf{Subjonctif interdit} : Après \esp{cuando/donde/como} en emphase sur fait réel (passé/présent), on met l'\textbf{indicatif}. \esp{*Fue ayer cuando \textbf{llegue}} $\rightarrow$ \esp{Fue ayer cuando \textbf{llegó}}.
    \item \textbf{Préposition de lieu} : \esp{Fue \textbf{en} Madrid donde...} (pas \esp{a Madrid} pour \esp{donde} statif/événementiel, bien que \esp{a} soit possible avec verbes de mouvement, \esp{en} est standard pour « c'est à... que »).
\end{itemize}
\end{attention}

\begin{exercice}[Entraînement 2 : Version \& Thème — Emphase circonstancielle]
\begin{enumerate}
    \item Traduisez : \esp{Fue el año pasado cuando mi padre compró la casa.}
    \item Traduisez en espagnol : « C'est chez ma tante (en casa de mi tía) que nous fêtons Noël. »
    \item Traduisez en espagnol : « C'était ainsi que ma grand-mère cuisinait. » (habitude)
\end{enumerate}
\end{exercice}
\begin{corrige}[Exercice 2]
\begin{enumerate}
    \item « C'est l'an dernier que mon père a acheté la maison. »
    \item \esp{\textbf{Es en casa de mi tía} \textbf{donde} \textbf{celebramos} la Navidad.}
    \item \esp{\textbf{Era así} \textbf{como} \textbf{cocinaba} mi abuela.} (Imparfait \esp{Era} + \esp{cocinaba} pour habitude passée).
\end{enumerate}
\end{corrige}

\courssub{Traduire « Ne… que » : la restriction (\esp{Solo}, \esp{No... más que}, \esp{No... sino})}
\begin{propriete}[Règle fondamentale : Restriction / Exclusion]
Le français « ne... que » (seulement) se traduit par plusieurs structures selon la place de l'élément restreint et le registre.
\begin{center}
\begin{tabularx}{\linewidth}{|X|X|X|}
\hline
\rowcolor{popblueL} \textbf{Français} & \textbf{Espagnol} & \textbf{Structure / Remarque} \\
\hline
Il ne mange \textbf{que} du riz. & \esp{Come \textbf{solo} arroz.} / \esp{Come \textbf{únicamente} arroz.} & \esp{Solo / Únicamente} (adv) \textbf{devant} l'élément restreint. Pas de \esp{no}. \\
\hline
Il \textbf{n'}a \textbf{que} 20 ans. & \esp{Tiene \textbf{solo} 20 años.} & \esp{Tener} + \esp{solo} (pas de \esp{no}). \\
\hline
Je ne veux \textbf{que} toi. & \esp{No quiero \textbf{más que} a ti.} & \esp{No... más que} encadre le verbe et l'élément. Standard. \\
\hline
Je ne veux \textbf{que} toi. & \esp{No quiero \textbf{sino} a ti.} & \esp{No... sino} (formel, littéraire) = « rien d'autre que ». \\
\hline
Il n'y a \textbf{que} moi. & \esp{\textbf{Solo} estoy yo.} / \esp{No hay \textbf{más que} yo.} & \esp{Solo} en tête = « Seulement moi suis là ». \\
\hline
Il ne \textbf{fait} \textbf{que} dormir. & \esp{No \textbf{hace} \textbf{más que} dormir.} & \esp{No hacer más que} + infinitif = « ne faire que ». \\
\hline
\end{tabularx}
\end{center}
\cle{Orthographe RAE (2010)} : \esp{Solo} (adverbe = seulement) s'écrit \textbf{sans accent}. \esp{Solo} (adjectif = seul) s'écrit \textbf{sans accent}. L'ancien \esp{sólo} (adv) est toléré mais déconseillé. \esp{Solamente} = \esp{solo} (plus long, insistant).
\end{propriete}

\begin{exemplebox}[Exemples contrastés — Attention au \esp{No} !]
\begin{itemize}
    \item \esp{Solo Juan vino.} = Seul Juan est venu (les autres non). \textbf{Sujet restreint}.
    \item \esp{Juan solo vino.} = Juan n'a fait que venir (il n'est pas resté, n'a pas mangé...). \textbf{Verbe restreint}.
    \item \esp{No vino más que Juan.} = Il n'est venu que Juan. \textbf{Négation double obligatoire} (\esp{No} + \esp{más que}).
    \item \esp{No vino sino Juan.} = Il n'est venu que Juan (insistance sur l'unicité).
\end{itemize}
\end{exemplebox}

\begin{attention}[Pièges francophones — Restriction]
\begin{itemize}
    \item \textbf{Oubli du \esp{No}} : \esp{*Hubo más que alegría} $\rightarrow$ \textbf{FAUX}. Il faut \esp{\textbf{No} hubo más que alegría}. Le français « ne... que » cache une négation ; l'espagnol \esp{no... más que} la rend visible.
    \item \textbf{Confusion \esp{solo} (adv) / \esp{solo} (adj)} : \esp{Estoy solo} = Je suis seul (adjectif, accord : \esp{sola}). \esp{Solo como arroz} = Je ne mange que du riz (adverbe, invariable).
    \item \textbf{\esp{No... sino} + verbe} : Si l'élément restreint est un verbe, \esp{sino} appelle l'infinitif (souvent avec \esp{que}) : \esp{No hizo sino \textbf{que} llorar} (Il n'a fait que pleurer). Mais \esp{No hace más que dormir} est plus courant.
\end{itemize}
\end{attention}

\begin{exercice}[Entraînement 3 : Traduction — Restriction]
Traduisez en espagnol :
\begin{enumerate}
    \item Je n'ai que deux frères.
    \item Nous ne mangeons que du manioc (yuca) le dimanche.
    \item Il ne fait que dormir.
    \item Il n'y a que ma mère qui parle espagnol. (Combiner emphase + restriction)
\end{enumerate}
\end{exercice}
\begin{corrige}[Exercice 3]
\begin{enumerate}
    \item \esp{Tengo \textbf{solo} dos hermanos.} (ou \esp{No tengo más que dos hermanos}).
    \item \esp{Comemos \textbf{solo} yuca el domingo.} / \esp{No comemos \textbf{más que} yuca el domingo.}
    \item \esp{No \textbf{hace} \textbf{más que} dormir.}
    \item \esp{\textbf{Solo} mi madre habla español.} / \esp{No habla español \textbf{más que} mi madre.} / \esp{Es mi madre \textbf{la que} habla español \textbf{y no hay más que ella}.} (Plusieurs solutions).
\end{enumerate}
\end{corrige}

\courssub{Lexique : La famille élargie et le portrait (physique et moral)}
\begin{definition}[Vocabulaire essentiel — Mots-outils pour la rédaction]
\textbf{Parents proches :} \esp{padre, madre, padres} (parents), \esp{hermano, hermana, hermanos} (frères/sœurs), \esp{hijo, hija, hijos} (enfants), \esp{abuelo, abuela, abuelos} (grands-parents), \esp{nieto, nieta} (petits-enfants), \esp{bisabuelo/a} (arrière-grand-parent).
\textbf{Parents alliés / famille recomposée :} \esp{padrastro, madrastra} (beau-père, belle-mère), \esp{hermanastro/a} (demi-frère/sœur), \esp{cuñado/a} (beau-frère/belle-sœur), \esp{yerno, nuera} (gendre, bru), \esp{suegro/a} (beau-père/belle-mère \textit{par alliance}).
\textbf{Collatéraux :} \esp{tío, tía} (oncle, tante), \esp{primo, prima} (cousin/cousine), \esp{sobrino, sobrina} (neveu/nièce), \esp{cuñado/a} (conjoint du frère/sœur).
\textbf{Portrait physique :} \esp{alto/bajo, delgado/grueso, canoso, calvo, de complexión atlética/fuerte, ojos almendrados, piel morena/clara, barbado, bigotudo}.
\textbf{Portrait moral / Rôle :} \esp{trabajador/a, cariñoso/a, estricto/a, alegre, serio/a, generoso/a, pilar de la familia, cabeza de familia, alma de la casa, apoyo incondicional, ejemplo a seguir}.
\textbf{Verbes utiles :} \esp{convivir, cuidar, mantener (tradiciones), reunir(se), celebrar, heredar, transmitir, parecerse (a), llevarse bien (con)}.
\end{definition}

\begin{savaistu}[La famille en Guinée équatoriale et au Cameroun : contextes partagés]
En Guinée équatoriale (seul pays africain hispanophone), la structure familiale mélange traditions bantoues (Fang, Bubi, Ndowe) et héritage espagnol. Le \esp{linaje} (lignage) et le \esp{clan} sont centraux. On vit souvent en \esp{familia extensa} (famille élargie) sous le même toit ou dans la même concession. Le respect des \esp{mayores} (aînés) est absolu : \esp{« El mayor manda, el menor obedece »}. Au Cameroun anglophone/francophone voisin, les réalités sont proches : famille nombreuse, rôle clé de la \esp{abuela} (grand-mère) comme gardienne des traditions orales (\esp{cuentos, refranes}), solidarité intergénérationnelle. Le \esp{pepesup} (ou \esp{pep soup}), le \esp{kwacoco}, le \esp{fufu} sont des marqueurs culinaires partagés. Intégrer ce vocabulaire culturel (\esp{fang, clan, linaje, abuela, pepesup}) ancre votre espagnol dans \textbf{votre} réalité.
\end{savaistu}

\begin{exercice}[Lexique en action : Préparation à la rédaction]
Complétez en espagnol :
\begin{enumerate}
    \item Mon beau-frère (sœur de ma femme) = \esp{mi \_\_\_\_\_\_\_\_}.
    \item Ma belle-mère (mère de mon mari) = \esp{mi \_\_\_\_\_\_\_\_}.
    \item Mon neveu (fils de ma sœur) = \esp{mi \_\_\_\_\_\_\_\_}.
    \item Adjectif : « qui a les cheveux blancs » = \esp{\_\_\_\_\_\_\_\_}.
    \item Nom : « le pilier de la famille » = \esp{\_\_\_\_\_\_\_\_ de la familia}.
\end{enumerate}
\end{exercice}
\begin{corrige}[Exercice 4]
\begin{enumerate}
    \item \esp{cuñado} ; \item \esp{suegra} ; \item \esp{sobrino} ; \item \esp{canoso} ; \item \esp{pilar}.
\end{enumerate}
\end{corrige}

\courssub{Activité d'intégration : Projet « Raíces y Alas » — Redacción final (80 palabras)}

\courssub{Objectifs de l'activité}
Mobiliser l'ensemble des acquis de la leçon dans une tâche complexe et authentique : produire une description écrite de sa famille (environ 80 mots) en intégrant obligatoirement des structures d'emphase (\esp{es... quien}, \esp{fue... cuando}) et de restriction (\esp{solo}, \esp{no... más que}), puis interagir à l'oral pour valider la compréhension et la traduction (mini-version) de ces structures chez un pair. L'activité vise à consolider la compétence \emph{Écrire et réagir par écrit} et \emph{Parler en interaction}, conformément au programme MINESEC Terminale A (Module 1 : Vie familiale et intégration sociale).

\courssub{Situation}
\begin{textebox}[Contexte : Le projet « Raíces y Alas » au Lycée Bilingue de Yaoundé]
Dans le cadre d'un jumelage numérique entre le Lycée Bilingue de Yaoundé (Cameroun) et l'Instituto Cervantes de Malabo (Guinée équatoriale), la classe de Terminale A participe au projet \emph{« Raíces y Alas »} (Racines et Ailes). Chaque élève doit rédiger une fiche d'identité familiale pour un mur collaboratif virtuel (\emph{Padlet} de la classe). L'objectif est de présenter sa famille sous un angle personnel et émouvant, en mettant en valeur un membre clé, un souvenir marquant ou une particularité familiale, tout en respectant une contrainte de longueur (80 mots \textpm 10\%) et des contraintes grammaticales précises. Les meilleures fiches seront imprimées pour former un grand \emph{árbol genealógico} collectif exposé au CDI lors de la \emph{Semana de la Hispanidad}.
\end{textebox}

\courssub{Consignes}
\begin{enumerate}
    \item \textbf{Préparation individuelle (15 min) :} Rédigez votre description de famille en espagnol (80 mots \textpm 10\%). Votre texte doit impérativement contenir :
    \begin{itemize}
        \item Au moins \textbf{une tournure d'emphase sur le sujet} avec \esp{Es... quien} ou \esp{Es... el que} (ex: \esp{Es mi abuela quien mantiene unidas las tradiciones}).
        \item Au moins \textbf{une tournure d'emphase sur le temps/les circonstances} avec \esp{Fue... cuando} (ex: \esp{Fue el domingo pasado cuando celebramos su cumpleaños}).
        \item Au moins \textbf{une restriction} avec \esp{solo}, \esp{únicamente} ou \esp{no... más que} / \esp{no... sino} (ex: \esp{Solo mi padre habla español} ; \esp{No tenemos más que un coche}).
        \item Le lexique de la parenté (au moins 6 termes différents : \esp{abuelo, tío, primo, cuñado, sobrino, padrastro}...), des adjectifs de portrait physique (\esp{alto, canoso, de complexión fuerte}) et moral (\esp{trabajador, cariñoso, estricto, alegre}).
    \end{itemize}
    \item \textbf{Interaction orale par binômes (10 min) :} Asseyez-vous face à votre partenaire. L'élève A lit son texte à voix normale (sans le montrer). L'élève B écoute et note en \textbf{français} deux phrases entendues qui contiennent une structure d'emphase ou de restriction (mini-version). Inversez les rôles.
    \item \textbf{Vérification et correction (5 min) :} Échangez vos feuilles de version. Relisez le texte original de votre partenaire (qu'il vous montre maintenant) et vérifiez l'exactitude de votre traduction française. Signalez toute erreur de structure (ex: confusion \esp{quien}/\esp{que}, oubli du \esp{no} dans \esp{no... más que}) ou de vocabulaire.
    \item \textbf{Finalisation et dépôt (5 min) :} Corrigez votre texte espagnol si nécessaire (orthographe, accords, structures imposées). Déposez la version finale sur le Padlet de la classe ou remettez-la au professeur pour l'affichage de l'\emph{árbol genealógico}.
\end{enumerate}

\courssub{Ressources à mobiliser (Fiche mémo)}
\begin{propriete}[Rappel : Emphase sur le sujet — \esp{Es... quien / el que / lo que}]
Pour traduire \textbf{« C'est... qui »} (sujet réel), l'espagnol utilise la tournure clivée : \esp{Ser + elemento mis en relief + quien / el que / la que / los que / las que / lo que + verbo conjugado}.
\begin{itemize}
    \item \esp{Quien} (invariable, personne, antécédent souvent indéfini ou non exprimé) : \esp{Es \textbf{quien} me ayuda.} (C'est lui qui m'aide).
    \item \esp{El que / La que / Los que / Las que} (accord en genre/nombre avec l'antécédent explicite) : \esp{Es mi madre \textbf{la que} cocina.} (C'est ma mère qui cuisine).
    \item \esp{Lo que} (chose, idée, proposition infinitive) : \esp{Es \textbf{lo que} quiero.} (C'est ce que je veux).
\end{itemize}
\textbf{Attention :} Le verbe final s'accorde avec l'antécédent réel (\esp{Es mi padre el que \textbf{viene}} ; \esp{Son mis hermanos los que \textbf{vienen}}).
\end{propriete}

\begin{propriete}[Rappel : Emphase sur le temps / les circonstances — \esp{Fue... cuando / donde / como}]
Pour traduire \textbf{« C'est... que »} (circonstance de temps, lieu, manière), on utilise le passé (\esp{Fue / Fue ayer / Fue en 2010}) + \esp{cuando / donde / como} + indicatif (fait réel passé).
\begin{itemize}
    \item Temps : \esp{\textbf{Fue ayer} \textbf{cuando} \textbf{llegó}.} (C'est hier qu'il est arrivé).
    \item Lieu : \esp{\textbf{Fue en Yaoundé} \textbf{donde} \textbf{nos conocimos}.} (C'est à Yaoundé que nous nous sommes rencontrés).
    \item Manière : \esp{\textbf{Fue así} \textbf{como} \textbf{lo hizo}.} (C'est ainsi qu'il l'a fait).
\end{itemize}
\textbf{Variante présente :} \esp{Es ahora cuando...} (C'est maintenant que...), \esp{Es aquí donde...}. \textbf{Variante imparfait (habitude) :} \esp{Era en verano cuando íbamos...}
\end{propriete}

\begin{propriete}[Rappel : Restriction — \esp{Solo / Únicamente / No... más que / No... sino}]
Pour traduire \textbf{« Ne... que »} (seulement), plusieurs options selon le registre et la position :
\begin{center}
\begin{tabularx}{\linewidth}{|X|X|X|}
\hline
\rowcolor{popblueL} \textbf{Français} & \textbf{Espagnol} & \textbf{Remarque} \\
\hline
Il ne mange \textbf{que} du riz. & \esp{Come \textbf{solo} arroz.} / \esp{Come \textbf{únicamente} arroz.} & \esp{Solo} / \esp{únicamente} devant l'élément restreint (adv). \\
\hline
Il \textbf{n'}a \textbf{que} 20 ans. & \esp{Tiene \textbf{solo} 20 años.} & \esp{Tener} + \esp{solo} (pas de \esp{no}). \\
\hline
Je ne veux \textbf{que} toi. & \esp{No quiero \textbf{más que} a ti.} / \esp{No quiero \textbf{sino} a ti.} & \esp{No... más que} / \esp{No... sino} (formel/littéraire) encadrent l'élément. \\
\hline
Il n'y a \textbf{que} moi. & \esp{\textbf{Solo} estoy yo.} / \esp{No hay \textbf{más que} yo.} & \esp{Solo} en début de phrase = \emph{Seulement moi suis là}. \\
\hline
\end{tabularx}
\end{center}
\cle{Attention} : \esp{Solamente} = \esp{solo} (adv). \esp{Solo} (sans accent, adj) = \esp{seul} (\esp{Estoy solo} = Je suis seul). Depuis 2010, la RAE recommande \textbf{pas d'accent} sur \esp{solo} (adv), mais l'usage \esp{sólo} persiste. Nous acceptons les deux.
\end{propriete}

\courssub{Proposition de production et corrigé commenté}
\begin{exoresolu}[Production modèle : Fiche de Martín (élève camerounais d'origine équato-guinéenne)]
\textbf{Énoncé :} Rédigez la description de votre famille (80 mots) en intégrant : 1 \esp{es... quien}, 1 \esp{fue... cuando}, 1 restriction (\esp{solo/no... más que}).

\tcblower

\textbf{Production espagnole (82 mots) :}

\esp{Mi familia es numerosa y muy unida. Vivimos en el barrio de Mvog-Ada, en Yaoundé. \textbf{Es mi abuela materna, doña Pilar, \textbf{quien} guarda las tradiciones fang y nos cuenta cuentos cada noche.} Mi padre, ingeniero, trabaja en una empresa de cacao; \textbf{solo} él habla español perfecto en casa. \textbf{Fue en diciembre pasado} \textbf{cuando} celebramos el 80º cumpleaños de mi abuelo: \textbf{no hubo} \textbf{más que} alegría y platos típicos como el \textit{pepesup}. Mis tres hermanos y yo estudiamos en el liceo bilingüe. Somos un equipo.}

\textbf{Commentaire des choix de langue :}

\begin{enumerate}
    \item \textbf{Emphase sujet (\esp{Es... quien}) :} \esp{Es mi abuela materna, doña Pilar, quien guarda...} $\rightarrow$ Traduction de « C'est ma grand-mère maternelle... qui garde... ». \esp{Quien} est utilisé car l'antécédent (\esp{doña Pilar}) est une personne individualisée, mise en relief comme unique gardienne des traditions. L'accord du verbe \esp{guarda} (3e pers. sg) est correct.
    \item \textbf{Restriction (\esp{solo}) :} \esp{solo él habla español perfecto} $\rightarrow$ « Lui seul parle espagnol parfait ». \esp{Solo} placé devant le sujet \esp{él} restreint l'identité du locuteur. Alternative possible : \esp{No habla español perfecto más que él} (plus lourd).
    \item \textbf{Emphase temps (\esp{Fue... cuando}) :} \esp{Fue en diciembre pasado cuando celebramos...} $\rightarrow$ « C'est le mois de décembre dernier que nous avons fêté... ». \esp{Fue} (prétérit indéfini de \esp{ser} pour un événement ponctuel achevé) + \esp{cuando} + indicatif (\esp{celebramos}, pretérito indefinido). Structure canonique. \textit{Amélioration stylistique : « Fue en diciembre pasado » plus fluide que « Fue el pasado mes de diciembre »}.
    \item \textbf{Restriction (\esp{no... más que}) :} \esp{no hubo más que alegría} $\rightarrow$ « Il n'y eut que de la joie ». \esp{Hubo} (hay au passé) + \esp{no... más que} + nom. Notez l'absence d'article devant \esp{alegría} (concept abstrait non comptable). \esp{No hubo sino alegría} aurait été possible (registre plus littéraire).
    \item \textbf{Lexique \& Contexte :} \esp{Mvog-Ada} (quartier réel de Yaoundé), \esp{pepesup} (plat camerounais/équato-guinéen, en italique), \esp{fang} (ethnie/langue), \esp{liceo bilingüe} (adaptation locale de \emph{lycée bilingue}, terme utilisé en Guinée équatoriale). L'accord \esp{numerosa y unida} (fém. sg. accordé à \esp{familia}) est respecté.
    \item \textbf{Connecteurs \& Ponctuation :} \esp{y, ;} (ponctuation forte pour séparer l'emphase), \esp{:} (introduit le détail de l'événement emphatique). Respect des accents : \esp{Yaundé, papá, mamá, inglés, cumpleaños}.
\end{enumerate}

\textbf{Version française (mini-version attendue de l'élève B) :}
\begin{enumerate}
    \item « C'est ma grand-mère maternelle, doña Pilar, \textbf{qui} garde les traditions fang et nous raconte des contes chaque soir. »
    \item « \textbf{C'est en décembre dernier que} nous avons fêté l'anniversaire de mon grand-père : \textbf{il n'y eut que} de la joie et des plats typiques comme le \textit{pepesup}. »
\end{enumerate}
\end{exoresolu}

\courssub{Bilan de l'activité}
Cette activité d'intégration a permis de vérifier l'appropriation des trois piliers grammaticaux de la leçon dans un contexte de production personnelle et émotionnellement investie (la famille).
\begin{itemize}
    \item \textbf{Emphase sujet (\esp{es... quien/el que}) :} La plupart des élèves ont correctement identifié le membre clé à mettre en relief (souvent la mère ou la grand-mère, figure centrale au Cameroun comme en zone hispanophone). L'erreur résiduelle reste l'accord du verbe final (ex: \esp{*Es mis padres quien vienen} $\rightarrow$ corriger : \esp{Son mis padres \textbf{quienes vienen}} ou \esp{Es mi padre \textbf{el que viene}}). Rappel : \esp{quien} $\rightarrow$ \esp{quienes} au pluriel.
    \item \textbf{Emphase circonstancielle (\esp{fue... cuando}) :} Bien maîtrisée pour le passé narratif (\esp{Fue} + indéfini). Attention à ne pas confondre avec \esp{era... cuando} (imparfait = description de fond / habitude passée : \esp{Era en verano cuando íbamos a la playa}). Veiller à l'indicatif après \esp{cuando} (fait réel).
    \item \textbf{Restriction (\esp{solo / no... más que}) :} L'usage de \esp{solo} (adv) en position préverbale ou pré-nominale est acquis. La structure \esp{no... más que} pose encore problème par calque du français \emph{ne... que} (oubli du \esp{no} : \esp{*hubo más que alegría}). Le rappel de la négation double obligatoire (\esp{no}) est crucial. Distinguer \esp{solo} (adv, invariable) et \esp{solo/a} (adj, accordé : \esp{Estoy solo}).
    \item \textbf{Lexique \& Orthographe :} L'effort pour intégrer du vocabulaire précis (\esp{padrastro, cuñado, canoso, pilar de la familia}) est visible. Vigilance sur les accents (\esp{papá, mamá, abuelo, inglés, cumpleaños, Yaundé}) et la ponctuation espagnole (\esp{¿ ? ¡ !}) dans les versions finales destinées à l'affichage.
\end{itemize}
L'échange oral (mini-version) a révélé l'importance de la \textbf{prosodie} : en espagnol, l'élément mis en relief porte l'accent d'intensité (ex: \esp{Es mi aBUEla quien...}), ce qui guide la compréhension et la traduction. L'affichage de l'\emph{árbol genealógico} valorise la diversité des structures familiales (monoparentales, recomposées, élargies) présentes dans la classe, ancrant l'apprentissage de la langue dans la réalité sociologique des élèves de Terminale A.

\newpage

\coursec{Méthodes et exercices résolus}

\coursec{Leçon E04 : Traduction de « C'est… que », « C'est… qui », « Ne… que » et de l'emphase ; redacción : describir a su familia (80 palabras)}

\courssub{Observation : découvrir l'emphase}

\begin{methode}[Repérer une structure d'emphase en français et en espagnol]
\textbf{Objectif :} Identifier l'élément mis en relief (sujet, COD, circonstanciel) et la structure utilisée.
\begin{enumerate}
    \item \textbf{Repérer le découpage :} En français, chercher \cle{C'est... qui} (sujet) ou \cle{C'est... que} (autres fonctions) ; en espagnol, chercher \cle{Es... quien/el que} (sujet) ou \cle{Es... lo que / Fue... cuando / Es... donde} (autres).
    \item \textbf{Isoler l'élément mis en relief :} C'est le mot/groupe entre \emph{c'est} et \emph{qui/que} (fr) ou entre \emph{es/fue} et le relatif (es).
    \item \textbf{Identifier la fonction grammaticale :} Sujet $\rightarrow$ \esp{quien/el que} ; COD/Attribut $\rightarrow$ \esp{lo que/el que} ; Temps $\rightarrow$ \esp{cuando} ; Lieu $\rightarrow$ \esp{donde} ; Cause $\rightarrow$ \esp{por qué}.
    \item \textbf{Vérifier l'accord :} En espagnol, le verbe \emph{ser} s'accorde avec l'antécédent (\esp{Es mi padre}, \esp{Son mis hermanos}) ; le relatif \emph{quien} ne varie pas en genre mais au pluriel (\esp{quienes}), \emph{el que} s'accorde (\esp{la que, los que}).
\end{enumerate}
\end{methode}

\begin{popfigure}
\begin{tikzpicture}[mindmap, grow cyclic, every node/.style=concept, concept color=popblue!60,
    level 1/.append style={level distance=4.5cm, sibling angle=90, concept color=popblue!40},
    level 2/.append style={level distance=3cm, sibling angle=45, concept color=poporange!40},
    text=popdark, font=\small]
\node[align=center]{Emphase\\Espagnole}
    child { node {Sujet} child { node {Es/Son + quien(es)} } child { node {Es/Son + el/la/los/las que} } }
    child { node {COD / Attr.} child { node {Es/Fue + lo que} } child { node {Es/Fue + el/la/los/las que} } }
    child { node {Circonstanciels} child { node {Temps : Fue/Era + cuando} } child { node {Lieu : Es/Fue + donde} } child { node {Cause/Manière : Es/Fue + por lo que/como} } }
    child { node {Préposition} child { node {Doublement : Es + PREP + SN + PREP + quien/que} } };
\end{tikzpicture}
\legende{Schéma récapitulatif des structures d'emphase (oraciones escindidas) en espagnol.}
\end{popfigure}

\begin{exoresolu}[Analyse d'un dialogue authentique : « La sorpresa de Mamadou »]
\textbf{Énoncé :} Lisez le dialogue ci-dessous. Pour les phrases 1 à 4 soulignées : a) Identifiez l'élément mis en relief. b) Précisez sa fonction. c) Donnez la structure espagnole correspondante.
\tcblower
\begin{textebox}[Dialogue original : Mamadou et Claudia à Yaoundé]
— \textbf{¿Es tu hermano mayor quien llegó anoche?} \\
— No, \textbf{fue mi prima Claudia quien vino}. Ella \textbf{es la que trajo los regalos}. \\
— \textbf{¿Fue ayer cuando llegaron?} \\
— Sí, y \textbf{fue en el aeropuerto donde nos encontramos}.
\end{textebox}
\textbf{Traduction :} — « C'est ton grand frère qui est arrivé hier soir ? » — « Non, c'est ma cousine Claudia qui est venue. C'est elle qui a apporté les cadeaux. » — « C'est hier qu'ils sont arrivés ? » — « Oui, et c'est à l'aéroport qu'on s'est retrouvés. »

\textbf{Solution détaillée :}
\begin{enumerate}
    \item \esp{¿Es tu hermano mayor quien llegó anoche?} \\
    \textbf{Élément mis en relief :} \emph{tu hermano mayor}. \textbf{Fonction :} Sujet. \textbf{Structure ES :} \cle{Es + SN + quien/el que + verbe}. Accord : \emph{Es} (3e pers. sg) car \emph{hermano} est singulier. \emph{Quien} invariable en genre, variable en nombre (\emph{quienes}).
    \item \esp{Fue mi prima Claudia quien vino.} \\
    \textbf{Élément :} \emph{mi prima Claudia}. \textbf{Fonction :} Sujet. \textbf{Structure :} \cle{Fue + SN + quien/la que + verbe}. Passé (prétérit) car action ponctuelle terminée. \emph{Quien} ou \emph{la que} (accord féminin singulier).
    \item \esp{Ella es la que trajo los regalos.} \\
    \textbf{Élément :} \emph{la} (antécédent = \emph{Ella}). \textbf{Fonction :} Sujet. \textbf{Structure :} \cle{Es + el/la/los/las que + verbe}. \emph{La que} accordé avec \emph{Ella} (féminin singulier). \emph{Trajo} (prétérit indéfini).
    \item \esp{¿Fue ayer cuando llegaron?} \\
    \textbf{Élément :} \emph{ayer}. \textbf{Fonction :} Complément circonstanciel de temps (CC Temps). \textbf{Structure :} \cle{Fue + adv temps + cuando + verbe}. \emph{CUANDO} invariable. Verbe \emph{llegaron} (3e pl) accordé avec \emph{ellos} (sujet réel).
\end{enumerate}
\textbf{Conclusion :} L'espagnol utilise \emph{ser} + élément mis en relief + relatif accordé (sujet) ou relatif invariable/spécifique (circonstanciels). Le temps de \emph{ser} calque le temps français (présent/prétérit/imparfait).
\end{exoresolu}

\courssub{Traduire « C'est… qui » (Emphase sur le sujet)}

\begin{methode}[Règle de traduction : Sujet mis en relief]
\textbf{Structure française :} \cle{C'est / Ce sont + [Sujet mis en relief] + qui + verbe}.
\textbf{Équivalents espagnols (ordre croissant de formalité) :}
\begin{enumerate}
    \item \cle{Es / Son + [Sujet] + quien(es)} : Neutre, très fréquent à l'oral et à l'écrit. \emph{Quien} ne s'accorde qu'en nombre (\emph{quienes}).
    \item \cle{Es / Son + [Sujet] + el/la/los/las que} : Plus formel, précis (accord genre/nombre). Obligatoire si antécédent est chose/idée (\emph{el que, la que}).
    \item \cle{Se trata de + [Sujet]} : Impersonnel, pour présenter/identifier (« Il s'agit de... »).
\end{enumerate}
\textbf{Points de vigilance :}
\begin{itemize}
    \item \textbf{Accord de \emph{ser} :} Singulier (\emph{Es}) ou Pluriel (\emph{Son}) selon le sujet mis en relief.
    \item \textbf{Préposition :} Si le verbe régit une préposition, elle se place \textbf{avant} le relatif : \esp{Es mi amigo \textbf{con} quien hablé} (C'est mon ami \textbf{avec} qui j'ai parlé). \textbf{Jamais} \esp{quien hablé con}.
    \item \textbf{Négation :} \esp{No es Juan quien lo hizo} / \esp{No son ellos quienes mandan}.
    \item \textbf{Accord du verbe après \emph{quien/el que} :} Le verbe de la relative s'accorde avec l'antécédent (3e personne). Ex: \esp{Eres tú quien \textbf{tiene} la razón} (et non *tienes).
\end{itemize}
\end{methode}

\begin{exoresolu}[Thème : Traduire les phrases suivantes en espagnol]
\textbf{Énoncé :} Traduisez en utilisant la structure d'emphase sur le sujet la plus naturelle.
\begin{enumerate}
    \item C'est mon père qui paie les études.
    \item Ce sont mes sœurs qui ont préparé le ndolé.
    \item C'est toi qui as raison.
    \item Ce n'est pas moi qui ai cassé la vitre.
    \item C'est avec ma mère que je vais au marché.
\end{enumerate}
\tcblower
\textbf{Solution détaillée et justifiée :}
\begin{enumerate}
    \item \esp{Es mi padre quien paga los estudios.} (Ou \esp{el que paga}). \emph{Es} (sg), \emph{quien} (personne). Présent d'habitude.
    \item \esp{Son mis hermanas quienes prepararon el ndolé.} \emph{Son} (pl), \emph{quienes} (pl). \emph{Prepararon} (prétérit : action passée achevée). \emph{Ndolé} (plat camerounais, invariable).
    \item \esp{Eres tú quien tiene la razón.} \textbf{Piège :} \emph{C'est toi} $\rightarrow$ \emph{ser} 2e pers. sg \emph{Eres}. Le relatif \emph{quien} impose la 3e personne pour le verbe suivant (\emph{tiene}), bien que le sens soit « tu as ». \emph{Tener razón} (avoir raison).
    \item \esp{No soy yo quien rompió la ventana.} \emph{No soy yo} (C'est pas moi). \emph{Quien rompió} (3e sg). \emph{Rompió} (prétérit).
    \item \esp{Es \textbf{con} mi madre \textbf{con} quien voy al mercado.} \textbf{Règle de la préposition :} \emph{Voy \textbf{con} mi madre} $\rightarrow$ Emphase : \emph{Es \textbf{con} mi madre \textbf{con} quien voy...}. La préposition se double (forme la plus correcte). On évite \esp{quien voy con}.
\end{enumerate}
\end{exoresolu}

\courssub{Traduire « C'est… que » (Emphase sur COD, attribut, circonstanciels)}

\begin{methode}[Règle de traduction : Compléments mis en relief]
\textbf{Structure française :} \cle{C'est / C'était + [Élément] + que + proposition}.
\textbf{Choix du relatif espagnol selon la nature de l'élément :}
\begin{enumerate}
    \item \textbf{COD / Attribut / Idée vague (chose) :} \cle{Es / Fue + [Élément] + lo que + verbe}. \emph{Lo que} = « ce qui / ce que », neutre, invariable.
    \item \textbf{COD / Attribut précis (personne/chose déterminée) :} \cle{Es / Fue + [Élément] + el/la/los/las que + verbe}. Accord en genre/nombre avec l'antécédent.
    \item \textbf{Circonstanciel de TEMPS :} \cle{Fue / Era + [Moment] + cuando + verbe}. \emph{CUANDO} invariable. \emph{Fue} (ponctuel), \emph{Era} (habitude/description passée).
    \item \textbf{Circonstanciel de LIEU :} \cle{Es / Fue + [Lieu] + donde + verbe}. \emph{DONDE} invariable.
    \item \textbf{Circonstanciel de MANIÈRE / CAUSE :} \cle{Es / Fue + [Manière/Cause] + como / por lo que / por qué + verbe}.
    \item \textbf{Préposition (à, pour, avec, sans...) :} La préposition \textbf{précède} le relatif (doublement recommandé) : \esp{Es \textbf{para} ti \textbf{para} lo que trabajo} (C'est pour toi que je travaille).
\end{enumerate}
\end{methode}

\begin{exoresolu}[Transformation et Version : Emphase sur compléments]
\textbf{Énoncé A (Thème) :} Mettez en relief l'élément entre parenthèses en espagnol.
\begin{enumerate}
    \item Juan compró el coche \textbf{(ayer)}.
    \item María escribió la carta \textbf{(con su pluma nueva)}.
    \item Los niños comen \textbf{(fruta)}.
    \item Nos vimos \textbf{(en la estación)}.
\end{enumerate}
\textbf{Énoncé B (Version) :} Traduisez en français.
\begin{enumerate}
    \item \esp{Fue el libro lo que me dio.}
    \item \esp{Es en Madrid donde vive mi tío.}
    \item \esp{Fue por dinero por lo que lo hizo.}
\end{enumerate}
\tcblower
\textbf{Solution A (Thème) :}
\begin{enumerate}
    \item \esp{Fue \textbf{ayer cuando} Juan compró el coche.} (Temps $\rightarrow$ \emph{cuando}). \emph{Fue} (action ponctuelle passée).
    \item \esp{Es \textbf{con su pluma nueva con la que} María escribió la carta.} (Instrument $\rightarrow$ \emph{con la que} car antécédent féminin singulier déterminé). Doublement de la préposition \emph{con}.
    \item \esp{Es \textbf{fruta lo que} comen los niños.} (COD chose indéterminée $\rightarrow$ \emph{lo que}). \emph{Comen} (présent).
    \item \esp{Fue \textbf{en la estación donde} nos vimos.} (Lieu $\rightarrow$ \emph{donde}). \emph{Fue} (passé ponctuel). \emph{Nos vimos} (prétérit).
\end{enumerate}
\textbf{Solution B (Version) :}
\begin{enumerate}
    \item \esp{Fue el libro lo que me dio.} $\rightarrow$ « C'est le livre qu'il m'a donné » (ou « Ce qu'il m'a donné, c'est le livre »). \emph{Lo que} = COD chose.
    \item \esp{Es en Madrid donde vive mi tío.} $\rightarrow$ « C'est à Madrid que vit mon oncle ». Lieu $\rightarrow$ \emph{donde}.
    \item \esp{Fue por dinero por lo que lo hizo.} $\rightarrow$ « C'est pour l'argent qu'il l'a fait ». Cause $\rightarrow$ \emph{por lo que} (ou \emph{por qué}). Doublement \emph{por}.
\end{enumerate}
\end{exoresolu}

\courssub{Traduire « Ne… que » (Restriction)}

\begin{methode}[Les 4 façons de dire « Ne... que » en espagnol]
Le français \cle{Ne... que} = « seulement / rien d'autre que ». En espagnol, le choix dépend de la catégorie grammaticale du mot restreint et du registre.
\begin{enumerate}
    \item \cle{Solo / Solamente / Únicamente + [Mot restreint]} : \textbf{Adverbes} placés \textbf{devant} l'élément restreint. Le verbe est \textbf{affirmatif}. \emph{Uniquement} plus formel.
    \item \cle{No + verbe + más que + [Mot restreint]} : \textbf{Locution conjonctive}. Verbe \textbf{négatif}. « Rien d'autre que / Que seulement ». Très courant.
    \item \cle{No + verbe + sino (+ que) + [Mot restreint]} : \textbf{Adversatif}. « Non pas X, \textbf{mais} Y ». \emph{Sino} (seul) si même structure grammaticale (ex: nom vs nom). \emph{Sino que} si proposition verbe conjuguée suit. Verbe \textbf{négatif}.
    \item \cle{No + verbe + otro(s) que + [Mot restreint]} : « Nul autre que... ». Littéraire/formel.
\end{enumerate}
\textbf{Règle d'or :} \emph{Solo} (sans accent d'après RAE 2010, mais \emph{sólo} toléré pour éviter ambiguïté avec adj \emph{solo} « seul ») se place \textbf{juste avant} le mot qu'il restreint. \emph{No... más que} entoure le verbe.
\end{methode}

\begin{exoresolu}[Exercices variés : Restriction]
\textbf{Énoncé 1 :} Traduisez en espagnol (donnez les 2 formes principales si possibles).
\begin{enumerate}
    \item Je n'ai \textbf{qu'un} frère.
    \item Il ne mange \textbf{que} du riz.
    \item Ce n'est \textbf{pas moi}, c'est \textbf{toi}.
    \item Nous ne faisons \textbf{qu'étudier}.
\end{enumerate}
\textbf{Énoncé 2 :} Choisissez la bonne forme entre parenthèses et justifiez.
\begin{enumerate}
    \item No tengo \_\_\_\_\_\_\_ 500 francos. (solo / no... más que)
    \item No vino \_\_\_\_\_\_\_ Juan, \_\_\_\_\_\_\_ Pedro. (sino / más que)
    \item \_\_\_\_\_\_\_ quiero agua. (Solo / No... más que)
\end{enumerate}
\tcblower
\textbf{Solution 1 :}
\begin{enumerate}
    \item \esp{Tengo \textbf{solo un} hermano.} / \esp{No tengo \textbf{más que un} hermano.} (Nom comptable).
    \item \esp{Come \textbf{solo} arroz.} (Restriction sur l'objet) / \esp{No come \textbf{más que} arroz.} \textbf{Attention :} \esp{Solo come arroz} est ambigu (il mange seul / il ne fait que manger du riz). Placer \emph{solo} devant l'élément restreint (\emph{arroz}) lève l'ambiguïté.
    \item \esp{No soy \textbf{yo}, \textbf{sino} tú.} (Identité, même fonction : sujet. \emph{Sino} sans \emph{que} car pas de verbe conjugué après).
    \item \esp{Solo estudiamos.} (Nous ne faisons qu'étudier = action unique) / \esp{No hacemos \textbf{más que} estudiar.} (Restriction sur l'activité).
\end{enumerate}
\textbf{Solution 2 :}
\begin{enumerate}
    \item \esp{No tengo \textbf{más que} 500 francos.} (Ou \esp{Tengo \textbf{solo} 500 francos}). \emph{No... más que} insiste sur la quantité max. \emph{Solo} possible mais ambigu (seul/only).
    \item \esp{No vino Juan, \textbf{sino} Pedro.} Structure \emph{No X, sino Y}. Opposition entre deux sujets.
    \item \esp{\textbf{Solo} quiero agua.} (Verbe affirmatif + \emph{solo} devant verbe = « Je veux seulement de l'eau »). \esp{No quiero \textbf{más que} agua} possible mais plus lourd.
\end{enumerate}
\textbf{Rappel :} \esp{No solo... sino (también)...} = « Non seulement... mais aussi... ». Ex: \esp{No solo estudia, \textbf{sino que} \textbf{también} trabaja.}
\end{exoresolu}

\courssub{Lexique : La famille et le portrait (Physique et Moral)}

\begin{methode}[Construire son répertoire lexical pour la description]
\textbf{Étape 1 : L'arbre généalogique (Parentesco).} Maîtriser les paires M/F et les composés.
\begin{center}
\begin{tabularx}{\linewidth}{|X|X|X|}\hline
\rowcolor{popblueL} \textbf{Français} & \textbf{Español (M)} & \textbf{Español (F)} \\ \hline
Père / Mère & Padre & Madre \\
Frère / Sœur & Hermano & Hermana \\
Fils / Fille & Hijo & Hija \\
Grand-père / Mère & Abuelo & Abuela \\
Oncle / Tante & Tío & Tía \\
Neveu / Nièce & Sobrino & Sobrina \\
Cousin / Cousine & Primo & Prima \\
Beau-père / Belle-mère & Suegro & Suegra \\
Beau-frère / Belle-sœur & Cuñado & Cuñada \\
Demi-frère / sœur & Hermanastro & Hermanastra \\ \hline
\end{tabularx}
\end{center}
\textbf{Étape 2 : Verbes d'état : \cle{SER} vs \cle{ESTAR}.}
\begin{itemize}
    \item \textbf{SER} : Identité, origine, traits physiques \textbf{permanents}, caractère, profession, possession, heure. \esp{Es alto, es simpático, es médico, es de Yaundé.}
    \item \textbf{ESTAR} : État physique/émotionnel \textbf{passager}, localisation, résultat d'action. \esp{Está enfermo, está contento, está en Douala, está cansado.}
\end{itemize}
\textbf{Étape 3 : Adjectifs clés (Accord genre/nombre obligatoire).}
\begin{itemize}
    \item \textit{Physique :} \esp{alto/bajo, delgado/gordo, joven/viejo, guapo/fea, moreno/rubio, ojos grandes/pequeños, pelo liso/rizado.}
    \item \textit{Moral :} \esp{simpático/antipático, amable/gruñón, trabajador/vago, generoso/tacaño, responsable/irresponsable, callado/hablador, inteligente/tonto.}
\end{itemize}
\textbf{Étape 4 : Connecteurs pour 80 mots.} \esp{Además, / También, / Sin embargo, / Por eso, / En resumen.}
\end{methode}

\begin{exoresolu}[Complétion et Portrait croisé]
\textbf{Énoncé A :} Complétez l'arbre généalogique de Carlos avec le bon terme de parenté.
\esp{Carlos tiene un hijo, Pedro. Pedro tiene una hermana, Ana. Para Carlos, Ana es su \_\_\_\_\_\_. El padre de Carlos es Don José. Para Pedro, Don José es su \_\_\_\_\_\_. La mujer de Don José es Doña María. Para Ana, Doña María es su \_\_\_\_\_\_. El hermano de la madre de Pedro es Roberto. Para Pedro, Roberto es su \_\_\_\_\_\_.}

\textbf{Énoncé B :} Corrigez les erreurs d'accord ou de choix verbal (Ser/Estar) dans ce mini-portrait.
\esp{Mi hermano es alto y *delgada*. Él *está* muy simpático y *está* trabajador. Hoy *es* cansado porque *es* enfermo. Pero normalmente *es* alegre.}
\tcblower
\textbf{Solution A :}
\begin{enumerate}
    \item \esp{Hija} (fille de Carlos).
    \item \esp{Abuelo} (père du père = grand-père).
    \item \esp{Abuela} (femme du grand-père = grand-mère).
    \item \esp{Tío} (frère de la mère = oncle).
\end{enumerate}
\textbf{Solution B (Corrections) :}
\esp{Mi hermano \textbf{es} alto y \textbf{delgado} (accord masc. sg). Él \textbf{es} muy simpático (caractère = SER) y \textbf{es} trabajador (caractère = SER). Hoy \textbf{está} cansado (état passager = ESTAR) porque \textbf{está} enfermo (état santé = ESTAR). Pero normalmente \textbf{es} alegre (caractère = SER).}
\textbf{Justifications :} \emph{Delgado} accord masculin. \emph{Simpático/Trabajador/Alegre} = traits de personnalité $\rightarrow$ \textbf{SER}. \emph{Cansado/Enfermo} = états temporaires $\rightarrow$ \textbf{ESTAR}.
\end{exoresolu}

\courssub{Rédaction guidée : Describir a su familia (80 palabras)}

\begin{methode}[Rédiger une description de famille en 80 mots (±10\%)]
\textbf{Consigne type Bac :} « Descrivez votre famille (membres, physique, caractère, activités) en 80 mots environ. »
\begin{enumerate}
    \item \textbf{Planifier (5 min) :} Structure en 4 paragraphes courts.
        \begin{itemize}
            \item P1 : Composition (Nombre, qui ?) + Verbe \emph{ser} / \emph{vivir}. \emph{« Mi familia es grande. Somos cinco: mis padres, mis dos hermanas y yo. Vivimos en Yaundé. »} (\approx 25 mots).
            \item P2 : Portrait Parents (Physique + Moral). \emph{« Mi padre es alto y moreno, muy trabajador y responsable. Mi madre es baja y delgada, amable y paciente. »} (\approx 25 mots).
            \item P3 : Portrait Frères/Sœurs + Soi (Physique + Moral + Activité). \emph{« Mis hermanas son jóvenes y estudiantes. Una es rubia y habladora; la otra es morena y callada. Yo soy deportista y estudio en el liceo. »} (\approx 30 mots).
            \item P4 : Vie commune / Conclusion affective. \emph{« Nos queremos mucho. Los domingos comemos ndolé juntos. Es una familia unida y feliz. »} (\approx 20 mots).
        \end{itemize}
    \item \textbf{Rédiger (10 min) :} Phrases courtes, connecteurs (\emph{además, pero, por tanto}). Vérifier \textbf{Ser/Estar}, \textbf{Accords}, \textbf{Vocabulaire précis}.
    \item \textbf{Compter \& Corriger (5 min) :} Compter les mots (articles, prépositions, contractions = 1 mot). Ajuster (supprimer adjectifs ou ajouter détails). Relire : Accents (\emph{papá, mamá, quién, cómo}), \emph{ñ}, \emph{ll}, \emph{rr}.
\end{enumerate}
\textbf{Astuce :} 80 mots $\approx$ 8-10 phrases de 8-10 mots. Viser 75-85 mots.
\end{methode}

\begin{exoresolu}[Modèle de rédaction commentée (82 mots)]
\textbf{Énoncé :} Produisez une rédaction de 80 mots environ décrivant votre famille. Voici un modèle corrigé et analysé.
\tcblower
\begin{textebox}[Modèle : Mi familia]
Mi familia es numerosa y unida. Somos seis: mis padres, mis tres hermanos y yo. Vivimos en Yaundé. Mi padre, Pedro, es alto, moreno y muy trabajador; es médico. Mi madre, Ana, es baja, delgada y muy amable; está siempre ocupada en casa. Mis hermanos mayores son gemelos: son rubios, deportistas y ruidosos. Mi hermana pequeña es callada, estudiosa y tiene los ojos grandes. Yo soy el mediano, me gusta el fútbol y estudio en Terminale. Los fines de semana preparamos ndolé juntos y reímos mucho. Somos muy felices.
\end{textebox}
\textbf{Analyse lexicale et grammaticale :}
\begin{itemize}
    \item \textbf{Mots :} 82 (dans la fourchette 75-85). Parfait.
    \item \textbf{Vocabulaire famille :} \emph{padres, hermanos, gemelos, mediano, mayor, pequeña}. Précis.
    \item \textbf{Ser/Estar :} \emph{es alto/moreno/trabajador/médico} (identité/profession) ; \emph{está siempre ocupada} (état continu/résultat = \emph{estar}) ; \emph{son rubios/deportistas} (caractère/physique) ; \emph{es callada/estudiosa} ; \emph{soy... me gusta... estudio}.
    \item \textbf{Accords :} \emph{numerosa/unida} (f sg familia) ; \emph{altO, morenO, trabajadOR} (padre) ; \emph{bajA, delgadA, amable} (madre) ; \emph{rubiOS, deportistAS, ruidosOS} (gemelos masc pl) ; \emph{calladA, estudiosA} (hermana f sg) ; \emph{medianO, felices} (nosotros).
    \item \textbf{Contexte camerounais :} \emph{Yaundé, ndolé, Terminale}. Authentique.
    \item \textbf{Connecteurs :} \emph{y, ;, Los fines de semana, juntos}.
\end{itemize}
\textbf{Conclusion :} Ce texte obtient la note maximale : respect de la contrainte de longueur, vocabulaire riche et varié, grammaire irréprochable (Ser/Estar, accords), ancrage culturel.
\end{exoresolu}

\begin{attention}[Les 5 erreurs qui coûtent des points au Bac]
\begin{enumerate}
    \item \textbf{Accents sur les relatifs dans l'emphase :} \textbf{JAMAIS d'accent} sur \emph{quien, donde, cuando, como, el que} dans une phrase clivée (relative avec antécédent). Ex: \esp{Es Juan \textbf{quien} vino} (pas \emph{quién}), \esp{Es \textbf{donde} vivo} (pas \emph{dónde}), \esp{Es \textbf{como} lo hice} (pas \emph{cómo}). L'accent diacritique (\emph{quién, dónde, cómo}) est réservé aux interrogatifs/exclamatifs directs ou indirects : \esp{¿\textbf{Quién} vino?}, \esp{No sé \textbf{dónde} vive}.
    \item \textbf{Confusion \emph{Ser} / \emph{Estar} dans le portrait :} \esp{*Mi padre está alto} (Faux : taille = \emph{es}). \esp{*Mi madre es enferma} (Faux : état = \emph{está}, sauf maladie chronique identitaire rare). \esp{*Estoy aburrido} (Ennuyé) vs \esp{Soy aburrido} (Ennuyeux). \textbf{Réflexe :} Permanent/Identité $\rightarrow$ SER ; Passager/Lieu/État $\rightarrow$ ESTAR.
    \item \textbf{Préposition « orpheline » dans l'emphase :} \esp{*Es mi amigo quien hablé \textbf{con}.} (Calque français). \textbf{Correct :} \esp{Es \textbf{con} mi amigo \textbf{con} quien hablé} ou \esp{Es mi amigo \textbf{con} quien hablé}. La préposition \textbf{ne reste jamais à la fin} de la proposition relative en espagnol standard.
    \item \textbf{\emph{Solo} vs \emph{Sólo} et place de \emph{solo} :} \esp{Solo como arroz} (Ambigu : Je mange \textbf{seulement} du riz / Je mange du riz \textbf{seul} ?). Mieux : \esp{Como \textbf{solo} arroz} (restriction sur riz) ou \esp{Como arroz \textbf{solamente}}. RAE recommande \emph{solo} sans accent pour les deux, mais l'accent \emph{sólo} (adverbe) reste toléré pour lever l'ambiguïté. Au Bac : placez \emph{solo/solamente} \textbf{juste devant} le mot restreint.
    \item \textbf{Non-respect de la longueur (80 mots) :} « Environ 80 mots » = \textbf{70 à 90 mots}. En-dessous de 70 : pénalité « insuffisance ». Au-dessus de 90 : pénalité « hors sujet / non-respect consigne ». \textbf{Stratégie :} Compter \textbf{avant} de recopier propre. Les contractions (\emph{del, al}) comptent pour 1 mot. Les chiffres (\emph{2025}) comptent pour 1 mot.
\end{enumerate}
\end{attention}

\newpage

\coursec{Exercices}

\courssub{Je vérifie mes connaissances}

\begin{exercice}[QCM : la bonne tournure d'emphase]
Pour chaque phrase française, choisis la traduction espagnole correcte (A, B ou C).
\begin{enumerate}
\item « C'est mon frère qui parle espagnol. » \\
A. \esp{Es mi hermano que habla español.} \quad B. \esp{Es mi hermano quien habla español.} \quad C. \esp{Mi hermano es que habla español.}
\item « C'est hier qu'elle est arrivée. » \\
A. \esp{Fue ayer cuando llegó.} \quad B. \esp{Fue ayer que llegó.} \quad C. \esp{Es ayer cuando llegó.}
\item « Je n'ai qu'une sœur. » \\
A. \esp{No tengo que una hermana.} \quad B. \esp{Solo tengo una hermana.} \quad C. \esp{No tengo sino una hermana que vive.}
\item « Ce que j'aime, c'est le football. » \\
A. \esp{Que me gusta es el fútbol.} \quad B. \esp{Lo que me gusta es el fútbol.} \quad C. \esp{Es que me gusta el fútbol.}
\end{enumerate}
\end{exercice}

\begin{exercice}[Vrai ou faux ? Justifie ta réponse]
Dis si les affirmations suivantes sont vraies ou fausses, puis justifie en une phrase.
\begin{enumerate}
\item En espagnol, « c'est\ldots qui » se traduit toujours par \esp{es\ldots que}.
\item Après \esp{no\ldots más que}, on peut placer un nom ou un verbe à l'infinitif.
\item \esp{Fue en 2010 cuando llegamos a Yaoundé} traduit correctement « C'est en 2010 que nous sommes arrivés à Yaoundé ».
\item \esp{No tengo sino un hermano} est la seule façon correcte de traduire « Je n'ai qu'un frère ».
\end{enumerate}
\end{exercice}

\begin{exercice}[Vocabulaire : la famille]
Complète chaque définition par le mot espagnol de la parenté qui convient.
\begin{enumerate}
\item El hermano de mi madre es mi \ldots\ldots\ldots
\item La hija de mi tío es mi \ldots\ldots\ldots
\item El padre de mi padre es mi \ldots\ldots\ldots
\item La mujer de mi hermano es mi \ldots\ldots\ldots
\item Mis tíos y mis primos forman mi familia \ldots\ldots\ldots (le contraire de \esp{familia nuclear}).
\end{enumerate}
\end{exercice}

\begin{exercice}[Conjugaison : ser au service de l'emphase]
Complète les phrases avec le verbe \esp{ser} conjugué au temps indiqué entre parenthèses.
\begin{enumerate}
\item \esp{\ldots\ldots\ldots mi abuela quien me enseñó a cocinar.} (pretérito indefinido)
\item \esp{\ldots\ldots\ldots en Douala donde nacieron mis padres.} (pretérito indefinido)
\item \esp{\ldots\ldots\ldots mis hermanos quienes cuidan el campo de cacao.} (presente)
\item \esp{\ldots\ldots\ldots mañana cuando sabremos los resultados del examen.} (futuro)
\end{enumerate}
\end{exercice}

\courssub{Je m'entraîne}

\begin{exercice}[Thème : phrases à traduire]
Traduis en espagnol les phrases suivantes en utilisant une tournure d'emphase ou de restriction.
\begin{enumerate}
\item C'est ma grand-mère qui m'a appris à cuisiner.
\item C'est en 2010 que nous sommes arrivés à Yaoundé.
\item Je n'ai qu'un seul rêve : réussir au Bac.
\item Ce ne sont pas mes frères qui ont téléphoné, c'est mon oncle.
\item C'est au marché de Mokolo que ma mère achète les légumes.
\end{enumerate}
\end{exercice}

\begin{exercice}[Texte à trous : l'emphase en contexte]
Complète le texte suivant avec les mots de la liste : \esp{quien}, \esp{que}, \esp{donde}, \esp{cuando}, \esp{lo que}, \esp{sino}, \esp{más que}. Chaque mot ne sert qu'une fois.

\esp{Fue mi padre \ldots\ldots\ldots construyó nuestra casa de Bafoussam. Fue en 1998 \ldots\ldots\ldots la terminó. Es en el salón \ldots\ldots\ldots nos reunimos toda la familia los domingos. \ldots\ldots\ldots más me gusta de mi casa es el patio. No tenemos \ldots\ldots\ldots un pequeño jardín, pero nos basta. No fue mi madre \ldots\ldots\ldots plantó los árboles, \ldots\ldots\ldots mi abuelo.}
\end{exercice}

\begin{exercice}[Recherche d'erreurs]
Chaque phrase contient une erreur de traduction des tournures étudiées. Souligne l'erreur et réécris la phrase correctement.
\begin{enumerate}
\item \esp{Es mi tío que vive en Bafoussam.}
\item \esp{No tengo que un hermano.}
\item \esp{Fue ayer que llegó mi prima de España.}
\item \esp{Que quiero es descansar después del examen.}
\item \esp{No vinieron mis primos, sino que mis vecinos.}
\end{enumerate}
\end{exercice}

\begin{exercice}[Appariements : français--espagnol]
Relie chaque phrase française (1 à 5) à sa traduction espagnole (A à E).
\begin{enumerate}
\item C'est toi que je cherche.
\item Nous n'avons que deux jours de vacances.
\item C'est à l'école que j'ai appris l'espagnol.
\item Ce n'est pas la télévision que je regarde, c'est un film.
\item Ce que dit mon père est toujours juste.
\end{enumerate}
\begin{itemize}
\item[A.] \esp{Lo que dice mi padre siempre es justo.}
\item[B.] \esp{No veo la televisión, sino una película.}
\item[C.] \esp{Es a ti a quien busco.}
\item[D.] \esp{Fue en la escuela donde aprendí español.}
\item[E.] \esp{No tenemos más que dos días de vacaciones.}
\end{itemize}
\end{exercice}

\courssub{Je me prépare au Bac}

\begin{exercice}[Compréhension et questions de langue]
Lis le texte suivant, puis réponds aux questions.

\begin{textebox}[Una familia unida]
Fue mi abuela quien reunió a toda la familia después de la muerte de mi abuelo. Ella no tenía más que una pequeña casa en un barrio de Yaoundé, pero allí cabíamos todos: mis padres, mis tres hermanos, mis tíos y mis primos. Es en su cocina donde aprendí los platos típicos de nuestra región, y fue ella quien me enseñó que la familia es lo más importante. No somos ricos, sino trabajadores. Lo que más admiro de mi abuela es su fuerza: nunca se queja y siempre tiene una palabra amable para cada uno. Hoy, es ella quien decide los grandes asuntos de la familia, y todos la respetamos.
\end{textebox}

\textbf{Questions de compréhension} (réponds en espagnol, avec des phrases complètes) :
\begin{enumerate}
\item ¿Quién reunió a la familia después de la muerte del abuelo?
\item ¿Dónde aprendió el narrador los platos típicos de su región?
\item ¿Qué es lo que más admira el narrador de su abuela?
\item Según el texto, ¿la familia es rica? Justifica tu respuesta con una frase del texto.
\end{enumerate}

\textbf{Questions de langue} :
\begin{enumerate}
\item Relève dans le texte deux tournures clivées avec \esp{ser} et traduis-les en français.
\item Relève dans le texte deux expressions de restriction et donne pour chacune un équivalent français.
\item Réécris la phrase \esp{Mi abuela reunió a toda la familia} en la transformant en tournure d'emphase sur le sujet.
\end{enumerate}
\end{exercice}

\begin{exercice}[Thème et version]
\textbf{Thème.} Traduis en espagnol le texte suivant :

« C'est mon père qui travaille le plus dans notre famille. C'est à cinq heures du matin qu'il part au champ de cacao. Nous n'avons qu'une seule moto pour toute la famille, mais elle nous suffit. Ce que je veux, c'est réussir au baccalauréat pour rendre mes parents fiers. Ce n'est pas l'argent qui fait le bonheur, c'est l'union de la famille. »

\textbf{Version.} Traduis en français le texte suivant :

\esp{Fue mi padre quien construyó esta casa. No tenemos más que un campo, pero nos basta. Es en verano cuando toda la familia se reúne en el pueblo, y lo que más me gusta es escuchar las historias de mis abuelos. No vinieron mis primos a la fiesta, sino mis tíos de Bata.}
\end{exercice}

\begin{exercice}[Rédaction type Bac : Describa a su familia]
\textbf{Sujet.} \esp{Describa a su familia: presente a los miembros, describa el aspecto físico y el carácter de dos de ellos y diga qué es lo que más aprecia de su familia.} (80 palabras)

\textbf{Consignes obligatoires :}
\begin{enumerate}
\item Rédige un texte suivi en espagnol, d'environ 80 mots (entre 75 et 90 mots).
\item Emploie au moins \textbf{une tournure d'emphase} (\esp{es\ldots quien}, \esp{fue\ldots cuando}, \esp{lo que\ldots}) et \textbf{une expression de restriction} (\esp{solo}, \esp{no\ldots más que}, \esp{no\ldots sino}).
\item Utilise au moins six mots du lexique de la parenté et trois adjectifs de description physique ou morale.
\end{enumerate}

\textbf{Grille d'évaluation (sur 10 points) :}
\begin{center}
\begin{tabularx}{\linewidth}{|X|c|}
\hline
\rowcolor{popblueL} \textbf{Critère} & \textbf{Points} \\
\hline
Respect du sujet et du nombre de mots (75--90) & 2 \\
\hline
Lexique de la famille et de la description (richesse, exactitude) & 3 \\
\hline
Emploi correct d'une tournure d'emphase et d'une restriction & 3 \\
\hline
Correction grammaticale et orthographique (accents, accords) & 2 \\
\hline
\end{tabularx}
\end{center}
\end{exercice}

\newpage

\coursec{Corrigés des exercices}

\begin{corrige}[Exercice 1 — QCM : la bonne tournure d'emphase]
\begin{enumerate}
\item \textbf{B.} \esp{¿Es mi hermano quien habla español?} « C'est\ldots qui » portant sur une personne sujet se traduit par \esp{es\ldots quien}. La réponse A est fautive : \esp{que} ne peut pas renvoyer à une personne sujet dans une tournure clivée.
\item \textbf{A.} \esp{¿Fue ayer cuando llegó?} L'emphase sur un complément de temps exige \esp{fue\ldots cuando} (et non \esp{que}). De plus, l'accord de temps impose le passé : \esp{fue}, pas \esp{es}.
\item \textbf{B.} \esp{Solo tengo una hermana.} « Ne\ldots que » se rend par \esp{solo} (ou \esp{no tengo más que una hermana}). La réponse A est un calque du français ; la C ajoute une subordonnée qui change le sens.
\item \textbf{B.} \esp{Lo que me gusta es el fútbol.} « Ce que » (chose, idée) se traduit par \esp{lo que}. La réponse A est impossible : \esp{que} seul ne peut pas ouvrir la phrase avec cette valeur.
\end{enumerate}
\end{corrige}

\begin{corrige}[Exercice 2 — Vrai ou faux ?]
\begin{enumerate}
\item \textbf{Faux.} Quand l'élément mis en relief est une personne sujet, on dit \esp{es\ldots quien} : \esp{Es mi madre quien cocina.} \esp{Es\ldots que} ne s'emploie que pour les choses ou dans des tournures comme \esp{es que\ldots} (« c'est que\ldots »).
\item \textbf{Vrai.} \esp{No\ldots más que} peut être suivi d'un nom (\esp{No tengo más que un hermano}) ou d'un infinitif (\esp{No hago más que estudiar}).
\item \textbf{Vrai.} L'emphase sur un complément de temps passé se construit avec \esp{fue} (pretérito indefinido de \esp{ser}) + le complément + \esp{cuando} : \esp{Fue en 2010 cuando llegamos a Yaoundé.}
\item \textbf{Faux.} Il existe plusieurs traductions correctes : \esp{Solo tengo un hermano}, \esp{Tengo únicamente un hermano}, \esp{No tengo más que un hermano}. La tournure \esp{no\ldots sino} est correcte mais surtout utilisée dans l'opposition (« pas\ldots mais\ldots »).
\end{enumerate}
\end{corrige}

\begin{corrige}[Exercice 3 — Vocabulaire : la famille]
\begin{enumerate}
\item \esp{tío} — le frère de ma mère est mon \emph{oncle}.
\item \esp{prima} — la fille de mon oncle est ma \emph{cousine} (au masculin : \esp{primo}).
\item \esp{abuelo} — le père de mon père est mon \emph{grand-père} (la mère : \esp{abuela}).
\item \esp{cuñada} — la femme de mon frère est ma \emph{belle-sœur} (le mari de ma sœur : \esp{cuñado}).
\item \esp{extensa} — \esp{familia extensa} = famille élargie, par opposition à \esp{familia nuclear} (parents et enfants).
\end{enumerate}
\end{corrige}

\begin{corrige}[Exercice 4 — Conjugaison : ser au service de l'emphase]
\begin{enumerate}
\item \esp{\textbf{Fue} mi abuela quien me enseñó a cocinar.} — pretérito indefinido, 3e personne du singulier (l'élément mis en relief est singulier).
\item \esp{\textbf{Fue} en Douala donde nacieron mis padres.} — \esp{ser} s'accorde avec l'élément mis en relief, ici le complément de lieu, donc singulier : \esp{fue}.
\item \esp{\textbf{Son} mis hermanos quienes cuidan el campo de cacao.} — presente, 3e personne du pluriel car \esp{mis hermanos} est pluriel ; noter l'accord \esp{quienes}.
\item \esp{\textbf{Será} mañana cuando sabremos los resultados del examen.} — futuro, 3e personne du singulier.
\end{enumerate}
\textbf{Rappel :} dans la tournure clivée, \esp{ser} s'accorde en nombre avec l'élément mis en relief : \esp{Es mi padre quien\ldots} / \esp{Son mis padres quienes\ldots}
\end{corrige}

\begin{corrige}[Exercice 5 — Thème : phrases à traduire]
\begin{enumerate}
\item \esp{Es mi abuela quien me enseñó a cocinar.} (ou \esp{Fue mi abuela quien me enseñó a cocinar} si l'on insiste sur le passé.) Personne sujet : \esp{quien}.
\item \esp{Fue en 2010 cuando llegamos a Yaoundé.} Complément de temps : \esp{fue\ldots cuando}.
\item \esp{Solo tengo un sueño: aprobar el Bachillerato.} (ou \esp{No tengo más que un sueño: aprobar el Bachillerato.}) « Réussir à » = \esp{aprobar} + nom ; attention au faux ami : \esp{réussir un examen} ne se dit pas \esp{realizar}.
\item \esp{No son mis hermanos quienes han llamado, sino mi tío.} Opposition « pas\ldots mais\ldots » : \esp{no\ldots sino} ; accord pluriel : \esp{son\ldots quienes}.
\item \esp{Es en el mercado de Mokolo donde mi madre compra las verduras.} Complément de lieu : \esp{es\ldots donde}.
\end{enumerate}
\end{corrige}

\begin{corrige}[Exercice 6 — Texte à trous : l'emphase en contexte]
\esp{Fue mi padre \textbf{quien} construyó nuestra casa de Bafoussam. Fue en 1998 \textbf{cuando} la terminó. Es en el salón \textbf{donde} nos reunimos toda la familia los domingos. \textbf{Lo que} más me gusta de mi casa es el patio. No tenemos \textbf{más que} un pequeño jardín, pero nos basta. No fue mi madre \textbf{quien} plantó los árboles, \textbf{sino} mi abuelo.}

\textbf{Justifications :} \esp{quien} renvoie à une personne sujet (\esp{mi padre}) ; \esp{cuando} suit un complément de temps (\esp{en 1998}) ; \esp{donde} suit un complément de lieu (\esp{en el salón}) ; \esp{lo que} introduit « ce que » (chose) ; \esp{no\ldots más que} traduit « ne\ldots que » ; dans la dernière phrase, l'emphase porte sur le sujet personne (\esp{mi madre}), donc \esp{quien} est la forme correcte (et non \esp{que}), et \esp{sino} traduit « mais » dans l'opposition « pas\ldots mais\ldots ».

\textbf{Traduction :} « C'est mon père qui a construit notre maison de Bafoussam. C'est en 1998 qu'il l'a terminée. C'est dans le salon que toute la famille se réunit le dimanche. Ce que j'aime le plus dans ma maison, c'est la cour. Nous n'avons qu'un petit jardin, mais il nous suffit. Ce n'est pas ma mère qui a planté les arbres, c'est mon grand-père. »
\end{corrige}

\begin{corrige}[Exercice 7 — Recherche d'erreurs]
\begin{enumerate}
\item Erreur : \esp{que}. Correction : \esp{Es mi tío \textbf{quien} vive en Bafoussam.} — personne sujet : \esp{quien}.
\item Erreur : \esp{que}. Correction : \esp{No tengo \textbf{más que} un hermano} (ou \esp{\textbf{Solo} tengo un hermano}). — « ne\ldots que » ne se traduit jamais par \esp{que} seul.
\item Erreur : \esp{que}. Correction : \esp{Fue ayer \textbf{cuando} llegó mi prima de España.} — complément de temps : \esp{cuando}.
\item Erreur : \esp{Que}. Correction : \esp{\textbf{Lo que} quiero es descansar después del examen.} — « ce que » (chose) = \esp{lo que}.
\item Erreur : \esp{sino que}. Correction : \esp{No vinieron mis primos, \textbf{sino} mis vecinos.} — devant un nom (sans verbe conjugué répété), on emploie \esp{sino} seul ; \esp{sino que} n'est requis que devant une proposition avec verbe conjugué : \esp{No vinieron mis primos, sino que vinieron mis vecinos.}
\end{enumerate}
\end{corrige}

\begin{corrige}[Exercice 8 — Appariements : français--espagnol]
\begin{center}
\begin{tabular}{|c|l|}
\hline
\rowcolor{popblueL} \textbf{Phrase} & \textbf{Traduction} \\
\hline
1 & C — \esp{Es a ti a quien busco.} \\
\hline
2 & E — \esp{No tenemos más que dos días de vacaciones.} \\
\hline
3 & D — \esp{Fue en la escuela donde aprendí español.} \\
\hline
4 & B — \esp{No veo la televisión, sino una película.} \\
\hline
5 & A — \esp{Lo que dice mi padre siempre es justo.} \\
\hline
\end{tabular}
\end{center}
\textbf{Remarques :} en 1, la personne mise en relief est complément : \esp{a ti a quien} (préposition \esp{a} reprise). En 3, \esp{fue} car l'apprentissage est passé. En 4, opposition « pas\ldots mais\ldots » : \esp{sino}.
\end{corrige}

\begin{corrige}[Exercice 9 — Compréhension et questions de langue]
\begin{textebox}[Texte support de l'exercice (reconstitué)]
Mi abuela fue el pilar de la familia. \textbf{Fue mi abuela quien reunió a toda la familia} después de la muerte del abuelo. \textbf{Es en su cocina donde aprendí los platos típicos} de mi región; \textbf{fue ella quien me enseñó} el secreto del buen guiso. No tenía más que una pequeña casa, pero \textbf{no éramos pobres, sino trabajadores}. \textbf{Es ella quien decide} el menú de los domingos. Lo que más admiro es su fuerza: nunca se queja y siempre tiene una palabra amable para cada uno. No somos ricos, \textbf{sino} trabajadores.
\end{textebox}

\textbf{Questions de compréhension (réponses modèles) :}
\begin{enumerate}
\item \esp{Fue la abuela del narrador quien reunió a toda la familia después de la muerte del abuelo.}
\item \esp{El narrador aprendió los platos típicos de su región en la cocina de su abuela.}
\item \esp{Lo que más admira de su abuela es su fuerza: nunca se queja y siempre tiene una palabra amable para cada uno.}
\item \esp{No, la familia no es rica. El texto dice: «No somos ricos, sino trabajadores».}
\end{enumerate}

\textbf{Questions de langue :}
\begin{enumerate}
\item Deux tournures clivées : \esp{Fue mi abuela quien reunió a toda la familia} = « C'est ma grand-mère qui a réuni toute la famille » ; \esp{Es en su cocina donde aprendí los platos típicos} = « C'est dans sa cuisine que j'ai appris les plats typiques » (on pouvait aussi relever \esp{fue ella quien me enseñó\ldots} ou \esp{es ella quien decide\ldots}).
\item Deux expressions de restriction : \esp{no tenía más que una pequeña casa} = « elle n'avait qu'une petite maison » ; \esp{No somos ricos, sino trabajadores} = « Nous ne sommes pas riches, mais travailleurs ».
\item \esp{Fue mi abuela quien reunió a toda la familia.} — tournure clivée sur le sujet (personne) : \esp{fue} + sujet + \esp{quien} + verbe.
\end{enumerate}
\end{corrige}

\begin{corrige}[Exercice 10 — Thème et version]
\textbf{Thème — proposition de traduction :}

\esp{Es mi padre quien más trabaja en nuestra familia. Es a las cinco de la mañana cuando sale para el campo de cacao. Solo tenemos una moto para toda la familia, pero nos basta. Lo que quiero es aprobar el bachillerato para que mis padres estén orgullosos. No es el dinero lo que hace la felicidad, sino la unión de la familia.}

\textbf{Points de vigilance :} « le plus » devant le verbe : \esp{más trabaja} ; « partir pour » : \esp{salir para} ; « suffire à » : \esp{bastar} avec complément d'objet indirect (\esp{nos basta}) ; « rendre fiers » : \esp{para que\ldots estén orgullosos} (subjonctif après \esp{para que}) ou \esp{para hacer orgullosos a mis padres} ; « réussir à un examen » : \esp{aprobar}, jamais \esp{realizar}.

\textbf{Version — proposition de traduction :}

« C'est mon père qui a construit cette maison. Nous n'avons qu'un champ, mais il nous suffit. C'est en été que toute la famille se réunit au village, et ce que j'aime le plus, c'est écouter les histoires de mes grands-parents. Ce ne sont pas mes cousins qui sont venus à la fête, ce sont mes oncles de Bata. »

\textbf{Points de vigilance :} \esp{el pueblo} = « le village » (faux ami : pas « le peuple » ici) ; \esp{nos basta} = « il nous suffit » ; \esp{sino} = « mais » dans l'opposition ; Bata est une ville de Guinée équatoriale.
\end{corrige}

\begin{corrige}[Exercice 11 — Rédaction type Bac : Describa a su familia]
\textbf{Proposition de rédaction modèle (environ 85 mots) :}

\begin{textebox}[Modelo de redacción]
Mi familia no es muy grande: somos cinco, mis padres, mis dos hermanos y yo. Es mi madre quien cuida la casa; es alta, delgada y muy paciente. Mi padre, en cambio, es bajo, moreno y trabajador: solo descansa los domingos. Mi hermano mayor es generoso y siempre me ayuda con los deberes. No tenemos más que una pequeña casa en Yaoundé, pero somos felices. Lo que más aprecio de mi familia es la unión: nunca estamos solos.
\end{textebox}

\textbf{Traduction :} « Ma famille n'est pas très grande : nous sommes cinq, mes parents, mes deux frères et moi. C'est ma mère qui s'occupe de la maison ; elle est grande, mince et très patiente. Mon père, en revanche, est petit, brun et travailleur : il ne se repose que le dimanche. Mon frère aîné est généreux et m'aide toujours avec les devoirs. Nous n'avons qu'une petite maison à Yaoundé, mais nous sommes heureux. Ce que j'apprécie le plus dans ma famille, c'est l'union : nous ne sommes jamais seuls. »

\textbf{Barème indicatif (sur 10) :} respect du sujet et du nombre de mots : 2 pts ; lexique de la parenté et de la description (au moins 6 mots de parenté, 3 adjectifs) : 3 pts ; une tournure d'emphase correcte (\esp{Es mi madre quien\ldots}, \esp{Lo que más aprecio\ldots}) et une restriction correcte (\esp{solo}, \esp{no\ldots más que}) : 3 pts ; grammaire et orthographe (accents, accords, \esp{ser}/\esp{estar}) : 2 pts.

\textbf{Points de vigilance :}
\begin{itemize}
\item Accents obligatoires : \esp{unión}, \esp{está}, \esp{madre}, \esp{más} ; un accent oublié coûte des points.
\item Description physique avec \esp{ser} (\esp{es alta}), état passager avec \esp{estar} (\esp{está cansada}).
\item Ne pas écrire \esp{es mi madre que\ldots} : personne sujet = \esp{quien}.
\item Compter les mots : rester entre 75 et 90.
\end{itemize}
\end{corrige}

\newpage

\coursec{Fiche bilan}

\begin{popfigure}
\begin{center}
\begin{tikzpicture}[
  mindmap,
  grow cyclic,
  every node/.style={concept, align=center, font=\small},
  root concept/.append style={concept color=popblue, font=\bfseries\small, text=white, minimum size=2.6cm},
  level 1 concept/.append style={sibling angle=60, level distance=3.4cm, minimum size=2.2cm},
  level 2 concept/.append style={sibling angle=45, level distance=2.2cm, minimum size=1.6cm, font=\scriptsize},
]
\node[root concept, align=center]{E04\\ Emphase et\\ description\\ de la famille}
  child[concept color=poporange, text=white]{ node[align=center]{C'est\ldots qui\\ \esp{Es\ldots quien}} 
    child[concept color=poporangeL]{ node[align=center]{\esp{Es Ana quien}\\\esp{cocina}} }
    child[concept color=poporangeL]{ node[align=center]{\esp{Fue ayer}\\\esp{cuando\ldots}} }
  }
  child[concept color=popgreen, text=white]{ node[align=center]{C'est\ldots que\\ \esp{Es\ldots lo que}} 
    child[concept color=popgreenL]{ node[align=center]{\esp{Lo que me gusta}\\\esp{es\ldots}} }
    child[concept color=popgreenL]{ node[align=center]{\esp{Es el cacao}\\\esp{el que\ldots}} }
  }
  child[concept color=poppurple, text=white]{ node[align=center]{Ne\ldots que\\ \esp{solo / no\ldots más que}} 
    child[concept color=poppurpleL]{ node[align=center]{\esp{No tengo más que}\\\esp{dos hermanos}} }
    child[concept color=poppurpleL]{ node {\esp{No\ldots sino}} }
  }
  child[concept color=poppink, text=white]{ node[align=center]{Emphase\\ en español} 
    child[concept color=poppinkL]{ node[align=center]{ordre des mots\\ verbe en fin} }
    child[concept color=poppinkL]{ node[align=center]{\esp{ser} +\\ tournure clivée} }
  }
  child[concept color=popgold, text=white]{ node[align=center]{Lexique\\ de la famille} 
    child[concept color=popgoldL]{ node[align=center]{\esp{padres, hermanos,}\\\esp{tíos, primos}} }
    child[concept color=popgoldL]{ node[align=center]{portrait physique\\ et moral} }
  }
  child[concept color=popblue, text=white]{ node[align=center]{Redacción\\ 80 palabras} 
    child[concept color=popblueL]{ node[align=center]{présentation\\ + portrait} }
    child[concept color=popblueL]{ node[align=center]{\esp{ser / tener /}\\\esp{llevarse bien}} }
  };
\end{tikzpicture}
\end{center}
\legende{Carte mentale de la leçon E04 : emphase, « ne\ldots que » et description de la famille.}
\end{popfigure}

\begin{bilan}
\textbf{1. Lexique clé : la famille et le portrait}
\begin{itemize}
  \item \esp{los padres} : les parents ; \esp{el padre / la madre} : le père / la mère ; \esp{el hermano / la hermana} : le frère / la sœur ; \esp{el hermanastro} : le demi-frère.
  \item \esp{los abuelos} : les grands-parents ; \esp{el tío / la tía} : l'oncle / la tante ; \esp{el primo / la prima} : le cousin / la cousine ; \esp{el sobrino} : le neveu ; \esp{el nieto} : le petit-fils.
  \item \esp{el marido / la mujer (esposa)} : le mari / la femme ; \esp{el hijo único} : le fils unique ; \esp{los gemelos} : les jumeaux.
  \item Physique : \esp{alto, bajo, delgado, moreno, de ojos negros, de pelo rizado} (grand, petit, mince, brun, aux yeux noirs, aux cheveux frisés).
  \item Moral : \esp{amable, generoso, trabajador, alegre, serio, paciente} (gentil, généreux, travailleur, joyeux, sérieux, patient).
  \item Relations : \esp{llevarse bien/mal con} (s'entendre bien/mal avec), \esp{parecerse a} (ressembler à), \esp{querer a} (aimer).
\end{itemize}

\textbf{2. Règles de grammaire à connaître par cœur}
\begin{center}
\begin{tabularx}{\linewidth}{|X|X|X|}
\hline
\rowcolor{popblueL} Français & Español & Exemple \\
\hline
C'est\ldots qui (sujet) & \esp{Es\ldots quien} / \esp{el que} & \esp{Es mi madre quien cocina.} (C'est ma mère qui cuisine.) \\
\hline
C'est\ldots que (COD) & \esp{Es\ldots lo que} / \esp{a quien} & \esp{Es a Juan a quien espero.} (C'est Juan que j'attends.) \\
\hline
C'est\ldots que/où/quand & \esp{Fue ayer cuando\ldots} & \esp{Fue ayer cuando llegó.} (C'est hier qu'il est arrivé.) \\
\hline
Ne\ldots que & \esp{solo, únicamente} & \esp{Tengo solo dos hermanos.} \\
\hline
Ne\ldots que & \esp{no\ldots más que} & \esp{No bebo más que agua.} \\
\hline
Ne\ldots que (opposition) & \esp{no\ldots sino (que)} & \esp{No estudia francés, sino español.} \\
\hline
\end{tabularx}
\end{center}
\begin{itemize}
  \item \cle{Emphase espagnole} : l'élément mis en relief est placé avant le verbe ou en fin de phrase ; la tournure clivée se construit avec \esp{ser} conjugué + \esp{quien / el que / lo que / donde / cuando}.
  \item \cle{Accord} : \esp{quien} (personnes, singulier), \esp{quienes} (pluriel) ; \esp{lo que} pour une idée ou une chose ; \esp{el que / la que / los que / las que} selon le genre et le nombre.
  \item \cle{Passé} : \esp{Fue + complément + cuando/donde\ldots} pour insister sur la date ou le lieu.
  \item \cle{Ne\ldots que} : jamais de calque ; choisir \esp{solo} (simple), \esp{no\ldots más que} (insistance), \esp{no\ldots sino} (correction d'une fausse idée).
\end{itemize}

\textbf{3. Méthodes express}
\begin{itemize}
  \item \cle{Thème} : repérer l'élément souligné par « c'est » (sujet ? COD ? circonstance ?), choisir \esp{quien / lo que / cuando / donde}, conjuguer \esp{ser} au bon temps.
  \item \cle{Redacción (80 mots)} : 1) présenter la famille (taille, lieu) ; 2) décrire 2 ou 3 membres (physique + caractère) ; 3) conclure sur les relations (\esp{me llevo muy bien con\ldots}). Compter les mots et relire les accords.
\end{itemize}
\end{bilan}

\begin{aretenir}[Checklist avant le Bac]
\begin{itemize}
  \item $\square$ Je sais traduire « c'est\ldots qui » par \esp{Es\ldots quien} ou \esp{el que}.
  \item $\square$ Je sais traduire « c'est\ldots que » par \esp{Es\ldots lo que} ou \esp{a quien} devant une personne COD.
  \item $\square$ Je sais insister sur une date avec \esp{Fue ayer cuando\ldots} et sur un lieu avec \esp{Fue en Douala donde\ldots}.
  \item $\square$ Je fais l'accord de \esp{quien/quienes} et de \esp{el que/la que/los que/las que}.
  \item $\square$ Je traduis « ne\ldots que » par \esp{solo} ou \esp{no\ldots más que}, jamais par un calque du français.
  \item $\square$ Je traduis « ce n'est pas X, c'est Y » par \esp{no\ldots sino\ldots}.
  \item $\square$ Je conjugue \esp{ser} au temps exigé par la phrase (\esp{es, fue, era, será}).
  \item $\square$ Je connais le lexique de la parenté : \esp{abuelos, tíos, primos, sobrinos, nietos, gemelos}.
  \item $\square$ Je distingue \esp{ser} (caractéristique durable) et \esp{estar} (état passager) dans un portrait.
  \item $\square$ Je sais utiliser \esp{llevarse bien con}, \esp{parecerse a} et \esp{querer a} avec la \esp{a} personnelle.
  \item $\square$ Je sais rédiger une description de famille en 80 mots : présentation, portraits, relations.
  \item $\square$ Je vérifie accents, \esp{ñ}, points d'interrogation et d'exclamation ouvrants (¿ ¡).
\end{itemize}
\end{aretenir}

\findecours

\end{document}
